% Автоматически перенесено из ГЛАВА_2_ВЕРСИЯ_3.docx

\chapter{Устный, письменный и ученый языки в каталонском обычном праве XII--XIII~вв.}

Под устным языком права в настоящей главе понимаются юридически значимые высказывания участников правового взаимодействия: обвинения, клятвы, отказы, распоряжения, требования и ответы. Письменный язык охватывает способы фиксации и текстового упорядочения этих действий в сводах обычного права и актовом материале. Учёный язык представлен терминологией, конструкциями, цитатами и способами толкования, связанными с профессиональной латинской письменностью и традицией римского и канонического права. Это разграничение основано на способах функционирования правовой речи и не совпадает с делением памятников на латинские и старокаталанские.

Цель главы состоит в том, чтобы установить, как эти формы правовой коммуникации взаимодействовали в каталонских памятниках XII--XIII~вв. Сначала рассматривается письменное представление словесных действий участников правового взаимодействия. Затем исследуются старокаталанские элементы в латиноязычных памятниках, латинские заимствования и цитаты в каталаноязычных сводах. Наконец, сопоставление латинской и старокаталанской версий «Барселонских обычаев» позволяет определить терминологические, синтаксические и смысловые изменения, возникавшие при переводе.

В центре анализа находится выбор участниками правового взаимодействия формы юридически значимого высказывания. Этот выбор доступен исследователю только в письменной передаче, поэтому его необходимо отличать от последующей работы составителей сводов, писцов, нотариусов и переводчиков, отбиравших, упорядочивавших и перерабатывавших сведения о словесных действиях.

\section{Диалектика устного и письменного в текстах каталонского обычного права}

\subsection{Устная традиция и её отражение в записях обычного права}

Основное внимание в параграфе уделено тому, как в сводах обычного права и актовом материале представлен выбор участниками правового взаимодействия слов и форм высказывания и как этот выбор фиксировался в письменном тексте. Рассматриваются речевые формулы и иные типовые реплики, связанные с обвинением, клятвой, отказом или распоряжением, указания на произнесение и восприятие слов, а также способы письменного подтверждения совершённого действия.

При источниковедческой интерпретации материала необходимо различать три уровня. Первый составляет словесное действие участника правового взаимодействия, совершавшееся или предполагавшееся в конкретной ситуации. Второй представляет его нормативную модель в записи обычного права, то есть предписание о том, что, кем и в какой последовательности должно быть сказано. Третий уровень составляет документальная фиксация отдельного действия в грамоте или судебном акте. Непосредственно исследованию доступны второй и третий уровни; устная практика реконструируется по ним с необходимой осторожностью.

Под формулой в настоящей главе понимается закреплённая в правовом тексте относительно устойчивая словесная форма, посредством которой выражалось содержание нормы, задавался образец высказывания участника правового взаимодействия или письменно фиксировалось юридически значимое действие. В зависимости от контекста одна и та же конструкция могла выступать как речевая, нормативная или документальная формула; эти характеристики не являются взаимоисключающими. Формула могла быть представлена в виде прямой речи, однако не всякая приведённая в источнике реплика имела формульный характер. Тем самым формула рассматривается не только как повторяющийся оборот, но и как способ нормативного или документального оформления правового действия.

Это разграничение не позволяет непосредственно отождествлять сохранившиеся формулы со словами, произнесёнными в конкретной ситуации. Своды обычного права показывают, какие словесные действия составители считали значимыми и как описывали порядок их надлежащего совершения. Грамоты фиксируют отдельные правовые действия при помощи документального формуляра и в некоторых случаях специально указывают, что слова были произнесены, услышаны или подтверждены свидетелями.

При анализе необходимо учитывать профессиональную и книжную среду, в которой создавались записи обычного права. Составители памятников были связаны с латинской письменностью, нотариальной практикой и в ряде случаев с учёным правом. Они систематизировали местный обычай средствами сложившейся культуры правового письма, отбирали материал, выбирали терминологию и синтаксические конструкции в соответствии с жанром и назначением памятника. Поэтому письменный текст сохранял сведения об устной практике в переработанном виде, а включённые в норму слова и формулы получали устойчивую грамматическую форму и соотносились с определёнными юридически значимыми действиями.

Соотношение устного слова, действия и письменной фиксации занимает важное место в исследованиях средневековой коммуникации. М.~Баньяр в монографии «Живая речь: письменная и устная коммуникация на латинском Западе с IV по IX~вв.» показал, что письменная коммуникация на латинском Западе сохраняла тесную связь с практиками произнесения, публичного чтения и восприятия текста на слух\footnote{\emph{Banniard M.} Viva voce: communication écrite et communication orale du IVe au IXe siècle en Occident latin. Paris: Institut des études augustiniennes, 1992. P.~46--47.}. Сформулированный им подход позволяет учитывать, что записанный текст мог сохранять связь с практиками произнесения, слушания и запоминания. В исследованиях М.Т.~Клэнчи, Р.Д.~Маккиттерик и М.~Зиммерманна эта проблематика получила развитие применительно к документальной культуре Средневековья, публичному оглашению актов и взаимодействию памяти с письменным удостоверением\footnote{\emph{Clanchy M.T. }Remembering the Past and the Good Old Law // History. 1970. Vol.~60. P.~165; \emph{McKitterick R.D.}~The Written Word and Oral Communication: Rome’s Legacy to the Franks // Latin Culture and Medieval Germanic Europe: Germania Latina I / ed. by R. North, T. Hofstra. Groningen: Egbert Forsten, 1992. P.~89--112; \emph{Zimmermann M.} Écrire et lire en Catalogne (IXe--XIIe siècle). Madrid: Casa de Velázquez, 2003. Vol.~1. P.~43; \emph{Сафронова Д.П. }Устная речь в правовой и документальной практике Леонского королевства X--XI~вв. // Диалог со временем. 2020. Вып.~72. С.~124--125.}. М.~Мостерт, обращаясь непосредственно к средневековому судебному процессу, подчеркнул необходимость совместного анализа речи, ритуального действия и текста\footnote{Mostert M. Introduction // Medieval Legal Process: Physical, Spoken and Written Performance in the Middle Ages / ed. by M. Mostert, \emph{P.S.~Barnwell}. Turnhout: Brepols, 2011. P.~3--5.}. В настоящем исследовании эти наблюдения используются для анализа того, как источники представляют выбор участниками правового взаимодействия словесной формы действия и как этот выбор оформлялся в нормативном и документальном тексте.

Устная составляющая представлена в интересующих нас источниках несколькими взаимосвязанными формами. Так, отдельные статьи записей обычного права могли прямо требовать произнесения обвинения, клятвы, отказа или распоряжения; предписывать назвать стороны, предмет и основание требования; приводить типовую реплику; распределять вопросы и ответы между участниками процедуры. В актовом материале словесное действие передавалось посредством документальных формул, прямой или косвенной речи, указаний на способность лица говорить, а также свидетельств о том, что присутствующие видели и слышали совершённое действие. Совместное рассмотрение этих форм позволяет проследить, как устное высказывание закреплялось в нормативных текстах и документах.

Прямая речь составляет одну из наиболее наглядных форм письменного представления словесного действия. Под прямой речью далее понимаются реплики участников либо образцы слов, которые надлежало произнести в определённой правовой ситуации. Такое оформление не означает, что текст дословно передаёт реальное высказывание: соответствующий фрагмент в своде обычного права мог представлять нормативную модель, а в грамоте --- письменную передачу слов, приписанных конкретному участнику. Если реплика имела относительно устойчивую словесную форму и была связана с определённой правовой функцией, она рассматривается как речевая формула; единичные реплики без этих признаков формулами не называются.

Актовый и нормативный материал раскрывают разные стороны этого процесса. Грамоты позволяют проследить, как сведения об отдельном словесном действии включались в документ и какими средствами подтверждалось его совершение: указанием на присутствие свидетелей, на услышанные ими слова либо на нотариальную или судебную запись. Своды обычного права показывают, каким образом словесное действие представлялось как обобщённое предписание и образец для повторного применения. Сопоставление этих источников позволяет установить, какие слова, состав участников и последовательность реплик связывались с надлежащим совершением правового действия.

В каталонских памятниках подобные конструкции распределены неравномерно. «Барселонские обычаи» преимущественно содержат краткие указания на необходимость произнести обвинение, клятву или иное заявление. «Обычаи Льейды», «Обычаи Орты» и «Обычаи Миравета» фиксируют отдельные формулы либо обязательный состав высказывания. В «Кутюмах Перпиньяна» нормы регулировали пределы словесного изложения требования при наличии публичного документа. Наиболее подробно речевые действия представлены в «Обычаях Тортосы», где прямые реплики включены в описание судебной процедуры, завещательных распоряжений и других правовых действий. Сопоставление этих памятников составляет основу дальнейшего анализа.

\subsection{Прямая речь в описании судебной процедуры}

Количественное сопоставление показывает, что прямая речь и близкие к ней словесные модели распределялись по каталонским правовым источникам XII--XIII~вв. крайне неравномерно. Наиболее репрезентативным в этом отношении остаётся текст «Обычаев Тортосы», где все случаи прямой речи сравнительно легко поддаются внутренней классификации. Все случаи использования «прямой речи» в тексте «Обычаев Тортосы» (99 в обеих редакциях памятника) довольно легко разграничить по обстоятельствам их использования -- можно выделить 3 категории, неравномерно представленные в тексте:

\begin{itemize}
\item Процессуальные формулы и реплики --- 72 случая употребления;
\item Типовые формулы (иногда в виде примеров-шаблонов) для составления завещаний --- 21 случай;
\item Клятвы и присяги --- 6 случаев употребления\footnote{Мы будем касаться преимущественно двух первых категорий, поскольку проблематика, связанная с принесением клятвы, её ролью и местом в средневековой каталонской правовой культуре требуют отдельного обстоятельного обсуждения, в том числе с привлечением других нормативных источников, а также по возможности актового материала. См. Главу 3 настоящего исследования.}.
\end{itemize}
В остальных памятниках соотношение оказывается иным. В «Обычаях Льейды» выявлено 5 случаев прямой речи или близких к ней типовых формул, из которых 1 относится к обвинительным формулам, 1 --- к завещательным распоряжениям, 3 --- к категории клятв и присяг. В «Обычаях Орты» употребление прямой речи, по-видимому, было сведено к минимуму: в настоящее время можно с уверенностью говорить об одном примере, связанном с возбуждением уголовного обвинения. Для «Обычаев Миравета» более характерны не готовые реплики, а нормы, предписывающие назвать сторону, предмет и основание требования. В «Обычаях Тарреги» и «Кутюмах Перпиньяна» прямой подсчёт случаев прямой речи затруднителен, поскольку в этих памятниках она почти не представлена как самостоятельный текстовый элемент. В первом случае процессуальное заявление передаётся глаголом «жаловаться» (лат. --- \emph{conqueri}) и его последующую квалификацию в гражданском или уголовном порядке. Во втором памятнике регулируются пределы словесного изложения требования при наличии публичного документа. Поэтому для этих источников более уместно учитывать не случаи прямой речи как таковой, а нормы, определяющие словесное оформление жалобы, иска или требования.

Для актового материала XII~в. количественная картина должна строиться несколько иначе, чем для сводов обычного права. Здесь следует различать, с одной стороны, устойчивые формулы отказа и уступки, зафиксированные в письменной форме, а с другой --- сравнительно редкие случаи, когда документ воспроизводит произнесённые слова и специально отмечает, что они были услышаны и увидены присутствующими. Гораздо чаще в грамотах встречаются устойчивые формулы типа \emph{«окончательный отказ от притязаний и уступка права»} (\emph{diffinitio et evacuatio}), \emph{«мы окончательно отказываемся от притязаний и уступаем право» }(\emph{difinimus et evacuamus}) и \emph{«я окончательно отказываюсь от притязаний и полностью уступаю право»} (\emph{diffinio atque modis omnibus evacuo}). В данном случае перед нами повторяющаяся формула-перформатив, присущая практике составления документов. Случаи, в которых прямая речь выводится в текст эксплицитно, следует рассматривать как особую группу, поскольку они позволяют наблюдать не только письменную фиксацию акта, но и момент его публичного совершения.

В «Обычаях Тортосы» прямая речь преимущественно включена в описание судебной процедуры: с её помощью передаются типовые требования, ответы, возражения и заявления участников процесса. В более кратких сводах, таких как «Обычаи Льейды», «Обычаи Орты» и «Обычаях Миравета», она либо используется на отдельных этапах судебной процедуры, либо уступает место метаязыковым указаниям на то, что должно быть сказано.

В королевской привилегии «Знатные мужи Барселоны признали…» (Recognoverunt proceres, 1284) устное высказывание уже было тесно связано с нотариальным заверением и словом свидетелей. Согласно главе XXV, нотариус мог составить завещание наедине с завещателем, однако затем должен был \emph{«призвать свидетелей, перед которыми он объявит, что составил завещание этого завещателя»}\footnote{Recognoverunt proceres. [XXV]: «…vocet testes coram quibus dicat se fecisse testamentum ipsius testatoris». P.~595}. Такое заявление нотариуса имело ту же силу, как если бы свидетели сами слышали волеизъявление завещателя. В главе XLVIII признавалось действительным устное завещательное распоряжение, сделанное в отсутствие нотариуса, при условии, что свидетели в течение шести месяцев приносили клятву о том, \emph{«что они именно так видели и слышали, как оно было записано или произнесено»}\footnote{Ibid. [XLVIII]: «…quod ipsi testes ita viderunt et audiverunt scribi seu dici». P.~599}. Согласно главе LI, сообщению сайона о вызове стороны в суд и назначении дня заседания доверяли при условии, \emph{«что соответствующая запись обнаруживается в реестре или в судебных актах»}\footnote{Ibid. [LI]: «creditur et statur relationi unius sagionis, dummodo reperiatur scriptura in capibrevio vel in actis iudicis». P.~600.}. Таким образом, произнесённое слово сохраняло самостоятельное правовое значение, однако его действительность обеспечивалась присутствием свидетелей, клятвой и нотариальной либо судебной записью. В грамотах же устойчивые формулы в большинстве случаев принадлежали, скорее, письменному формуляру, чем реконструируемой живой речи.

\subsection{Формализованный судебный диалог и речевые роли его участников}

Наибольший интерес для нас в данном случае представляют процессуальные элементы, снабжённые вставками с потенциальными репликами участников заседания, которые создают «эффект присутствия», как, например, в 1.3.14 статье «Обычаев Тортосы»: \emph{«Названые судьи должны сказать кредитору “был ли получен залог?”, и если он говорит: “Да”, то [они.--- А.К.] должны сказать: }\emph{“Сколько дней назад был получен [залог.--- А.К.]?”, и если он ответит и это будет правдой: “Три дня назад и более”, то они должны ему сказать: “Во сколько ты его оцениваешь?”, и он должен сказать: “Столько-то и столько”, то есть назвать ту цену, в которую он его [залог.--- А.К.] оценивает и, после оглашения цены, судьи должны спросить: “Кто является его покупателем?”, и посредник должен их уведомить, что покупатель определён и будет в суде в назначенный день и час, когда будет суд. Выше названные судьи должны сказать вегеру следующее: “О вегер, сообщи покупателю и спроси, собирается ли он присутствовать [на суде.--- А.К.]”, и если он откажется, то вегеру в первый день и час суда, когда будет суд, выше упомянутые судьи должны сказать и спросить у покупателя: “Купили ли вы эту вещь и дали ли за неё эту цену?”, и он должен сказать и ответить “да”, -- и ответить “да, дал эту цену”. Выше упомянутые судьи должны сказать вегеру следующее: “О вегер, сообщи должнику и спроси, собирается ли он присутствовать [на суде.--- А.К.]”, и если он откажется, то вегеру в первый день и час суда, когда будет суд, и на вопрос посредника должен сказать выше упомянутым судьям “Я нашёл этого человека и, как вы присудили, сказал <…> ему о залоге, о котором сообщил его кредитор, который потратил время и дал этому человеку достаточно. Должник ответил, что не придёт”, --- в таком случае, судьи верят вегеру и, услышав это, должны сказать писцу: “Сделай ему [кредитору. ---}\emph{ }\emph{А.К.] грамоту”...»}\footnote{\textsuperscript{ }Costums de Tortosa. 1.3.14~: «~Los quals jutges deven dir al corredor: “as correguda aquesta penyora?”, e si diu: “hoc”, deven dir: “quantz dies la as correguda?”, e si él respon e sia ver: “tres dies e plus”, deven-li dir: “què i trobes?”,e él deu dir: “aitant y trop”, ço és aquel preu que i troba e, dit lo preu, los jutges deven demanar: “qual és lo conprador?”, e el corredor deu-los assignar, lo qual comprador assignat e que sia present en la cort a dia e a ora de cort e que cort se tenga, los ditzjutges deven dir al veger enaxí: “en veger, fetz-ho saber al deutor, e deitz-li si hi vol acórrer”, e si diu que no, lo veger al primer dia de cort, dia e ora de cort e que cort se tenga, los ditz jutges deven dir e demanar al comprador: “comprat vós aquesta cosa e datz-hi aquest preu?”, e él deu dir e respondre als jutges: “hoc”, e fet lo respost: “hoc donar”, los jutges deven dir al veger enaxí: “en veger, fetz-ho saber al deutor e deitz-li si hi vol accórrer”, e si diu que no, lo veger al primer dia de cort e ora de cort e que cort se tinga, a demanda del creedor, \textbf{deu dir} als jutges: “jo trobe aital hom e axí com vós avíetz jutyat dix-li e fiu-li saber si volia acórrer a aital peynora que avia assignada a sson creedor, que corregut ha tot son temps e dóna-y hom aitant, e él respòs que no i volia acórrer”, e en aquest cas és creegut lo veger, els jutges, açò oït, deven dir a l'escrivà: “feiz-li la carta”». P.~30--31.}.

В данном случае мы имеем дело с нормами, регулирующими распродажу заложенного имущества по решению суда. Такие установления вполне могли бы быть записаны и без применения прямой речи, однако по какой-то причине составители первой редакции обычаев предпочли привести её именно в таком виде. Из 12 случаев употребления глаголов речи в приведённом фрагменте 10 сопровождаются глаголом долженствования «долженствовать» (кат. --- \emph{deure}), то есть нормативность выражена наиболее простым аналитическим способом. Вопросы со стороны судей и ответы посредника на них также сформулированы предельно кратко и в отдельных случаях, когда требуют конкретных сведений, представлены абстрактными формулировками (например, оценка стоимости заложенного имущества).

Стоит также отметить, что даже в рамках относительно небольшого по объёму фрагмента нормативного текста наблюдается повторяемость реплик отдельных сторон судебного заседания, особенно судей и должностных лиц. Это позволяет говорить об их формульной организации: повторяющиеся вопросы и ответы одновременно делают текст «звучащим» и нормативно упорядочивают словесное взаимодействие участников. Вместе с тем такая степень регламентации не позволяет непосредственно возводить приведённую прямую речь к реальной процессуальной практике.

Для того, чтобы подтвердить или опровергнуть совершение какого-либо действия, обвиняемый должен был, согласно тексту «Обычаев Тортосы», использовать строго определённые фразы, указания на действия или предметы. Для этого применялись различные формы выражения долженствования\footnote{\emph{Варьяш О.И. }Язык средневекового права // Пиренейские тетради: право, общество, власть и человек в Средние века / Сост. и отв. ред. \emph{И.И.~Варьяш}, \emph{Г.А.~Попова}. М.: Наука, 2006. С.~90.}. Получается, что «роли» всех сторон в ходе суда заранее предписываются всем его участникам: от обвиняемого до судей и королевских официалов. Сам судебный процесс предстаёт перед нами в качестве социального представления с детально прописанными диалогами\footnote{\emph{Mostert}\emph{ }\emph{M}\emph{. }Op. cit. P.~5--6.}. Отдельно стоит отметить, что особо объёмные вставки «прямой речи» (в том числе и приведённый выше) были исключены из текста в ходе составления второй редакции «Обычаев Тортосы»\footnote{На избыточность и перенасыщенность прямой речью мусульманских судебников, написанных на старокаталанском языке в более поздний период, указывает \emph{И.И.~Варьяш} в работе «Мусульманская Европа: сигналы идентичности»: \emph{Варьяш И.И.} Мусульманская Европа. Сигналы идентичности. СПб.: Наука, 2020. C.~20.}.

Показательно, что статья 2.7.2 «Обычаев Тортосы» прямо упоминает заявления, которые адвокаты произносили перед судьёй \emph{«на простом языке, то есть на романсе» }(старокат. --- \emph{en pla o en romanç})\footnote{Costums de Tortosa. 2.7.2~: «…e en plan o en romans diran e preposaran en pleyt denant lo jutgie…». P.~96. Как указывают \emph{Х.Х.~Чао} Фернандес, \emph{Х.Ф.~Меса} Санс, \emph{М.К.~Пуче} Лопес, сходную формулу можно встретить в привилегии Жауме I, выданной Валенсии 4 июня 1264~г. См. \emph{Chao Fernández J.J., Mesa Sanz J.F., Puche López M.C.} Latín y vernáculo en los documentos de Jaime I «El Conquistador». P.~4--5. URL: \url{https://rua.ua.es/server/api/core/bitstreams/1d1715b4-03df-44af-8f8f735d42598d67/content} (дата обращения: 29.07.2026).}. Содержащееся в критическом издании предположение, будто эта формула косвенно указывает на латынь как обычный язык устной судебной практики, представляется недостаточно обоснованным\footnote{Ibid. 2.7.2~: « L’expressió «en pla o en romanç diran» sembla indicativa que la llengua oral en la pràctica forense era la llatina». P.~96.». P.~96. N. 2.}. В средневековом каталанском слово со значением «простой, понятный» (старокат. --- \emph{pla}) употреблялось как обозначение народного языка и могло выступать синонимом обозначения «романс» (старокат. --- \emph{romanç})\footnote{Pla // \emph{Alcover A.M., Moll F. de B.}\emph{ }Diccionari català-valencià-balear. URL: \url{https://dcvb.iec.cat/results.asp?word=pla} (дата обращения: 29.07.2026); \emph{Mas i Miralles A. }Al voltant del mot explanar, ‘passar d’una llengua a una altra’ // Zeitschrift für Katalanistik. 2019. Bd. 32. S. 192.}. К.~Дуарте-и-Монсеррат также отмечал, что именно выражением «на простом языке или на романсе» (старокат. --- \emph{en pla o en romanç}) рукопись называет собственный язык\footnote{\emph{Duarte i Montserrat C.} Notes sobre el Llibre de les Costums de Tortosa // Revista de Llengua i Dret. 1985. Núm. 6. P.~23--25.}. Союз «или» (старокат. --- \emph{o}) имеет здесь, вероятнее всего, пояснительное значение: «на простом языке, то есть на романсе». Сама норма не определяет общий языковой режим суда, а устанавливает процессуальные последствия слов адвоката: то, что он «скажет и заявит» (старокат. --- \emph{diran e proposaran}) в присутствии и с согласия стороны, приравнивалось к заявлению самой стороны, тогда как очевидная ошибка могла быть отозвана в течение трёх дней.

Таким образом, статья свидетельствует прежде всего о правовом признании вернакуляра в судебном разбирательстве, но не позволяет непосредственно судить о языке письменного делопроизводства, которым, вероятно, оставался преимущественно латинский. В королевской канцелярии Жауме I Завоевателя латынь продолжала количественно преобладать. Так, например, исследованный Х.Х.~Чао Фернандесом, Х.Ф.~Месой Сансом, М.К.~Пуче Лопесом корпус документов канцелярии этого короля содержал около 96,5 \% латинских текстов\footnote{\emph{Chao Fernández J.J., Mesa Sanz J.F., Puche López M.C.} Latín y vernáculo en los documentos de Jaime I «El Conquistador». P.~5. URL: \url{https://rua.ua.es/server/api/core/bitstreams/1d1715b4-03df-44af-8f8f735d42598d67/content} (дата обращения: 29.07.2026).}.

\subsection{Формулы вызова и явки в суд в каталонском обычном праве}

До каталонских сводов XIII~в. словесная сторона обвинения уже была обозначена в более ранней правовой традиции, однако ещё не в виде краткой типовой реплики. Для этого этапа важнее были не сами слова обвинителя, а право предъявить обвинение, обращение к судье и условия, при которых дело переходило в судебное разбирательство.

Именно в таком виде эта процессуальная модель была закреплена в тексте «Вестготской правды». В 1-м титуле VI книги «Об обвинениях в преступлении» определялось, кто и в каком порядке мог выступать обвинителем в уголовных делах. Отдельно указывалось, что при недостатке доказательств обвинение могло быть переведено в письменную форму и удостоверено свидетелями\footnote{LV. VI.1: «De accusatoribus reorum». P.~246--256.}. В делах об убийстве в 5-м титуле VI книги право обвинять предоставлялось прежде всего родственникам убитого, но при их бездействии --- и другим лицам; если же никто не выступал с обвинением, судья должен был действовать сам\footnote{LV. VI.5.14--15: «Si homicidam nullus accuset, iudex, mox facti crimen agnoverit, licentiam habeat corripere criminosum <...> primo proximis occisi datur licentia accusandi; quod si iidem proximi<...> distulerint, tunc accusandi homicidam tam omnibus generaliter aliis parentibus quam exteris aditum pandimus<...> Nam homicidii reus numquam potest esse securus, cum contra eum accusationem deferre nulli penitus licentia denegetur». P.~280--281.}. Близкая логика прослеживалась и в 5 законе 1-го титула VII книги, где говорилось, что обвинитель по делу о краже, отравлении, колдовстве или ином преступлении должен явиться к городскому правителю или судье, чтобы дело было расследовано\footnote{LV. VII.1.5: «Quicumque accusatur in crimine, id est veneficio, maleficio, furto aut quibuscumque factis inlicitis, accusator eius concurrat ad comitem civitatis vel iudicem <...> ut ipsi secundum leges causam discutiant <...> Prius tamen pene non subiaceat, quam aut sub presentia iudicum manifestis probationibus arguatur, aut certe <...> eum accusator inscribat; et sic in presentia iudicum <...> questionis agitetur examen». P.~288.}. Таким образом, вестготская модель связывала начало уголовного преследования прежде всего с явкой обвинителя к судье и принятием на себя обязанности доказать обвинение. Вместе с тем закон уже допускал в ряде случаев предъявление обвинения в письменной форме (лат. --- \emph{per scripturam}).

На этом фоне «Барселонские обычаи» обозначили следующий шаг в области словесных форм возбуждения дела. В 90-й статье говорилось: \emph{«Обвинение не должно приниматься в письменной форме; обвинитель должен выдвигать его лично, своим голосом, если он обладает законным правом на обвинение, и делать это в присутствии того, кого намерен обвинить, поскольку никто не может ни обвинять, ни быть обвинённым заочно»}\footnote{Usatges de Barcelona. Us. B7 (90): «Per scripturam nullius accusacio suscipiatur, sed propria uoce accuset, si legitima et condigna accusatoris persona fiat, presente uidelicet eo quem accusare desiderat, quia nullus absens aut accusari potest aut accusare». P.~170.}\emph{.} Здесь ещё не приводилась типовая реплика обвинителя, однако сама норма прямо требовала, чтобы обвинение было произнесено вслух и совершалось при личном присутствии сторон. Тем самым уже в «Барселонских обычаях» устанавливался важный принцип: обвинение приобретало юридическую силу не через письменную запись, а через личное устное заявление, сделанное в надлежащей процессуальной ситуации.

Сопоставление двух памятников свидетельствует о различии в способах нормативного оформления обвинения, однако само по себе не позволяет говорить о прямом переосмыслении норм «Вестготской правды» составителями «Барселонских обычаев». Если в первой обвинение описывалось через обращение к судье и в отдельных случаях допускалось его предъявление в письменной форме, то последние закрепляли требование непосредственного устного предъявления обвинения самим обвинителем. Такая формулировка могла отражать как локальную судебную практику, так и переработку более ранней правовой традиции. В дальнейшем требование непосредственного предъявления обвинения получил более детальную разработку в других каталонских сводах обычного права.

В этом контексте важно обратиться к нормам «Обычаев Льейды», 162 статья которых предполагала, что обвинение не должно совершаться \emph{«ни с какой юридической торжественностью»} и не оформлялось письменно, а обычная его форма сводилась к краткой реплике, которую было необходимо произнести обвинителю или лицу, заявившему о преступлении: \emph{«Я обвиняю такого-то в том, что он убил такого-то»}\footnote{Costums de Lleida. [162]: «De accusationibus. Non fiunt hic accusaciones cum aliqua sollempnitate vel scriptura, set sic fieri consuevit: “Ego talis conqueror de tali qui occidit talem”». P.~146.}. Тем самым в «Обычаях Льейды» возбуждение дела связывалось с произнесением типовой формулы обвинения. Обстоятельство, что статья приводила эту формулу полностью, позволяет рассматривать последнюю, скорее, как нормативный словесный образ. В этой реплике в сжатом виде были обозначены все элементы, необходимые для обвинения: субъект высказывания, лицо, против которого оно направлено, и само преступное деяние. Значение статьи, таким образом, заключалось не только в признании устного обвинения, но и в упорядочении речевой формы, в которой оно должно было предъявляться.

Сходную картину можно наблюдать и в «Обычаях Орты» 1296~г. Здесь также специально оговаривалось, что при возбуждении дела не требовались ни письменная запись, ни иная юридическая формальность; было достаточно, чтобы обвинитель или доносчик произнёс установленную реплику: \emph{«Я предъявляю такому-то уголовное обвинение в том, что он убил такого-то»}\footnote{Costums de Orta. 77: «Item, quod accusatores criminales non habeant fieri cum inscrip- cione vel alia juris sollempnitate, set sufficiat quod accusator vel denunciator dicat: “conqueror criminaliter de tali, qui occidit talem”». P.~623.}\emph{. }Насколько можно судить по тексту памятника, это был единственный случай прямой речи во всём своде, и относился он именно к моменту возбуждения уголовного обвинения. Уже одно это показывало, какое значение составители придавали начальному процессуальному акту. Если в «Обычаях Льейды» формула ограничивалась указанием на самого обвиняемого и совершённое им деяние, то в «Обычаях Орты» внутрь реплики вводилось дополнительное уточнение --- “\emph{criminaliter”}. По-видимому, к концу XIII~в. этого потребовало само развитие процессуальной техники: теперь важно было не только назвать лицо и преступление, но и прямо обозначить характер совершаемого действия, то есть указать, что речь шла именно об уголовном обвинении. Если в «Обычаях Льейды» уголовно-судебный смысл ещё задавался общим контекстом статьи, то в «Обычаях Орты» он уже переносился внутрь формулы.

В этом отношении «Обычаи Льейды» и «Обычаи Орты» обнаруживали общую черту: в обоих случаях прямая речь появлялась лишь в единственной строго локализованной процессуальной точке, а именно в момент возбуждения дела. Её функция состояла в том, чтобы дать образец высказывания, после которого жалоба обретала юридическую силу и могла повлечь за собой дальнейшее разбирательство. Перед нами в данном случае крайне лаконичная словесная модель, необходимая для того, чтобы процесс вообще начался.

«Обычаи Тортосы» демонстрируют несколько иную ситуацию: здесь прямая речь уже не ограничивалась единичной формулой, открывающей разбирательство, а была встроена в само содержание процессуального права. Несколько рубрик девятой книги этого кодекса были посвящены вопросам вызова в суд, объявления тяжбы и судебного разбирательства, и в этих рубриках составители нередко включали в текст образцы слов, которые должны были быть произнесены участниками процедуры. Судебное заседание, например, могло состояться лишь после фиксации жалобы вегером или сайоном в строго определённой форме: \emph{«Вегер или сайон, как только будет объявлено о необходимости явиться в суд, должен сказать и назвать поимённо: “Такой-то человек обвиняет вас, в такой-то день [во второй редакции: в такой-то день и час] вам следует явиться в суд”}\emph{.}\emph{ }\emph{Е}\emph{сли же он не назовёт имени истца, этот приказ не будет иметь силы}\emph{, и ответчик не понесёт никакого наказания за то, в чём его обвинил истец»}\footnote{Costums de Tortosa. 1.3.2: «Lo veguer o el sag, en qual que manament fassen de venir a cort, deu dir e nomenar: “aytal hom se clama de vós, siatz aytal dia (o aytal hora) a la cort”, e si no li nomena lo nom del clamant, no val aquel manament ne cau en pena neguna aquel a qui hom auràfeyt lo manament”». P.~20.}. В отличие от «Обычаев Льейды» и «Обычаев Орты», где прямая речь лишь вводила обвинение в процесс, в «Обычаях Тортосы» она уже начинала регулировать и дальнейший ход судебного действия: вызов, явку, объявление тяжбы, фиксацию сторон.

Более того, сама структура «Обычаев Тортосы» показывала, что составителей интересовал не просто факт произнесения реплики, но и точность её содержания. Значимым оказывалось только имя истца, а все остальные слова, в том числе порядок участников тяжбы, а также день и час явки характеризовались меньшей значимостью.

На этом фоне особенно показательна и рубрика “\emph{De inquisitione”}\footnote{Ibid. P.~610.}, которая, напротив, была написана почти без использования прямой речи, хотя нормы розыскного процесса были непосредственно связаны с расспросом и устным выяснением обстоятельств дела. По-видимому, составители легче фиксировали в виде устойчивых словесных формул те элементы процесса, которые были связаны с обвинением, вызовом или обозначением позиции сторон. Напротив, розыскной процесс, в большей степени зависевший от обстоятельств конкретного дела, хуже поддавался такому описанию. Включение прямой речи в текст обычая зависело не просто от того, насколько важна была устная коммуникация, а от того, насколько данное действие могло быть сведено к повторяемой и юридически значимой формуле.

Таким образом, сопоставление этих памятников позволяет увидеть не только общую роль слова в судебной практике, но и различия в способах его фиксации. В «Обычаях Льейды» прямая речь была сведена к краткой формуле возбуждения обвинения. В «Обычаях Орты» эта формула получила ещё и дополнительное процессуальное уточнение через \emph{criminaliter}. В «Обычаях Тортосы», напротив, прямая речь использовалась значительно шире и служила уже не только для обозначения начала процесса, но и для описания ряда последовательных действий, от которых зависела юридическая действительность разбирательства. Поэтому составители «Обычаев Тортосы» зафиксировали целый набор речевых актов, встроенных в процессуальную процедуру. Именно это отличало её от более сжатых формулировок, представленных в сводах обычного права Льейды и Орты.

\subsection{Формулы передачи имущества и распоряжения им по завещанию}

Формулы передачи имущества и распоряжения им по завещанию позволяют увидеть, каким образом в правовой практике Каталонии XII--XIII~вв. юридическое действие связывалось с произнесением определённых слов или с их последующей письменной фиксацией. Более ранняя традиция уже предполагала, что распоряжение имуществом могло исходить из устного волеизъявления и лишь затем получать письменное оформление. Она исходила из того, что распоряжение имуществом могло совершаться не только письменно, но и устно, в присутствии свидетелей, а затем получать письменное или квазиписьменное подтверждение. Именно так эта модель была закреплена в 5-м титуле II книги «Вестготской правды»:

\begin{itemize}
\item 10-й закон допускал, что несовершеннолетние старше десяти лет в случае тяжёлой болезни или близкой смерти могли распорядиться своим имуществом \emph{«по завещанию или иным образом»} как письменно, так и в присутствии свидетелей\footnote{LV. II.5.10: «Morientium extrema voluntas... sive tantummodo verbis coram probatione... promulgata: hec ordinationum quatuor genera omni perenniter valere subsistant... Illa vero voluntas defuncti... que verbis tantummodo coram probatione promulgata patuerit... tunc robur plenissimum obtinebit, si testes ipsi, qui hoc audierint... infra sex mensuum spatium... sua coram iudice iuratione confirment». P.~112--113.};
\item в 11(13)-м законе того же титула прямо говорилось, что последняя воля умирающего могла быть выражена и устно, но тогда свидетели должны были в течение шести месяцев подтвердить под присягой перед судьёй то, что слышали от завещателя\footnote{LV. II.5.11(13): «In itinere pergens aut in expeditione publica moriens, si ingenuos secum non habeat, volumtatem suam propria manu conscribat. Quod si litteras nescierit aut pre langore scribere non potuerit, eandem volumtatem servis insinuet, quorum fidem episcopus adque iudex probare debebunt. Et si nullatenus antea fraudulenti fuisse patuerint, quod sub iuramenti taxatione protulerint, conscribatur, ut sacerdotis adque iudicis suscriptione firmetur; hac postmodum autoritate regia roboratum, firmum quod decreverit habeatur». P.~114.};
\item закон 13(15), в свою очередь, требовал публикации письменного завещания в течение того же срока\footnote{LV. II.5.13 (15): «Omnes scripture, quarum et auctor et testis defunctus est, in quibus tamen suscriptio vel signum conditoris adque testium firmitas repperitur, dum in audientia prolate constiterint, ex aliis cartarum signis vel suscriptionibus contropentur, suffi- ciatque ad firmitatem vel veritatis huiuis indaginem agnoscendam trium aut quattuor scripturarum similis et evidens prolata suscriptio. Quod si talibus scripturis legum tempora obviaverint, pro certo decernitur, quia valere non poterunt». P.~115.}.
\end{itemize}
Тем самым устное распоряжение не противопоставлялось письму, а включалось в общий порядок последующего удостоверения и обнародования. Однако в каталонском материале особенно важно то, что в ряде случаев до нас дошло не только указание на сам акт распоряжения, но и словесные модели, посредством которых он должен был совершаться.

Особое место в этом ряду занимали «Барселонские обычаи», поскольку именно в них достаточно рано обозначились основные требования к словесной определённости распоряжения. Хотя текст, как правило, ещё не воспроизводил готовую реплику завещателя или дарителя, он уже исходил из того, что подобное действие должно быть выражено ясно и полно. Так, в 120 статье «Барселонских обычаев», посвящённой лишению наследства, устанавливалось: \emph{«Если кто-либо захочет лишить наследства сына, дочь, внука или внучку, он должен сделать это поимённо, указать вину, по которой лишает их наследства, и назначить на их место другого наследника}\emph{. П}\emph{ри этом назначенный наследник должен доказать, что основание для лишения наследства действительно существует. Если хотя бы одно из этих условий не соблюдается, лишение наследства считается недействительным и не имеет силы»}\footnote{Usatges de Barcelona. Us. 120(78): «SI QVIS FILIVM SVVM uel filiam seu nepotem suum siue neptem exheredare uoluerit, nominatim illum exheredet, et culpam per quam exheredet dicat, et alium in loco suo instituat, et causa exheredationis ab eo qui instituitur heres uera esse probetur. Si unum ex hiis defuerit, exheredare filium suum uel filiam, nepotem uel neptem nullo modo poterit; et si presumpserit, irritum erit et nil ualebit». P.~154.}. Тем самым распоряжение считалось действительным лишь при условии, что лицо, основание и новый порядок наследования были обозначены явно.

О близкой установке свидетельствовала и норма, закреплённая в статье C3 (148) «Барселонских обычаев», согласно которой, если в ходе тяжбы предъявлялось завещание или грамота, а другая сторона не могла опровергнуть её ни свидетелями, ни иными подписанными документами, судья должен был вынести решение по справедливости и сохранить каждому его право\footnote{Ibid. Us. C3(148): «Si quis testamentum uel cartam firmatam de aliqua contencione in placito ostenderit, de alia parte neque per testes neque per firmas scripturas conuincere potuerit, iudex dicat quod rectum ei uidetur, et seruet unicuique suum directum». P.~171.}. В обоих случаях письменный акт уже занимал важное место, однако решающим оставалось то, насколько ясно в нём были выражены правовые притязания и воля распорядителя.

Каталонский актовый материал XII~в. показывает, что распоряжение правом, отказ от него или уступка имущества далеко не всегда оформлялись через прямую речь в строгом смысле слова. Неоднократно мы встречались в грамотах с формулами типа \emph{«окончательный отказ от притязаний и уступка права»} (\emph{diffinitio et evacuatio}), \emph{«мы окончательно отказываемся от притязаний и уступаем право» }(\emph{difinimus et evacuamus}) и \emph{«я окончательно отказываюсь от притязаний и полностью уступаю право»} (\emph{diffinio atque modis omnibus evacuo}). Такие сочетания повторяются в одном и том же правовом контексте, сохраняют устойчивое глагольное ядро и передают один и тот же тип действия --- окончательный отказ от притязаний, их снятие и передачу права другой стороне. Даже тогда, когда грамота строится от первого лица, это ещё не означает, что перед нами запись реальной устной реплики: в большинстве случаев речь идёт о письменной формуле, с помощью которой документ оформляет уже признанное действие.

На это указывает и частотность подобных оборотов: в выборке из 873 грамот формулы типа \emph{diffinitio et evacuatio} и \emph{diffinio atque modis omnibus evacuo} встречаются несколько сотен раз. Такая повторяемость позволяет видеть в них не случайные словесные решения отдельных писцов, а один из базовых приёмов при оформлении актов такого типа. Редакторы сборника J. M. Salrach i Marès и T. de Montagut i Estragués отдельно выделили акты окончательного отказа от притязаний, признания права другой стороны и прекращения спора как один из наиболее значимых типов документов XII~в. и подчеркнули их особую роль в практике разрешения конфликтов\footnote{Justícia i resolució de conflictes… 2024. P.~XIV.}.

Эту группу хорошо иллюстрирует грамота №~169 от 16 марта 1133~г. Бернат Адалберт де Навата и его сыновья Бернат и Арнау в ней признали, что всё принадлежавшее им в приходе Сан-Марк-де-Креспия по праву относилось к кафедральному собору Святой Марии города Жироны. После этого они формально отказались от своих притязаний и передали в руки епископа Жироны Беренгера церковь Сан-Пере-де-Навата с десятинами, приношениями первых плодов, аллодами и другими пожертвованиями, половину сборов с Баскары, а также синодальную подать с церкви Сан-Эстеве-д’эн-Бас\footnote{Ibid. №~169: «Prephata autem omnia ego predictus Bernardus cum filiis meis prenominatis Bernardo atque Arnallo diffinio Domino Deo et Sancte Marie ac modis omnibus derelinquo in manu dompni Berengarii, Gerundensis episcopi, successorumque suorum, non solum ea que supranominata sunt, sed etiam omnia illa que de canonicali honore Gerundensis ecclesie alicubi poterint repperiri, diffinio atque modis omnibus evacuo sicut ad utilitatem Gerundensis ecclesie persona aliqua sanius intelligere poterit. Quam diffinitionem sive evacuationem si quis infringere temptaverit, non valeat, sed duplam Deo et prephate Sancte Marie compositionem persolvat». P.~194.}. В центре такого акта стоял документально зафиксированный отказ. Сходным образом была составлена и грамота №~315 от 1 июля 1156~г., посвящённая уступке прав в пользу церкви: Оливарий отказался от прав на острова и от каких-либо требований к их жителям, оформив отказ той же устойчивой формулой\footnote{Ibid. №~315: «Ego sepedictus Olivarius, castellanus de Olone, consilio et voluntate domini mei Arnalli de Olone et uxoris eius Beatricis et filiorum suorum Petri et Guillemi ac bono animo et spontanea voluntate, diffinio et evacuo Domino Deo et Beate Marie Stagnensi et Berengario, ipsius ecclesie priori». P.~359.}. В грамоте 1156/1157~г. описана подобная ситуация, снабженная дополнительным пояснением: речь идёт о споре Беренгера де Англеса с Жиронской церковью по поводу церковных поборов и прав, связанных с церквами, входившими в его держание от Экстрании, его покойной жены. В этой грамоте сам Беренгер, оформляя отказ от притязаний, специально добавил, что \emph{«такой же отказ и уступку ранее совершила и моя жена Экстрания, находясь при смерти, но в здравом уме, твёрдой памяти и }\textbf{\emph{п}}\textbf{\emph{ри }}\textbf{\emph{сохранной речи}}\emph{»}\footnote{Ibid. №~327: «Hanc ipsam diffinitionem et evacuationem fecit dudum Extranea uxor mea in extremo mortis periculo posita, pleno tamen sensu, memoria ac loquela…». P.~374.}. В этих актах завершение спора закреплялось письменной формулой отказа от притязаний, при этом слова участников дословно не воспроизводились.

Переходный характер такого материала особенно хорошо виден в тех грамотах, где составитель подчеркнул способность лица говорить в момент совершения акта. В уже упомянутой грамоте 1156/1157~г. об отказе и уступке Экстрании важно именно указание на речь (лат.---\emph{loquela}) наряду с памятью и разумом: тем самым наличие речи включается в число условий действительности её отказа\footnote{Ibid. №~327: «Hanc ipsam diffinitionem et evacuationem fecit dudum Extranea uxor mea in extremo mortis periculo posita, pleno tamen sensu, memoria ac loquela…». P.~374.}.

Та же установка встречается и в грамоте №~556, согласно которой Пере Бернат де Ур, передавая своё тело для погребения и оставляя клиру кафедрального собора Святой Марии имущество в Аласе, специально заявил, что сделал это \emph{«никем не принуждаемый, по собственной воле, при целой памяти и }\textbf{\emph{сохранной речи}}\emph{»}\footnote{Ibid. №~556: «Notum sit cunctis presentibus et futuris quod ego Petrus Bernardus de Ur, nullo cogente set spontanea voluntate, memoria et loquela integra, propter remedium anime mee relinquo corpus meum orto Beate Marie et relinquo eius canonice, ob remissionem peccatorum meorum et parentum meorum, quicquid habeo in villa de Alas et in eius terminis et quicquid ibi aliquo titulo habere debeo». P.~677.}.

Наконец, в грамоте №~483 свидетели клялись, что они \emph{«видели и слышали»} и присутствовали при том, как Бернат, тяжело больной и лежащий в доме архидиакона в Жироне, будучи ещё \emph{«в твёрдой памяти и }\textbf{\emph{при сохранной речи}}\emph{»}, распорядился своей последней волей относительно замка Лорет\footnote{Ibid. №~483: «…dum adhuc esset in plena memoria ac loquela ordinavit suam extremam voluntatem de castro de Laureto, sicut superius scriptum est». P.~590.}. Здесь речь уже не воспроизводится в прямой форме, но способность завещателя говорить фиксируется как признак его возможности выразить последнюю волю. Составитель посчитал необходимым подчеркнуть, что завещатель не только помнил и понимал, что делает, но и мог выразить свою волю словами. Формулы, содержащие указания на сохранность \textbf{памяти }(лат. --- \emph{memoria}), \textbf{здравого рассудка} (лат. --- \emph{sensus}) и \textbf{способности говорить} (лат. --- \emph{loquela}), характеризовали возможность лица осознанно и словесно выразить свою волю.

Тем заметнее немногочисленные случаи, когда составитель грамоты включал в текст не только итоговую формулу, но и сам момент её произнесения. Особенно интересна в этом отношении грамота №~522. Она касается спора между приором церкви свв. Иоанна и Руфа в Льейде и Бертраном де Бенифальетом о неуплате ценза в шесть кувшинов масла с держания, прежде принадлежавшего Беренгеру де Торрохе. После начала разбирательства Бертран пытался добиться для себя более выгодных условий держания, однако судьи поставили его перед выбором: либо уплатить ценз, либо отказаться от держания. Затем, \emph{«посоветовавшись с друзьями, он вернулся в курию и }\textbf{\emph{громким голосом произнёс}}\emph{: }\textbf{“}\textbf{\emph{Я отказываюсь от всего этого держания и оставляю его}}\textbf{”}\emph{»}\footnote{Ibid. №~522: «Quo audito Bertrandus habuit consilium cum amicis, quo habito reddierint ad curiam et alta voce dixit: “Totum honorem illorum finio et desemparo”. Prior hoc audiens egit gratias omni conventui curie. Hanc diffinitionem audierunt et viderunt supradicti iudices et prior Poncius, Dertusensis ecclesie, et Theobaldus et Guillelmus de Ylerda, canonici supradicte ecclesie, et Moroni Ianuensis et Guillelmus Adei et Iordanus Garidelli et Raimundus de la Granada et multi alii». P.~637.}. Ещё важнее следующее обстоятельство: это действие \emph{«слышали и видели}» судьи, приор, каноники и многие другие присутствующие. Здесь словесное действие и его восприятие образуют единое целое: важно не только то, что Бертран публично произнёс формулу отказа, но и то, что это было публично услышано и увидено. Здесь грамота сохранила саму сцену отказа и коллективное свидетельство этого момента.

Таким образом, актовый материал XII~в. показывает два уровня фиксации распоряжения имуществом. В большинстве грамот словесное действие не выводится в виде прямой речи, а оформляется через устойчивую письменную формулу отказа или уступки. Формулы типа \emph{окончательный отказ от притязаний и уступка права} (лат. --- \emph{diffinitio et evacuatio}) следует рассматривать как один из базовых элементов документальной практики XII~в. Однако в отдельных случаях составитель сознательно отмечал сам момент произнесения слов, а также то, что они были услышаны и увидены присутствующими. Именно эти редкие случаи особенно важны для историко-антропологического анализа, поскольку позволяют увидеть, как между устным действием и письменной фиксацией распределялись роли: формула могла существовать как письменный шаблон, при этом в ряде ситуаций составитель документа счёл необходимым указать на публичность и коллективное восприятие совершаемого акта.

Если вернуться в поле нормативных источников, стоит упомянуть о 145-й статье «Обычаев Льейды», в которой специально оговаривалось, что наследники не назначались поимённо. Вместо этого в завещании назначались душеприказчики, к которым следовало обратиться в такой форме: \emph{«прошу, чтобы они разделили всё моё имущество [таким образом.---А.К.], как будет показано ниже»}\footnote{Costums de Lleida. [145]: «Non instituuntur heredes nominatim per consuetudinem, set fiunt manumissores in testamento qui rogantur sic: “Precor ut dividant omnia bona mea sicut inferius apparebit”». P.~141.}. Для данного контекста существенно, что текст приводил не просто указание на право распорядиться имуществом, а именно словесный образец поручения. Перед нами снова краткая формула, встроенная в завещание и рассчитанная на дальнейшее письменное развёртывание. Тем самым распоряжение имуществом приобретало двойной характер: оно ещё связывалось с произнесёнными словами завещателя, но сами эти слова уже тяготели к устойчивому формуляру.

Ещё яснее эта тенденция выступала в королевской привилегии «Знатные мужи Барселоны признали…» (Recognoverunt proceres). В 56-й статье, посвящённой составлению брачного договора, приводилась готовая модель распоряжения: \emph{«Я, такой-то, даю тебе, такой-то моей жене, столько-то в качестве твоего брачного дара»}\footnote{Recognoverunt proceres. [LVI]: «Item, quando instrumentum dotalium fit in hec verba: “Ego, talis, dono tibi tali uxori mee, tantum pro dote tua et sponsalicio”». P.~601.}. Здесь особенно заметно, что словесная модель уже почти совпадала с самим инструментом. Юридическое действие оформлялось через реплику от первого лица, но эта реплика была заранее приведена к устойчивому образцу. Если в «Обычаях Льейды» ещё сохранялась относительно краткая форма поручения, то в привилегии «Знатные мужи Барселоны признали…» (Recognoverunt proceres) перед нами был уже почти полный формуляр имущественного распоряжения, в котором говорящий и адресат были заданы заранее.

На этом фоне особое место занимали «Обычаи Тортосы», где подобные словесные модели уже не только предполагались, но и включались в нормативный текст как образец. В статье 5.6.21, посвящённой назначению опекуна детям, приводилась следующая формула: если кто-либо \emph{«даст или установит опекуна своим детям и скажет следующим образом: “Даю моим детям такого-то человека опекуном”»}\footnote{Costums de Tortosa. 5.6.21: «Si alcun hom dóna o establex tutor o curadora ssos fils e dirà enaxí: “do a mos fils aytal hom tutor o curador”, en aquest mot “mos fils”, enaxí bén y són enteses les files com los fils, e enaxí bén és tutor o curador de les files com dels fils, com en aquesta paraula “mos fils” s'entén files e fils e els pòstumos». P.~294.}. Такое назначение признавалось действительным. Существенно не только то, что в тексте воспроизводилась сама реплика, но и то, что она вводилась глаголом \emph{«говорить»} (кат.--- \emph{dir}), то есть прямо связывалась с актом говорения. Сразу вслед за этим составители разъясняли смысл выражения \emph{«мои дети»}\emph{ }(кат. --- \emph{mos fils}), включая в него не только уже рождённых детей, но и тех, кто мог появиться на свет после смерти отца. Следовательно, норма фиксировала не только образец высказывания, но и его обязательное юридическое толкование.

В следующей статье (5.6.22) дополнительно уточнялось, что выражение «мои дети» являлось наиболее широкой из возможных формулировок\footnote{Ibid. 5.6.22: «Testador que establex tutor a son fil e no nomena el nom del fil, és entés que a totz los altres fils lo fa tutor e no tant solament aquels fils és entés que l' establex tutor que a en son poder». P.~294.}, а потому при отсутствии специальных уточнений охватывало всех детей; тем самым сам выбор слов в подобных актах приобретал самостоятельное правовое значение.

По сходной модели строилась и статья 6.4.13 о назначении наследника с субституцией. Завещатель объявлял: \emph{«Я назначаю Жоана моим наследником, но если он не примет, или не сможет, или не захочет принять наследство, то есть скажет, что не хочет быть наследником, то пусть его заменит Бернат»}\footnote{Ibid. 6.4.13: «Si testador establex alcun hereu, pot fer substitut a aquel en aquesta forma: « “Johan hereu meu; e si no pren la heretat o no la pot pendre o no la vol pendre, so és que diga que no vol ésser hereu, substituex-li Bernat”». P.~310.}. Здесь норма вновь фиксировала не только сам результат распоряжения, но и его словесную форму. Особенно важно, что и возможный отказ наследника мыслился как действие, выражаемое через слова: отказ не просто предполагался, а связывался с произнесением определённой формулы. Тем самым текст воспроизводил уже не одно волеизъявление, а целую цепь речевых действий: назначение наследника, возможный отказ и последующую замену. В сравнении с более ранними каталонскими сводами именно «Обычаи Тортосы» наиболее полно сохраняли эту связь между юридическим действием и его словесным оформлением.

Таким образом, в каталонском обычном праве XII--XIII~вв. можно проследить постепенное оформление речевых моделей распоряжения имуществом относительно тех, которые использовались в предшествующий период в соответствии с установлениями «Вестготской правды». В «Барселонских обычаях» требование словесной определённости ещё не переходило в типовую реплику. В то же время в грамотах XII~в. устойчивые формулы уже передавали отказ, уступку и дарение как совершаемое словом действие; в «Обычаях Льейды» и в королевской привилегии «Знатные мужи Барселоны признали…» такие формулы всё заметнее приближались к готовому шаблону для завещания, уступки имущества или брачного договора. В «Обычаях Тортосы» типовая реплика была прямо включена в норму и становилась частью её содержания. В результате словесная формула выступала здесь не внешним сопровождением правового акта, а одной из форм его выражения.

\subsection{Соотношение речевой формулы и письменного удостоверения}

Если в предшествующем подпараграфе в центре внимания находились словесные модели передачи имущества и распоряжения им по завещанию, то здесь важнее проследить, каким образом подобные действия приобретали юридическую силу в случае спора.

Показателен и документ №328 от 2 апреля 1156/1157~г., связанный со спором о замке Гаварра. В нём вице-граф Жерал де Кабрера заявлял, что \emph{«видел, слышал и узнал»}, что его дед некогда подарил этот замок церкви \emph{«по дарственной грамоте»} и затем вновь оставил его той же церкви в своём завещании\footnote{Justícia i resolució de conflictes… №328: «Propterea veni ad Sedem, vidi, audivi et cognovi quod Gerallus Poncii, meus avus, donavit Deo et Sancte Marie Sedis eiusque canonice prenominatum castrum per scripturam donacionis et postea illud eidem ecclesie in testamento suo dimisit». P.~375.}. Здесь особенно важно сочетание нескольких уровней подтверждения: память о прежнем действии, ссылка на письменный дарственный акт и повторное завещательное распоряжение. Перед нами уже целая цепь удостоверений, где устное знание и документ взаимодействовали в системе доказательств.

Ещё яснее переход к доказательному использованию документа виден в грамоте 1186~г. о споре вокруг замка Иварс. Беатриса де Монтросс представила завещание своего мужа Пере де Павии, в котором было записано, что замок оставлялся их малолетнему сыну. Однако одной только записи оказалось недостаточно: затем сам Пере де Павия-младший, допрошенный судьями, заявил, что, согласно 88-й статье «Барселонских обычаев», он \emph{«видел и слышал»}, как его мать совершила дарение и передачу замка\footnote{Ibid. №328: «Beatrix vero de Montross produxit sacramentale testamenti viri sui Petri de Pavia, in quo continebatur quod ipse reliquerat castrum de Ivarz iam dicto pupillo, Beatrix, uxore sua, assenciente et subscribente in sacramentali. Hostendebat quoque Petrum de Pavia filium suum indicem. Qui interrogatus a iudicibus protestatus est se indicare secundum usaticum se vidisse et audisse quod mater sua fecit donacionem iam dicto filio de castro Ivarz». P.~839, 1249.}. В этом случае документ не отменял необходимости устного подтверждения. Напротив, право на имущество укреплялось соединением двух элементов --- завещательного текста и свидетельства о том, что соответствующее действие действительно было совершено. Иначе говоря, письменная фиксация и устное подтверждение здесь дополняли друг друга.

Если сопоставить эти случаи, становится заметно, что грамота XII~в. всё чаще выполняла двойную функцию. С одной стороны, она закрепляла сам акт распоряжения, отказа или передачи имущества. С другой --- функционировала как самостоятельное средство подтверждения этого действия в потенциальном споре. Именно поэтому составители документов придавали значение не только содержанию распоряжения, но и тому, каким образом оно было оформлено: кто его совершал, на какие прежние документы можно было сослаться, кто мог подтвердить совершённое действие и в каком виде оно могло быть предъявлено в суде.

Похожее перераспределение функций видно и в более поздних текстах. В «Кутюмах Перпиньяна» статья LXII устанавливала, \emph{«что истец не обязан подробно излагать основание иска или требования, если он требует уплаты денежной суммы на основании публичного документа, составленного в писцовой конторе Перпиньяна»}\footnote{Costums de Perpignan. [LXII]: «…ut nemo agens teneatur in tota causa exprimere causam seu peticionis, dum tamen eam de peccunia sibi solvenda faciat cum instrumento publico facto in scribania Perpinyani». P.~584--585.}. Здесь документ в известной мере уже заменял развёрнутое словесное изложение: он сам по себе служил достаточным основанием для дальнейшего движения дела. Тем самым письменный инструмент брал на себя часть той функции, которая в более раннем материале ещё принадлежала устному заявлению.

На этом фоне материал «Обычаев Тортосы» позволяет увидеть момент перехода особенно ясно. В статье 6.4.26 прямо оговаривалось, что завещание не может считаться законным, если существует только в устной форме. Для того чтобы грамота с завещанием обладала силой, она должна была быть записана, а если завещатель не умел писать, следовало обратиться к городскому писцу; далее требовались подписи свидетелей и самого завещателя. При этом текст приводил стандартную латинскую формулу свидетельской подписи: \emph{«Я такой-то по просьбе такого-то подписываюсь свидетелем и подпись свою ставлю»}\footnote{Costums de Tortosa. 6.4.26: «Et con testament o derrera volentat és feit, Axí con desús és dit, la subscription dels testimonis deu dir enaxí: “ego talis rogatus a tali mesubscribo pro teste et signum meum appono”». P.~314.}.

Наиболее полно это соотношение было сформулировано в королевской привилегии «Знатные мужи Барселоны признали…». В XXV статье специально оговаривалось, что нотариус мог составить завещание, находясь наедине с завещателем; после этого он должен был вызвать свидетелей и объявить при них, что составил завещание данного лица, и такой документ считался действительным, как если бы свидетели сами «слышали» завещание\footnote{Recognoverunt proceres. [XXV]: «Item, quod notarius potest facere testamentum ipso solo existente cum testatore, et quod, ipso faco sive notato in papiro, vocet testes coram quibus dicat se fecisse testamentum ipsius testatoris, et quod valet ac si audivissent testamentum ipsi testes». P.~595.}. В следующей статье дополнительно уточнялось, что завещание, при котором присутствовали два или три свидетеля, сохраняло силу и не могло быть оспорено по их числу\footnote{Ibid. [XXVI]: «Item, quod testamentum in quo sunt duo vel tres testes adhibiti, valet et non informatur racione testium». P.~595.}. Наконец, в главе XLVIII \emph{De testamento sacramentali} действительной признавалась последняя воля, совершённая \emph{«в письменном или неписанном виде»}, если свидетели затем приносили присягу, что они\emph{ «видели и слышали»}, как завещание было написано или сказано\footnote{Ibid. [XLVIII]: «Item, est consuetudo quod si aliquis fecerit testamentum presentibus testibus, in terra vel in mare ubicumque sit, in scriptis vel sine scriptis, seu suam voluntament etiam aliquo notario non presente in ipsa voluntate, verbotenus dicta vel scripta quod valeat ipsa ultima voluntas sive testamentum, dum testes qui interfuerint ipsi ultime voluntati vel testamento infra sex menses ex quo fuerint in Barchinona iurent in ecclesia sancti Iusti, super altare sancti Felicis martiris, presente notario qui talia testamenta conficit et aliis personis, quod ipse testes ita viderunt et audiverunt scribi seu dici, sicut in illa scriptura continetur sive ultima voluntate verbotenus ab ipso testatore dicta, et quod tale testamentum vocatur sacramentale». P.~599.}. Тем самым устное волеизъявление сохранялось, но его действительность всё теснее связывалась с записью, свидетельством и нотариальным подтверждением.

Таким образом, каталонский материал XII--XIII~вв. показывал не вытеснение устного волеизъявления письменным документом, а изменение самой социальной формы правового действия. Слова по-прежнему сохраняли значение при распоряжении имуществом, назначении наследника и отказе от права, однако их сила всё чаще зависела от того, были ли они записаны, подтверждены свидетелями и могли ли быть предъявлены в суде. Иначе говоря, речь шла не об исчезновении устной практики, а о включении её в более сложную процедуру удостоверения.

Для антропологического подхода важно, что менялся не только юридический статус акта, но и круг участников, обеспечивавших его действительность. Если раньше решающее значение имел сам момент произнесения слов, то теперь всё большую роль начинали играть писец, нотариус, свидетели, судья и сама ситуация публичного подтверждения.

В этом смысле документ не просто служил дополнительным подтверждением устного действия, а закреплял его, ограничивал возможные споры о его содержании и делал пригодным для повторного использования в суде. Поэтому соотношение речевой формулы и письменного удостоверения в каталонском праве XII--XIII~вв. следует понимать как переход от разового словесного акта к процедуре его фиксации и признания.

\subsection{Глаголы речи и стилистическая организация прямой речи в правовом тексте}

Сопоставление разных типов источников показывает, что выбор глаголов речи в каталонском правовом письме XII--XIII~вв. зависел не только от языка памятника, но и от того, какое именно процессуальное действие требовалось зафиксировать. В актовом материале XII~в. речь, как правило, не разворачивалась в виде полноценного диалога, а сводилась к нескольким опорным глаголам, обозначавшим позиции сторон и переход от спора к его разрешению. В одних случаях использовались нейтральные формы \emph{«говорили»} (лат. --- \emph{dicebant}) и \emph{«отвечали» }(лат. --- \emph{respondebant}), которые позволяли кратко передать столкновение притязаний, не воспроизводя высказывания целиком. В других употребляли формулы отказа и уступки, оформленные в первом лице, но уже настолько устойчивые, что их следует рассматривать как элементы письменного формуляра.

В «Обычаях Льейды» речевая организация текста была более сжатой. Показательна форма \emph{«жалуюсь»} (лат. --- \emph{conqueror}), которая обозначала формальный акт предъявления обвинения. Вокруг этого краткого высказывания, после которого сразу излагались процессуальные последствия, строилась логика данной нормы обычного права. При этом в «Обычаях Льейды» мы не находим развернутой сценической модели судебного общения, но уже выделяли юридически достаточный словесный минимум обвинения.

В «Кутюмах Перпиньяна» глагол «излагать» (лат. --- \emph{exprimere}) уже маркировал необязательное процессуальное действие, от которого можно отказаться или воздержаться при наличии надёжного письменного инструмента\footnote{Costums de Perpignan. [LXII]: «Item, statuimus ut nemo agens teneatur in tota causa \textbf{exprimere} causam seu peticionis...». P.~584.}. Следовательно, по сравнению с «Обычаями Льейды» в данном случае сфера применения устного заявления сокращалась, тяготея к письменной форме претензии.

В «Обычаях Тарреги» глагольная система отражала другой аспект правовой квалификации. В 1-й статье решающим представлялось разграничение между гражданским и уголовным порядком предъявления претензии: в гражданском деле ответчику даётся отсрочка для представления поручительства, тогда как в уголовном --- при очевидности преступления --- обеспечение доказательствами требуется немедленно\footnote{Consuetudines ville Tarege… [1]: «Quod si conqueritur aliquis de aliquo, civiliter, reus dilacionem unius diei habeat petendi fiduciam . In crirrinali autem causes, statim assecuret reus si crimen notorium fuerit et manifestum». P.~436.}. Форма\emph{ «обращается с жалобой»} (лат. --- \emph{conqueritur}) сохраняла значение процессуального заявления, содержание которого зависело от квалификации тяжбы как гражданской или уголовной. Связанный с действиями, включавшими прямую речь, глагол в данном случае встраивался в более сложную систему процессуальных различений.

Зачастую прямая речь могла вводиться глаголами или глагольными конструкциями, которые во многом влияли на восприятие такой вставки «читателем» и отражали интенцию составителя кодекса. Из 99 случаев употребления прямой речи в обеих версиях «Обычаев Тортосы» в 86 (64 случая в первой версии «Обычаев Тортосы» и 22 случая --- во второй версии) мы столкнулись с наиболее распространённым глаголом речи \emph{“dir”} в разных его формах. Этот глагол наиболее универсален, поскольку может в различных ситуациях вводить как слова должностного лица, так и непосредственно участников тяжбы. Например, в обычае 2.18.2 мы встречаем следующий пассаж: \emph{«… Ответчик может сказать следующее и поступить так, как захочет: “не хочу ничего доказывать…, хочу очиститься посредством клятвы”»}\footnote{Costums de Tortosa. 2.18.2: «...lo demanador pot dir e fa si·s vol: “no y vul ren provar, vàgia enant, e escondesca-m'o ab sagrament”...». P.~127.}.

В других случаях глаголом \emph{«}\emph{сказать}\emph{»} (кат. --- \emph{dir}) вводилась речь судей или должностных лиц: «\emph{Упомянутые судьи должны сказать вегеру…»}\footnote{Ibid. 1.3.14: «...os ditz jutges deven dir al veger enaxí...». P.~31.}. В то же время в «Обычаях Тортосы» присутствует ряд других глаголов, которые могут вводить прямую речь обеих сторон процесса и городских официалов: это глаголы общего значения\emph{ }\emph{«отвечать» }(лат. ---\emph{ respondere), «спрашивать, задавать вопросы»} (кат. \emph{--- demanare)}\emph{.}

Обращает на себя внимание использование глаголов, передающих прямую речь в описании процессуальных действий на судебном заседании. С одной стороны, использование общеупотребительной лексики могло быть попыткой упрощения сложной структуры правового текста. Это объяснение представляется довольно логичным, поскольку нормативный текст включал множество синтаксически сложных и насыщенных профессиональной лексикой предложений. С другой стороны, упрощение как практика работы с текстом было не очень распространено, особенно по отношению к текстам нормативного характера. Кроме того, в иных текстах нормативного характера, созданных с опорой на тексты «Барселонских обычаев» или «Фуэро Валенсии»\footnote{\emph{Cerdá Ruiz-Funes J.} La “Inquisitio” en los Furs de Valencia y en el “Llibre de las costums” de Tortosa // Anuario de Historia del Derecho Español. 1980. P.~563--586.}, такое обильное использование прямой речи не является распространённым и применяется в основном по отношению к нормам, регулирующим принесение судебной клятвы.

В королевской привилегии «Знатные мужи Барселоны признали…» организация речевого действия выглядит иначе. Здесь прямая речь встречается реже, зато особенно заметной становится связка глаголов восприятия и фиксации: свидетели должны подтвердить, что они «видели и слышали», как завещание было написано или сказано. В главе о составлении завещания\emph{ }действительной признаётся последняя воля, совершённая письменно или устно, если свидетели затем приносят присягу, что они «\emph{видели и слышали то, как оно было произнесено»}\footnote{Recognoverunt proceres. [XLVIII]: «…quod ipse testes ita viderunt et audiverunt scribi seu dici, sicut in illa scriptura continetur sive ultima voluntate verbotenus ab ipso testatore dicta, et quod tale testamentum vocatur sacramentale». P.~599.}. В таком контексте глагол “\emph{dicere”} уже не вводит реплику участника напрямую, а включается в формулу последующего подтверждения. Для языка памятника важно не столько воспроизвести слова завещателя, сколько обозначить, что эти слова были услышаны, а значит, могут быть удостоверены. Тем самым речевой акт здесь предстает не как типовая формула, а как объект памяти, свидетельства и нотариального подтверждения.

Наконец, в «Обычаях Орты» и «Обычаях Миравета» наблюдалось дальнейшее уплотнение процессуального языка. В первом из упомянутых кодексов в статье об уголовном обвинении специально оговаривается, что при возбуждении дела не требуются письменная запись и иная юридическая формальность; достаточно, чтобы обвинитель или доносчик сказал установленную формулу. В «Обычаях Миравета», напротив, акцент переносится на глагол\emph{ «называть» }(лат. ---\emph{ nominare)}\emph{\textsuperscript{ }}\footnote{Constitutiones Baiule Mirabeti. 77~:~ «~ltem si aliquis querelam preposueril de illo, debet personam \textbf{nominare}\textbf{ }de qua conqueritur et rem quam petit et qua ratione~». P.~21.}. Здесь закон уже не воспроизводит собственно реплику, а предписывает минимальный набор элементов, которые должны быть словесно обозначены. В результате и эти записи обычного права демонстрируют более сжатую модель, чем та, что запечатлена в «Обычаях Тортосы». Прямая речь либо сохраняется в единственной формуле, либо вовсе уступает место указанию на то, что именно должно быть названо.

В целом это сопоставление позволяет увидеть, что глаголы речи в каталонских правовых текстах выполняли не столько стилистическую, сколько структурную функцию. В грамотах они обозначали поворотные моменты спора и публичность действия; в «Обычаях Льейды» и «Обычаях Орты»--- минимальную и юридически достаточную форму обвинения; в «Кутюмах Перпиньяна» --- границу между тем, что ещё требуется выразить словами, и тем, что уже подтверждается документом; в «Обычаях Тарреги» --- процессуальную квалификацию жалобы; в «Обычаях Тортосы» --- нормативную организацию целых речевых сцен. В то же время в королевской привилегии «Знатные мужи Барселоны признали…» (Recognoverunt proceres) --- последующее удостоверение сказанного и написанного. В «Обычаях Миравета» --- обязанность назвать стороны, предмет и основание спора. Поэтому преобладание нейтральных глаголов типа \emph{dir}, \emph{dicere}, \emph{respondere}, \emph{conqueror} или \emph{nominare} следует объяснять не упрощением текста со стороны составителей и не ограниченными познаниями латинского языка, а стремлением придать речевому действию процессуальную определённость и встроить его в устойчивую юридическую форму.

\textbf{Выводы к параграфу §2.1}\textbf{.}

Рассмотрев ряд случаев применения прямой речи, встречающихся в сборниках каталонского обычного права, подведём итоги. Во-первых, несмотря на довольно большое число статей с элементами прямой речи (значительно большее, чем, например, в «Барселонских обычаях»): казуальные диалоги в рамках процессуальных действий, частое использование глаголов речи, большинство из приведённых «устных» элементов имеют прескриптивный характер, отражают стандартные, регламентированные правилами фразы и действия участников определённых правовых событий, будь то заседание суда или передача имущества.

Во-вторых, подавляющее большинство таких потенциально «устных» элементов записаны на латинском языке. В старокаталанских вариантах они даны лишь в «Обычаях Тортосы»; такие вставки с прямой речью действительно могли произноситься в определённых ситуациях, но в то же время большинство из них подвергались предварительной обработке и латинизации со стороны писца или нотария.

Что касается случаев, когда легисты обращались к использованию прямой речи в качестве средства передачи тех или иных правовых норм, они представлены по большей части нормами гражданского судопроизводства и гораздо реже касаются судебных процедур и дознания по уголовным делам. Таким образом, использование прямой речи в тексте «Обычаев Тортосы» можно считать свидетельством устной коммуникации лишь частично с учётом степени формализации при включении в текст законодательного памятника. Прямая речь включается легистами в запись обычного права через речевые формулы, диалогические реплики сторон и формулы вызова, произносимые судебными должностными лицами, которые являются литературно обработанными и даже стилизованными образцами речевых актов. Эти включения одновременно решают несколько задач: фиксируют перформативные высказывания, от которых зависела юридическая действительность; нормируют «правильные слова» для практикующих судей, обвинителей, нотариусов и др.; придают тексту авторитет посредством «голоса» власти и общины; усиливают дидактическую и мнемоническую выразительность кодекса. Материал «Обычаев Тортосы» показывает концентрацию прямой речи в пограничных моментах процесса (вызов и явка, признание и отказ, принесение присяги, оглашение приговора) и в казусах, где она моделирует переговорную динамику и иерархию ролей, превращая устные практики в нормализованное предписание. Следовательно, как в записях муниципальных обычаев, так и в актовом материале прямая речь предстает и как свидетельство живой судебно-устной культуры, и как сознательный литературно-юридический приём текстуализации.

\section{Иноязычные элементы в латиноязычных сводах местного обычного права и актовом материале}

В настоящем параграфе основным источником анализа выступает латинский текст «Барселонских обычаев» как наиболее цельный и внутренне связный памятник каталонской правовой культуры, позволяющий проследить способы включения иноязычных элементов в нормативный текст, их функции, степень морфологической и графической адаптации, а также случаи социальной маркировки разговорного словоупотребления. Именно данный свод даёт возможность рассматривать подобные явления не как изолированные лексические факты, а как часть более широкой языковой стратегии составителя, работавшего на стыке латинской письменной традиции и живого романского словоупотребления.

По этой причине материал «Барселонских обычаев» в данном разделе рассматривается как основной, но не как единственный. Для проверки типичности выявленных явлений, уточнения семантики отдельных терминов и определения степени их распространённости в юридической письменности привлекаются также другие латинские своды обычного права и, в необходимых случаях, актовый материал XII~в. При этом такое сопоставление не предполагает механического уравнивания источников различных жанров: если своды обычного права отражают сознательную редакторскую обработку и позволяют судить о принципах юридической фиксации иноязычной лексики, то грамоты и иные документы дают возможность проверить, в какой мере соответствующие термины и модели их включения были укоренены в повседневной правовой практике. Тем самым анализ строится по концентрической модели: «Барселонские обычаи» составляют исследовательское ядро, тогда как остальные источники служат сравнительным и верификационным фоном.

Поскольку «Барселонские обычаи» дают возможность наблюдать не только сами иноязычные элементы, но и способы их юридико-текстовой обработки, отправной точкой анализа должно стать рассмотрение механизмов их введения в латинский текст. Такой подход позволяет прежде всего выявить, в каких случаях составитель прямо маркирует нелатинское слово как принадлежащее сфере разговорного употребления, а в каких --- включает его в правовой контекст без специального пояснения, ограничиваясь его графической, морфологической или синтаксической адаптацией. Именно поэтому первый аспект, требующий специального рассмотрения, связан со способами введения иноязычных элементов в латинский правовой текст.

\subsection{Способы введения иноязычных элементов в латинский правовой текст}

В латинском тексте «Барселонских обычаев» локально маркированные языковые элементы вводятся в правовой текст не случайно, а по вполне различимым моделям. Наиболее явной из них является специальная пояснительная формула, когда составитель не просто употребляет нужное слово, но и заранее маркирует его как принадлежащее иной речевой сфере, подробнее таких формул мы коснёмся в §2.2.2. Так, в 93 статье «Барселонских обычаев» говорится: \emph{«…при помощи осадных орудий, которые крестьяне называют фунибула, госка и гатта}»\footnote{Usatges de Barcelona. Us. 70(75): «...cum ingeniis, quod rustici dicunt funibula, gosza et gatta …». P.~112.}. Уже сама конструкция \emph{“quod … dicunt”} показывает, что составитель осознаёт особый статус этих слов: они не растворены в латинской ткани текста, а введены в неё с оговоркой, как обозначения, существующие вне книжной нормы. При этом внутри той же фразы заметна разная степень формальной адаптации: \emph{funibula} уже оформлено по латинской морфологической модели, тогда как \emph{gosza} и \emph{gatta} сохраняют более непосредственную связь с местным словоупотреблением.

Та же техника фиксируется и в тексте 117-й статьи: «\emph{Если же крестьянин найдёт золото или серебро, то есть то, что в просторечии называется }\emph{бонас (кат. }\emph{bonas}\emph{)}\emph{…»}\footnote{Ibid. Us. 94(117): «Rusticus vero si inuenerit aurum vel argentum, quod vulgo dicitur bonas…». P.~130.}. И здесь составитель, предварительно маркируя местный термин как внешний по отношению к языку правовой фиксации элемент, включает его в текст нормы.

В то же время «Барселонские обычаи» знают и другую модель --- непосредственное включение местного слова в латинскую синтаксическую конструкцию без специальной пояснительной рамки. Особенно отчётливо это видно в статьях об оскорблениях. Так, в одной из них читаем: \emph{«если кто-либо станет попрекать крещёного иудея или сарацина их прежним законом или назовёт их вероотступником (кат. tressalidz) либо ренегатом…»}\footnote{Ibid. Us. 70(75): «Si quis Iudeo vel sarraceno baptizatis retraxerit illorum legem vel appellauerit eos tressalidz vel renegats…». P.~108.}. Здесь слова «перебежчики» (кат. --- \emph{tressalidz}) и «вероотступники» (кат. --- \emph{renegats}) были включены непосредственно в латинскую фразу без формулы «называется» (лат. --- \emph{dicitur}). Правовой текст не поясняет их, исходя из их понятности и значимости для адресата; это позволяет предположить, что для составителя было важно зафиксировать оскорбительное обозначение именно в той форме, в какой оно функционировало в живой речевой практике. Тем самым перед нами уже не прокомментированная «цитата» из живой речи, а прямое встраивание местных лексем в нормативную формулу.

Тот же приём обнаруживается и в 90 статье, где среди наказуемых словесных оскорблений фигурирует формула \emph{«назовёт кого-либо рогоносцем»}\footnote{Ibid.: «... uel appellauerit aliquem cuguz...». P.~108.}. В этом случае составитель вовсе не считает нужным подготовить читателя к появлению такого слова: оно включено в текст так, как если бы само по себе уже обладало достаточной юридической определённостью. Именно подобные случаи особенно показательны для анализа: латинский правовой текст принимает местную словоформу почти в готовом виде, не снимая с неё её нелатинского облика, но и не вынося её за пределы нормы.

К той же модели следует отнести и употребление термина \emph{exovar} (старокат. --- «приданое»). Он входит в латинский текст «Барселонских обычаев» без замещения классическим \emph{dos}. Особое значение в данном случае имеет сам характер заимствования и показательный выбор составителя: вместо традиционного латинского обозначения сохраняется местный правовой термин. Тем самым локальная форма получает в латинском тексте терминологически значимый статус, а не выступает случайным вкраплением. Такое решение свидетельствует о том, что составитель ориентировался на книжную норму в сочетании с реальной юридической практикой, где соответствующее обозначение, по-видимому, уже успело укорениться.

Каталаноязычные включения присутствуют и в королевской привилегии «Знатные мужи Барселоны признали…», например, в 58-й статье: \emph{«Также, каждый вправе сооружать “пристройку}\footnote{\textbf{Атанс}\textbf{ }(кат. atanç, atans) -- пристройка к стене. (См. ATANÇ, [ATANS]. // Vocabulari de la llengua catalana medieval de Lluís Faraudo de Saint-Germain. URL: \url{https://www.iec.cat/faraudo/results.asp} (поиск по слову: atanç; дата обращения: 14.04.2026).}\emph{ вдоль или поперёк” к соседской стене, если она не мешает окнам соседа, которые находятся там уже тридцать лет или [наличие которых.---А.К.] подтверждено документом»}\footnote{Recognoverunt proceres. [LVIII]: «Item, quod quilibet potest habere “atans per larc e per traves” in pariete vicino, sine impedimento lucernarum vicini que erunt ibi per triginta annos, vel quod habeat ipsas lucernas cum instrumento». P.~601.}. В данном случае мы, очевидно, имеем дело с обозначением местной хозяйственной реалии, название которой составители кодекса стремились привести в оригинальном виде. Этот термин в таком же виде встречается в 64-й статье «О пристройках к городским стенам Барселоны»: \emph{«Также постановлено, что никто не вправе возводить пристройки (кат.}\emph{~}\emph{atans}\emph{) к городской стене Барселоны, иначе как посредством “бурсеги” (кат.}\emph{~}\emph{burcega}\emph{)}\footnote{\textbf{Бурсега} (кат. burcega, burcegnus) -- строительный термин, значение которого не до конца изучено. Вероятно, означает особый вид стены или кладки. Ср.: BURCEGNUS // \emph{Du Cange et al.} Glossarium mediae et infimae latinitatis. URL: \url{https://ducange.enc.sorbonne.fr/BURCEGNUS} (дата обращения: 14.04.2026); burcega // Diccionari català-valencià-balear. URL: \url{https://dcvb.iec.cat/results.asp?word=burcega} (дата обращения: 14.04.2026).}\emph{ и лишь при согласии того, кому принадлежит эта стена»}\footnote{Recognoverunt proceres. [LXIV]: «Item, quod nemo habeat “atans” in muro civitatis Barchinone, nisi cum pa- riete burcega, nisi faciat cum voluntate eius cuius est murus». P.~602.}. В данном случае специальными маркировками данные хозяйственные термины не выделялись, они адаптировались и органично включались в латинский текст.

В «Барселонских обычаях» составитель чаще останавливался на слове, как бы выводил его на поверхность текста и маркировал его как принадлежащее иной речевой среде. В других кодексах и особенно в грамотах XII~в. та же по типу лексика уже не требовала подобного комментария. Следовательно, можно говорить не об одном, а как минимум о двух базовых способах введения локально маркированного элемента в латинский правовой текст: во-первых, через метаязыковую формулу, во-вторых, через прямое немаркированное включение в саму норму или в деловую запись.

Таким образом, правовая культура Каталонии XII--XIII~вв. демонстрирует по меньшей мере две устойчивые техники обращения с нелатинским или локально окрашенным словом. В одних случаях составитель специально оговаривает, что перед ним обозначение, принадлежащее иной речевой практике, в других --- вводит такую единицу непосредственно, как уже привычный термин нормы или акта. Именно это сочетание особенно важно для дальнейшего анализа: оно показывает, что латинский юридический текст не был замкнутой и однородной системой, а постоянно взаимодействовал с живым локальным узусом, принимая его элементы внутрь собственных нормативных формул почти без остатка.

\subsection{Социальная маркировка терминов вернакуляра в латинском тексте}

Особый интерес представляет не только сам факт включения вернакулярных элементов в латинский правовой текст, но и сам способ их введения. В ряде случаев составитель прямо указывал на ту социальную среду, в которой употреблялся тот или иной вернакулярный элемент, либо, напротив, обозначал его как общеупотребительное наименование. Тем самым лексика, используемая крестьянами или горожанами, получила в тексте дополнительную характеристику: она была напрямую связана с определённым кругом носителей или с живой практикой повседневного употребления. Такая маркировка особенно важна для историка, поскольку позволяет увидеть, каким образом автор правового памятника осмыслял соотношение книжного языка права и локальной речевой среды, а также какие социальные различия он считал значимыми при фиксации нормы.

Особенно наглядно такая маркировка проявляется в статье 94 (113) «Барселонских обычаев», где вернакулярный термин вводится как распространённое в устной речи обозначение: \emph{«…Если же крестьянин найдёт золото или серебро, которые в народе называют бонас»}\footnote{Usatges de Barcelona. Us. 90 (117): «Rusticus vero si inuenerit aurum vel argentum, quod vulgo dicitur bonas…». P.~130.}. Слово «крестьянин» (\emph{rusticus}) обозначает социальный статус участника, а формула «в народе называют» (\emph{vulgo vocatur}) характеризует употребление слова «бонас» (\emph{bonas})\footnote{\textbf{Бонас }(лат. bonas)--- локальное «народное» обозначение клада Ср.: BONAS // \emph{Du}\emph{ }\emph{Cange}\emph{ }\emph{et}\emph{ }\emph{al}\emph{.} Glossarium mediae et infimae latinitatis. URL: \url{https://ducange.enc.sorbonne.fr/BONAS} (дата обращения: 14.04.2026).}. Вероятно, составитель воспринимал его как принятое в обиходной речи локальное наименование, требовавшее пояснения в латинском тексте. При этом статья не даёт оснований связывать термин исключительно с речью крестьян: маркировка здесь имеет прежде всего узуальный характер.

Этот выбор показывает, что языковое сознание составителя было чувствительно к внутренней неоднородности самой разговорной среды. В одном случае он связывает слово с крестьянским словоупотреблением в более узком смысле, в другом --- с более широким «народным» обиходом. Напротив, сам текст позволяет увидеть внутри неё определённую стратификацию. Это особенно существенно для истории правовой культуры, поскольку показывает, что составитель осмыслял языковую вариативность не только как языковую, но и как социальную проблему.

Рассмотренная выше модель введения терминов на местном наречии в текст на латинском языке с указанием на социальную принадлежность носителей или в обобщённой форме не является уникальной для «Барселонских обычаев». В сборнике грамот «Правосудие и урегулирование конфликтов в Средневековой Каталонии» за XII~в. она встречается как минимум 8 раз. Так, в грамоте №19 упомянута \emph{«церковь блаженного Мартина, которая в народе называется Баамоль»}\footnote{Justícia i resolució de conflictes… 2024. №~19: «...ecclesia Beati Martini, \textbf{quod vulgo dicitur Baamol}». P.~24.}. В следующей грамоте с латинизированным названием этой же церкви --- «Бавамоль»\footnote{Ibid. №~20: «...et aecclesiam Sancti Martini in valle Stacionense, \textbf{quem dicitur Bavamolle}...». P.~25.} снова появляется аналогичная формула: \emph{«и на Петра Романс церковь Святого Стефана, которую в народе называют Петрабруна»}\footnote{Ibid. №~20: «...et in Petra Romans, \textbf{que vulgo dicitur Petrabruna}, aecclesiam Sancti Stephani». P.~25.}. В этих примерах упоминание о народном наречии выступает как средство точной локализации агиотопонимов\footnote{\textbf{Агиотопоним} --- топоним, образованный от имени святого, например, название церкви, часовни или монастыря. См.: Агиотопоним // \emph{Подольская Н.В.} Словарь русской ономастической терминологии / 2-е изд., перераб. и доп. М.: Наука, 1988. С.~136.}, важных для разрешения конкретной тяжбы.

В дальнейшем функция таких конструкций существенно расширилась. В грамоте №211 формула уже относится не к агиотопониму, а к антропонимам: \emph{«… со своими сыновьями Гильельмом и Пётром, которых называют Блигариями}»\footnote{Justícia i resolució de conflictes… 2024. №~19: «…cum filiis meis Guilielmo et Petro, \textbf{qui dicuntur Bligarii}…». P.~242.}. Здесь вернакулярная номинация служит дополнительным социальным и семейным идентификатором для конкретных лиц.

В грамоте №235 составитель ввёл в латинский текст имена нарицательные на местном наречии, специально помечая их как употребимые в народе: «\emph{…берега и батедó}\emph{c}\footnote{\textbf{Батед}\textbf{ó}\textbf{с} (кат. batedós) --- народное обозначение сукновальных мельниц (валялен); ср. batedor / batador // Vocabulari de la llengua catalana medieval de Lluís Faraudo de Saint-Germain. URL: \url{https://www.iec.cat/faraudo/results.asp?ActualPage=1&Id=3776688&Word=Llarg\%2C+\%5Blarch\%5D&optCriteria=1&optSearchType=2&search=balena} (дата обращения: 14.04.2026); molí draper // Gran Enciclopèdia Catalana. URL: \url{https://www.enciclopedia.cat/gran-enciclopedia-catalana/moli-draper} (дата обращения: 14.04.2026); batan // Gran Enciclopèdia Catalana. URL: \url{https://www.enciclopedia.cat/gran-enciclopedia-catalana/batan-0} (дата обращения: 14.04.2026).}\emph{, как их называют в народе…»}\footnote{Justícia i resolució de conflictes… 2024. №~235: «...las riberas y los batedós \textbf{que vulgo nominantur}...». P.~271.}. Таким образом, формула применялась гораздо шире, чем просто ономастический маркер.

Особенно показателен в этом отношении более поздний материал. В грамоте №478 читаем: \emph{«в месте, которое называется Сьюрана, а ныне и в народе именуется Вилет»}\footnote{Ibid. №~478: «…in loco ubi dicitur Siurana, qui nunc et vulgo appellatur Vilet». P.~583.}. Здесь вернакулярная формула уже прямо указывала на актуальный для носителей старокаталанского топоним, противопоставленный более старому названию. Наконец, в грамоте №~661 подобные конструкции стали частью развернутого описания имуществ: \emph{«места, которые в народе называют Гуардиолес и Менлеу»}\footnote{Ibid. №~478: «…in locis que vulgo appellantur Guardioles et Menleu…». P.~798.}. В результате такие формулы оказались встроены в латинский документ в качестве полноценного инструмента правовой индивидуализации вещей и лиц.

Формулы «в народе называется» (\emph{vulgo dicitur}), «в народе именуется» (\emph{vulgo appellatur}), «называют» (\emph{dicunt}), «именуется» (\emph{nominatur}) и «которые называются» (\emph{qui dicuntur}) позволяют проследить включение местного словоупотребления в латинский текст каталонских памятников XII~в. Обозначения «крестьян» (\emph{rustici}) и «народа» (\emph{vulgus}) уточняли область распространения отдельных слов и связывали их употребление с определённой социальной средой. Тем самым социальная маркировка позволяла не просто зафиксировать локальную словоформу, но и определить её место в языке правового памятника.

Вместе с тем местные слова могли включаться в латинский текст и без социальной маркировки. Такой случай представлен, например, обозначениями «перебежчиков» (\emph{tressalidz}), «вероотступников» (\emph{renegats}) и «рогоносцев» (\emph{cuguz}) в статье 70 «Барселонских обычаев» об оскорблениях и диффамации\footnote{Usatges de Barcelona. Us.70 (75): «...vel appellauerit eos tressalidz vel renegats <…> uel appellauerit aliquem cuguz...». P.~108.}. Социальное происхождение говорящих здесь специально не обозначается: составитель фиксирует конкретные термины как реально существующие и значимые для правовой квалификации словесного действия. Сопоставление этой нормы со статьями 73 и 90 позволяет заключить, что социальная маркировка не являлась обязательным условием включения локального слова в текст «Барселонских обычаев». Она появлялась там, где составитель считал необходимым обозначить социальную среду, с которой связывалось данное словоупотребление. Избирательность этого приёма свидетельствует о его смысловой значимости.

Рассмотренные формулы показывают, что в «Барселонских обычаях» и актовом материале происхождение отдельных слов могло связываться с социальной средой их употребления, особенно при обозначении конкретных объектов и действий. Они отражают способы осмысления языкового разнообразия в рамках общественной иерархии.

\subsection{Военная и инженерная лексика в системе иноязычных включений}

Военная и инженерная лексика образует в системе иноязычных включений особую группу. В отличие от слов, связанных с социальной характеристикой носителей или с общими формулами народного именования, здесь речь идёт прежде всего о терминах, обозначающих конкретные формы вооружённого действия, элементы осады, строительства и хозяйственной практики. Их появление в латинском тексте обусловлено, по-видимому, прежде всего предметной точностью: составитель сохраняет локальное обозначение там, где оно уже закрепилось за определённой реальностью и не имело столь же устойчивого книжно-латинского эквивалента. Поэтому анализ данной группы позволяет увидеть, как в юридический текст включались элементы языка практики, тесно связанного с военной и материальной стороной жизни средневекового общества.

В военной сфере составители латинского текста «Барселонских обычаев» допускали использование местных обозначений. Это связано, по-видимому, с тем, что соответствующая лексика принадлежала не столько книжной правовой традиции, сколько языку повседневных военных практик в условиях каталонского пограничья XI--XIII~вв. Поэтому в подобных случаях составитель не стремился заменить местный термин латинским эквивалентом, а, напротив, сохранял форму, закрепившуюся в употреблении. Такое обозначение было понятно адресату нормы именно в его локальном, практически освоенном виде. В этом отношении показателен текст 6(7--8) статьи «Барселонских обычаев»: \emph{«За гайту, энкаль}\emph{с}\emph{ всадника и штурм замка надлежит возмещать [ущерб] посредством оммажа и алискары}\emph{…»}\footnote{Usatges de Barcelona. Us. 6 (7--8): «GAITA ET ENCALZ de cauallario et assalt de castello emendetur per hominiaticum et per aliscaram...$\rightarrow$ [A]guayt e encalç de cavaler e assalt de casteyl sie emenat per homanatge ho per [al]is[ca]ra». P.~56--57. В латинском тексте оммаж и алискара соединены союзом «и» (лат. --- \emph{et}), в старокаталанской версии --- союзом «или» (старокат. --- \emph{ho}).\par \textbf{Гайта }(кат. gaita, guaita) --- военный дозор или сторожевая служба. См.: GUAITA, [GUAYTA] // Vocabulari de la llengua catalana medieval de Lluís Faraudo de Saint-Germain.  URL: \url{https://www.iec.cat/faraudo/results.asp?ActualPage=7&Id=3774246&Word=Guaita\%2C+\%5Bguayta\%5D&optCriteria=1&optSearchType=2&search=ralla} (дата обращения: 15.04.2026).\par \textbf{Энкальс} (кат. encalç) --- преследование, погоня. См.: ENCALÇ, [ENCALS] // Vocabulari de la llengua catalana medieval de Lluís Faraudo de Saint-Germain.  URL: \url{https://www.iec.cat/faraudo/results.asp?ActualPage=1&Id=3767891&Word=Encal\%EF\%BF\%BD\%2C+\%5Bencals\%5D&optCriteria=1&optSearchType=2&search=cal\%EF\%BF\%BDo} (дата обращения: 15.04.2026).\par \textbf{Осада }(кат. assalt) --- нападение, натиск, приступ; в сочетании “assalt de castello” --- штурм замка. См.: ASSALT // Vocabulari de la llengua catalana medieval de Lluís Faraudo de Saint-Germain.  URL: \url{https://www.iec.cat/faraudo/results.asp?ActualPage=11&Id=3756664&Word=Assalt&optCriteria=1&optSearchType=2&search=sego} (дата обращения: 15.04.2026).\par \textbf{Алискара} (лат.aliscara) --- особая позорящая санкция; обозначающие тяжёлое и унизительное взыскание. См.: HARMISCARA // \emph{Du Cange et al.} Glossarium mediae et infimae latinitatis / Éd. augm. Niort: L. Favre, 1883--1887. URL: \url{https://ducange.enc.sorbonne.fr/HARMISCARA} (дата обращения: 15.04.2026); ARMISCARA // Du Cange et al. Glossarium mediae et infimae latinitatis. / Éd. augm. Niort : L. Favre, 1883--1887. URL: \url{https://ducange.enc.sorbonne.fr/armiscara} (дата обращения: 15.04.2026).}. Здесь все три первых термина даны практически в неизменных старокаталанских формах внутри латинского предложения и обозначают различные виды вооружённого действия --- дозор, преследование и штурм укреплённого пункта. Это позволяет увидеть в них не случайные вкрапления разговорной лексики, а устойчивые элементы военно-правового словаря, сложившегося в среде знати, для которой война, осада и охрана замков были частью социальной и политической повседневности, поскольку указания на социолект крестьян, горожан или простолюдинов вообще отсутствует.

В той же норме употреблено обозначение оммажа (лат. --- \emph{hominiaticum}; старокат. --- \emph{homanatge})\footnote{Обозначение оммажа (лат. --- hominaticum, hominiaticum; старокат. --- homanatge, homenatge) восходит к слову «человек» (лат. --- homo), которое в феодальной лексике могло обозначать «человека сеньора». \textbf{Оммаж} --- церемония признания человеком своей личной связи с сеньором и к возникавшей вследствие этого зависимости. В данной статье «Барселонских обычаев» оммаж выступает формой удовлетворения за вооружённое действие. См.: \emph{Rodón Binué E.} El lenguaje técnico del feudalismo en el siglo XI en Cataluña: contribución al estudio del latín medieval. Barcelona: Consejo Superior de Investigaciones Científicas, Escuela de Filología, 1957. P.~136--138; \emph{Coderch M. }Entre l’amor i el dret: l’ús del lèxic feudal a la lírica amorosa medieval catalana i valenciana // Anuario de Estudios Medievales. 2015. Vol.~45, núm. 1. P.~220.}. Латинская форма получила соответствующее существительному среднего рода окончание и вошла в предложно-падежную конструкцию «посредством оммажа» (лат. --- \emph{per hominiaticum}). Термин был тем самым включён в грамматическую систему средневековой латыни при сохранении обозначения института, сложившегося в местной феодальной практике.

Статья 6(7--8) фиксирует несколько способов письменного оформления локальной лексики. Обозначения вооружённых действий --- гайта, энкальс и штурм (лат. --- \emph{gaita}, \emph{encalz}, \emph{assalt}) --- переданы в формах, близких к романскому употреблению. Обозначение оммажа получило латинское морфологическое оформление (лат. --- \emph{hominiaticum}). Соседство этих форм в пределах одной нормы показывает, что составители «Барселонских обычаев» могли сохранять местное слово без существенных изменений либо приспосабливать его к грамматике латинского текста. Оба способа служили письменной фиксации военно-правового словаря, сложившегося в практике отношений внутри знати.

С противоположной ситуацией мы имеем дело в уже рассмотренном фрагменте 73-й статьи «Барселонских обычаев»: \emph{«…и не осаждать [замок сеньора] с применением осадных орудий, которые крестьяне называют фунибулы, госки и гатты, поскольку это было бы великим позором для властей»}\footnote{Usatges de Barcelona. Us. 73 (93--94--95): «... nec debellare cum ingeniis, quod rustici dicunt funibula, gosza et gatta, quia magnum dedecus erit potestatibus». P.~112.\par \textbf{Фунибула} (лат. funibula/fundibula) --- осадная камнемётная машина. Ср.: FUNDIBULA // \emph{Du Cange et al.} Glossarium mediae et infimae latinitatis / éd. augm. Niort : L. Favre, 1883--1887. T.~3. URL: \url{https://ducange.enc.sorbonne.fr/FUNDIBULA} (дата обращения: 14.04.2026); FRANDEGULUM // \emph{Du Cange et al.} Glossarium mediae et infimae latinitatis / éd. augm. Niort : L. Favre, 1883--1887. URL: \url{https://ducange.enc.sorbonne.fr/FRANDEGULUM} (дата обращения: 14.04.2026).\par \textbf{Госка }(лат. gosza) --- осадная машина; по мнению \emph{К.А.~Прието} Эспиносы, термин следует интерпретировать как «осадная машина, использовавшаяся наподобие тарана». См.:\emph{ Prieto Espinosa C.A.} El lèxic dels oficis a la documentació llatina de la Catalunya altmedieval : doctoral thesis. Barcelona : Universitat de Barcelona, 2020. T.~1. P.~381.\par \textbf{Гатта} (лат. Gatta) --- осадная машина типа крытой подвижной галереи; латинские аналоги термина -- vinea или pluteus. По мнению \emph{К.А.~Прието} Эспиносы, термин следует интерпретировать как «осадная машина, использовавшаяся наподобие тарана». См.: \emph{Prieto Espinosa C.A.} El lèxic dels oficis a la documentació llatina de la Catalunya altmedieval : doctoral thesis. Barcelona: Universitat de Barcelona, 2020. T.~1. P.~381. См.: GATTA // \emph{Du Cange et al.} Glossarium mediae et infimae latinitatis / éd. augm. Niort : L. Favre, 1883--1887. URL: \url{https://ducange.enc.sorbonne.fr/GATTA} (дата обращения: 14.04.2026).}. Следовательно, речь идёт о вполне определённой социальной среде --- крестьянском, сельском населении. Локальная лексика здесь оказывается не просто внешней по отношению к книжной латыни, но ещё и социально «приземлённой»: составитель фиксирует её как речь незнатных носителей, находящихся вне круга латинской учёности и политико-правовой элиты\footnote{\emph{Zimmerman M. }La représentation de la noblesse... P.~25.}. При этом сам термин подвергается адаптации и латинизации: сохраняется классическая латинская форма множественного числа среднего рода 2-го склонения на \emph{-}\emph{a}.

При этом особенно важно, что речь идёт о военной и технической лексике, связанной с осадными орудиями. Сами реалии принадлежат сфере вооружённого конфликта и, следовательно, непосредственно касаются интересов и действий знати. Однако составитель не сообщает, какими терминами пользовались в аналогичном контексте представители высших слоёв. Он фиксирует лишь то, как эти орудия \emph{«называют крестьяне»}. Тем самым текст создаёт своеобразную асимметрию: язык социальной верхушки остаётся не названным и как бы подразумевается в качестве не нуждающейся в комментарии нормы, тогда как язык крестьян подлежит отдельному обозначению. Из этого следует, что социальная маркировка здесь служит не просто описанием речевой практики, но и способом выстраивания иерархии между различными языковыми регистрами. Следует также учитывать, что каталонское общество формировалось в условиях длительного соседства и борьбы с мусульманским югом, а сама письменная правовая традиция складывалась на фоне расширения власти, укрепления замков и освоения пограничных территорий; поэтому присутствие локальных военных терминов в правовом тексте выглядит не исключением, а закономерностью.

Таким образом, военная и техническая лексика в системе иноязычных включений занимает в «Барселонских обычаях» особое место. В отличие от тех случаев, где вернакулярное слово вводится через указание на его общеупотребительность или на социальную принадлежность носителей, здесь локальные обозначения нередко включаются в латинский текст как уже готовые и функционально точные термины. Это особенно заметно в нормах, связанных с осадой, охраной замков, преследованием противника и применением оружия: составитель сохраняет такие слова не потому, что они случайно проникли в текст, а потому, что именно в местной форме они были наиболее удобны для обозначения конкретных военных и технических реалий. Тем самым иноязычное включение в этой группе выполняет прежде всего предметно-уточняющую функцию и отражает тесную связь языка права с языком практики.

Вместе с тем рассмотренные примеры показывают, что внутри этой группы существовали разные способы легализации местной лексики в правовом тексте. В одних случаях, как в статье о «дозоре» (кат. --- \emph{gaita}), «преследовании» (кат. --- \emph{encalz}) и «нападении» (кат. --- \emph{assalt}), военные термины включались в текст без социальной маркировки, что позволяет отнести их к привычному военно-правовому словарю знати и обитателей замков. В других, как в статье об осадных орудиях, аналогичная по тематике лексика специально приписывается крестьянам, и тогда вернакулярное слово выступает уже не только как техническое обозначение, но и как социально маркированная речь. Наконец, сопоставление латинской и старокаталанской версий показывает, что переводчик мог перерабатывать и сам предметный ряд военной лексики, подбирая более близкие к повседневной практике эквиваленты. Всё это позволяет заключить, что военная и техническая лексика в «Барселонских обычаях» служит важным показателем взаимодействия латинской письменной нормы с локальным опытом войны, обороны и вооружённого конфликта, то есть с той сферой, где право особенно непосредственно соприкасалось с социальной реальностью.

\subsection{Оскорбительная лексика и конфессиональные обозначения как объект правовой квалификации}

В системе иноязычных включений в латинском тексте «Барселонских обычаев» значительное место занимает лексика, связанная с оскорблением, диффамацией и конфессиональными обозначениями. В подобных случаях слово представляет собой не только языковую единицу, но и непосредственный объект правовой квалификации, поскольку именно в словесной форме совершается действие, наносящее ущерб чести, репутации или общественному статусу лица. Этим, по-видимому, и объясняется сохранение местных обозначений в их узнаваемом виде: для составителя важна не только семантика высказывания, но и та конкретная форма, в которой оно функционировало в повседневной речевой практике. Рассмотрение этой группы позволяет проследить, каким образом правовой текст фиксировал и регулировал словесные формы бесчестия, включая те из них, которые были связаны с конфессиональными различиями и социальной напряжённостью.

Следы некнижного или даже разговорного словоупотребления обнаруживаются в тех статьях «Барселонских обычаев», где объектом правовой квалификации становилось само оскорбительное слово. В подобных случаях правовой текст фиксировал не только факт речевой агрессии, но и конкретную словесную форму, посредством которой наносится бесчестие. Показателен в этом отношении текст 70(90) статьи «Барселонских обычаев»: «\emph{Если кто-либо станет попрекать крещёного иудея или сарацина их прежним законом или назовёт их вероотступником (кат. tressalidz) либо ренегатом…}»\footnote{Usatges de Barcelona. Us. 70 (90): «“SI QVIS IVDEO uel sarraceno baptizatis retraxerit illorum legem uel appellauerit eos tressalidz uel renegats …». P.~108--109.}\emph{.}

Логично предположить, что в одной местности параллельно употреблялись обозначения «вероотступник» (кат. --- \emph{tressalidz}; лат. --- \emph{renegatus}), он был впервые зафиксирован на территории Каталонии в латинском и старокаталанском тексте в «Барселонских обычаях»\footnote{Ibid. Us. 39(42),70 (90). P.~82--83; 108--109. Renegatus // \emph{Du Cange C. et al. }Glossarium~mediae~et~infimae~latinitatis / éd.~augm., Niort~:~L.~Favre, 1883---1887, t.~7,~col.~123c. URL~: \url{http://ducange.enc.sorbonne.fr/RENEGATUS} (дата обращения: 27.03.2023); Renegat // Gran Diccionari de la llengua catalana (GDLC). URL~: \url{https://www.diccionari.cat/EC---GDLC---e00116932.xml\#} (дата обращения: 27.03.2023).}. В данном случае мы видим употребление двух практически синонимичных терминов рядом: действительно, в общем смысле их значение весьма сходно. Это может свидетельствовать о различительной способности, чувствительной к очень тонким оттенкам значений, связанных с межконфессиональным взаимодействием. Термин «еретик» (лат. --- \emph{hereticus}) мы встречаем и в тексте «Барселонских обычаев»: \emph{«…христиане и сарацины, иудеи и еретики, [чтобы государи.---А.К.] могли им доверять, вверяя им не только самих себя, но и свои города и замки}»\footnote{Usatges de Barcelona. Us. 60(64): «...christiani et sarraceni, iudei et heretici, possint se in illis fidare et credere non solum illorum personas...». P.~96--97; Hæreticus //\emph{ Du Cange C. et al. }Glossarium mediae et infimae latinitatis / éd. augm., Niort : L. Favre. 1883-1887. T.~4. Col.~155c. URL~: \url{http://ducange.enc.sorbonne.fr/HAERETICUS} (дата обращения: 27.03.2023).}\emph{.} В «Барселонских обычаях» он используется по отношению к катарам (альбигойцам) или вальденсам, а не к «новым христианам» или крещённым мусульманам и иудеям\footnote{Usatges de Barcelona. Us. 60(64). В данной статье «Барселонских обычаев» в рамках перечисления религиозных групп и видно, что «еретики» не относятся ни к иудеям, ни к мусульманам, ни к христианам, и представляют собой отдельную конфессиональную группу.}. Кроме того, именно в тот период на старокаталанском языке появился термин «еретики» (кат. --- els haeritici), который, как показывает Большая Каталонская энциклопедия, применялся в XIII~в. исключительно в отношении местных катарских общин\footnote{El catarisme als Països Catalans // Gran Enciclopédia Catalana. URL: \url{https://www.enciclopedia.cat/ec---gec---0089425.xml} (дата обращения: 19.04.2024)}. Они подвергались преследованиям, но при этом пользовались защитой северо-каталонской и окситанской знати\footnote{\emph{Ventura i Subirats~J. }El catarismo en Cataluña // Boletín de la Real Academia de Buenas Letras de Barcelona. 1960. T.~28. P.~75--158.}, «Барселонские обычаи» не сообщают ничего об их специальном преследовании со стороны властей графства Барселонского. Таким образом, эта норма свидетельствует о том, что в правовом сознании составителя «Барселонских обычаев» конфессиональная принадлежность, как сохраняемая, так и уже оставленная, оставалась значимым основанием для правовой квалификации словесного оскорбления. Соответственно, оскорбительная лексика, связанная с религией, осмыслялась как затрагивающая не только личную честь, но и социальную идентичность человека.

В городском пространстве в рамках христианского общества оскорблением считалось слово\emph{ “}\emph{cuguz}\emph{” --- }рогоносец\footnote{\textbf{Рогоносец}\textbf{ }(кат. \textbf{c}\textbf{uguz, cuguç)} --- оскорбительное обозначение мужа, жена которого была уличена в супружеской неверности, то есть «рогоносца»; ср. cuguç // Vocabulari de la llengua catalana medieval de Lluís Faraudo de Saint-Germain.  URL: \url{https://www.iec.cat/faraudo/} (поиск по слову: cuguç; дата обращения: 15.04.2026);: cuguç // Diccionari català-valencià-balear.  URL: \url{https://dcvb.iec.cat/results.asp?word=cugu\%25u00e7} (дата обращения: 15.04.2026).}. Обратимся снова к 70(90) статье «Барселонских обычаев»\footnote{Usatges de Barcelona. Us. 70(90). P.~108; Cuguç // \emph{Coromines J. }Diccionari etimològic i complementari de la llengua catalana. Barcelona: Curial Edicions Catalanes, 1981. Vol.~2. P.~1086.}\emph{.} Сравним эту статью в латинской и старокаталанской версиях «Барселонских обычаев»:

\emph{Таблица }\emph{X}\emph{. Сопоставление латинского текста и старокаталанских версий статьи 70 (75) «Барселонских обычаев»}

\begin{center}
\small
\setlength{\tabcolsep}{2pt}
\begin{longtable}{|p{\dimexpr0.94\textwidth/4\relax}|p{\dimexpr0.94\textwidth/4\relax}|p{\dimexpr0.94\textwidth/4\relax}|p{\dimexpr0.94\textwidth/4\relax}|}
\hline
\textbf{Латинский текст} & \textbf{Старокаталанский текст по рукописи K, изданной Ж. Бастардесом-и-Парерой} & \textbf{Старокаталанский текст Викской рукописи, изданной Ж. Гудиолем-и-Кунилем} & \textbf{Русский перевод} \\ \hline
\emph{…uel si quis infra menia nostra uel burgos primus eduxerit gladium contra alium uel appellauerit aliquem cuguz…} & \emph{…ho si negú dins los murs nostres primerament traurà coltel contra altre o l’apelarà cuguç…}\footnote{Usatges de Barcelona. Us. 70(90). P.~108--109.} & \emph{…o si algun enfre els murs nostres o els burcs primerament traura coltel a altre o apelara altre caguz…}\footnote{\emph{Gudiol i Cunill J. }Traducció dels Usatges, les més antigues constitucions de Catalunya y les Costumes de Pere Albert // Anuari de l’Institut d’Estudis Catalans. 1907. Vol.~1. P.~295.} & «…или если кто-либо в пределах наших стен \textbf{или в бургах}\textbf{ (кат1. -- опущ.)} первым обнажит \textbf{меч (кат. нож --- coltel)}\footnote{Об этом разночтении см. §2.4.4.}\textbf{ }против другого либо назовёт другого рогоносцем… \\ \hline
\end{longtable}
\end{center}

Сопоставление двух старокаталанских версий показывает вариативность передачи пространственной формулы. В рукописи K, текст которой был опубликован Ж.~Бастардесом-и-Парерой, слова «или в бургах» (лат. --- \emph{uel burgos}) опущены: \emph{«если кто-либо в пределах наших стен…»}. В то же время в тексте другой рукописи, опубликованной Ж.~Гулиолем-и-Кунилем, сохранились оба элемента латинского текста. Опущение относится, таким образом, к тексту рукописи K и не характеризует старокаталанскую переводческую традицию в целом. Сочетание «в пределах наших стен или в бургах» очерчивало область действия санкции\footnote{Старокаталанское слово «бург» (старокат. --- burc, burch) могло обозначать поселение, сложившееся возле укреплённого города или замка. Поэтому слова «или в бургах» (лат. --- uel burgos; старокат. --- o els burcs) расширяли территориальный охват нормы и не являлись повторением предшествующего указания на пространство внутри городских стен. См.: \emph{Jaime Moya J.M. }El lèxic d’origen germànic en el llatí medieval de Catalunya: tesi doctoral. Barcelona: Universitat de Barcelona, 2015. P.~63--64.}. Латинская версия и текст, опубликованный Ж.~Гудиолем-и-Кунилем, расширяли её на пространство внутри укреплений и на примыкавшие к нему поселения.

В одной формуле санкции во всех версиях были соединены обнажение оружия и произнесение оскорбительного наименования --- два действия, способные нарушить порядок в плотно заселённой среде. Социальный статус акторов в данном случае статья не регулировала, однако пространственные маркеры в латинской версии позволяют отнести её к праву городскому\footnote{\emph{Варьяш~О.И.} Средневековое право: Опыт пристального чтения некоторых памятников // Пиренейские тетради: право, общество, власть и человек в Средние века. / сост. и отв. ред. \emph{И.И.~Варьяш}, \emph{Г.А.~Попова}, М., 2006. С.~181.}. Оскорбительное обозначение «рогоносец» сохранялось во всех трёх версиях в близких формах (лат. --- cuguz; кат. --- cuguç, caguz), поскольку для правовой квалификации было существенно узнаваемое слово, посредством которого совершалось правонарушение. Его употребление было засвидетельствовано и в других каталонских правовых памятниках XII--XIII~вв., включая «Поселенную хартию Тортосы» (1149~г.) и «Обычаи Тортосы» (1270-е гг.)\footnote{Cuguç // \emph{Coromines J. }Diccionari etimològic i complementari de la llengua catalana. Barcelona: Curial Edicions Catalanes, 1981. Vol.~2. P.~1086.}.

Все перечисленные оскорбления вводятся глаголом «называть» (лат. --- \emph{appellare}), что важно для нас, поскольку он передаёт функцию называния кого-либо каким-либо словом, пусть и в пейоративном значении. Эта функция называния, по мнению лингвистов, на не доминирующем языке в рамках мультилингвизма зачастую передавалась максимальным приближением к звучанию данного слова или синтагмы на разговорном языке.

В случае с «Барселонскими обычаями» мы видим, что слова, хорошо вписывающиеся в старокаталанскую орфографию, попадают в латинский текст практически без изменений (ср.~лат. “\emph{cuguz}\emph{” --- }кат.\emph{ “cuguç”, }лат. “\emph{t}\emph{ressalidz” --- }кат.\emph{ “trespassalig”, лат. “}\emph{renegats}\emph{” --- кат. “}\emph{renegats}\emph{”}). Исходя из этого факта, мы можем предположить, что для составителя данного текста было важно сохранить данные оскорбительные и унизительные обозначения именно в их старокаталанской форме, т.е. в том виде, в каком они реально могли прозвучать на территории принципата Каталония и шире Арагонской короны.

«Обычаи Тортосы» демонстрируют прямую преемственность по отношению к более раннему памятнику (первые нормы о бесчестии встречались ещё в «Поселенной хартии Тортосы» 1149~г.)\footnote{Costums de Tortosa. Carta de l'acapte de la Ciutat de Tortosa. P.~531.}. В статье 9.4.2 приводится определение вербального бесчестья: \emph{«Оскорбление --- это когда кто-то называет другого предателем, лжецом, рогоносцем, вором, клятвопреступником, блудницей или развратницей, прелюбодеем, изменником или другими подобными им словами, независимо от того, правда это или нет»}\footnote{Ibid. 9.4.2: «Injúria és qui a altre cridarà traydor o bare, cugús, ladre, perjur, putana, corregut, aültre, renegat, o altres coses semblantz a aquestes, renegat, o altres coses semblants a aquestes, jasia so que sia ver o no». P.~419.}. Оскорбления обычно вводились глаголами «называть» (кат. --- \emph{appellare}) или «публично выкрикивать» (кат. --- \emph{cridare}). Это подчёркивало важность самого речевого акта: факт произнесения вслух также имел правовое значение (как, например, в ситуации с клятвой).

В данной группе иноязычных включений в латинском тексте «Барселонских обычаев» были зафиксированы слова, значимые не только как действия, подлежащие правовой оценке. Именно поэтому составители сохраняли их в форме, максимально близкой к живому употреблению. Рассмотренные примеры показывают, что правовое значение придавалось как конфессионально окрашенным обозначениям, затрагивавшим статус и идентичность человека, так и формам бытового словесного бесчестия, характерным, в частности, для городской среды. Тем самым данный материал показывает, что каталонская правовая культура включала в себя социально узнаваемые формы речевого конфликта, превращая их в объект нормативного регулирования.

\subsection{Локальные правовые реалии и терминологические особенности их обозначения}

В системе иноязычных включений в латинском тексте «Барселонских обычаев» заметное место занимают термины, связанные с локальными правовыми реалиями. В отличие от военной, оскорбительной или социально маркированной лексики, здесь речь идёт прежде всего о таких обозначениях, которые закрепились за определёнными институтами, формами владения, имущественными отношениями или обычаями и потому не всегда имели точный и равнозначный эквивалент в традиционном латинском юридическом языке. В подобных случаях составитель, как правило, не стремится полностью заменить местное слово латинским аналогом, а сохраняет его в тексте как наиболее точное обозначение конкретной правовой практики. Анализ этой группы позволяет увидеть, каким образом латинская письменная норма включала в себя элементы местного правового опыта и как именно такие термины закреплялись в тексте --- в неизменной форме, в частично адаптированном виде или в сочетании с латинскими словообразовательными средствами.

Один из наиболее показательных примеров такого рода представляет термин «приданое»\emph{ }(кат.\emph{ --- exovar)}\emph{\textsuperscript{ }}\footnote{Usatges de Barcelona. Us. 85 (108).}, включённый в латинский текст «Барселонских обычаев» практически без формального изменения.\emph{ }Исследователи старокаталанского языка сходятся во мнении, что данный термин происходит от арабского \arabictext{الشوار} (aš- šuwâr, šawâr, šiwâr)\footnote{\emph{Reig~E.S.} E de la paraula “aixovar” // Aula de lletres valencianes. Núm. 2. 2012. P.~95; \emph{Biosca i Bas~A.} Aproximación a los arabismos en la documentación de Jaime I // IV Congresso Internacional de Latim Medieval Hispânico. Lisboa, 12--15 de outubro de 2005 : actas. Lisboa, 2006. P.~219--226; \emph{To Figueras~L. }Las Mujeres en la sociedad catalana de los siglos IX al XI.... P.~352.}, хотя А.~Биоска-и-Бас зафиксировал латинизированную форму данного термина \emph{--- }\emph{exovariu}\emph{m}, встречающуюся в документах при Жауме I Завоевателе (1213--1276)\footnote{\emph{Biosca i Bas~A}.Op. cit. P.~221.}. Однако в «Барселонских обычаях» мы такой морфологической адаптации к латинскому тексту не наблюдаем; лексическая единица арабского происхождения через народное разговорное наречие включается в неизменной форме в латинский текст правового компендиума. Д.~Кэгей и Л.~То~Фигуэрас указывали, что термин «экзовар» был доминирующим для обозначения приданого\emph{ }на территории Каталонии в X--XI~вв., а в завещаниях практически не упоминаются другие термины\emph{ }латинского происхождения\footnote{\emph{To Figueras L. }Op. cit. P.~352.}, например, традиционный латинский термин\emph{ “}\emph{d}\emph{ō}\emph{s}\emph{”} (лат. --- приданое). Как считают вышеупомянутые исследователи, впоследствии значение термина «приданое» (кат. --- \emph{exovar}) расширилось, и он стал употребляться как синоним латинского термина «приданое» (лат. --- \emph{dōs}). Термин «приданое» (кат. --- \emph{exovar}) обозначал имущественный вклад, вносившийся при заключении брака со стороны родственников будущей супруги. В его состав могли входить деньги, движимые вещи, скот, земельные владения и права на недвижимость. В отдельных документах в качестве приданого дочери передавались земли и другие недвижимые объекты. Тем самым термин указывал на правовое назначение имущества\footnote{Aixovar // \emph{Coromines J. }Diccionari etimològic i complementari de la llengua catalana. Barcelona: Curial Edicions Catalanes, 1980. Vol.~1. P.~115--116; \emph{To Figueras L. }Les fonctions de la dot et du douaire dans la société rurale de la Catalogne (Xe--XIe siècle) // Dots et douaires dans le haut Moyen Âge. Rome: École française de Rome, 2002. P.~189--217.}.

К.А.~Прието~Эспиноса различает несколько стандартных моделей, согласно которым каталанские термины могли включаться в латинский текст источников различного рода, в частности, прибавление к каталанскому термину латинских словообразовательных суффиксов\footnote{\emph{Prieto Espinosa C.A}. Op.cit.Vol.~I.P.~72.}.

Хотя в статье №~86 (109) речь идёт об имуществе бездетного крестьянина\footnote{Usatges de Barcelona. Us. 86 (109): «DE REBVS et facultatibus de exorchis pagensi: bus illis ab hoc seculo discessis, eorum seniores habeant partem illam quam debent habere in simul filii, si ibi remansissent ab exorchiis procreati». P.~122.}, оно называется термином \emph{“}\emph{riquees}\emph{”} (кат.---богатство)\footnote{Riquesa, [rriquea] s.// Vocabulari de la llengua catalana medieval de Lluís Faraudo de Saint-Germain. URL~: \url{http://www.iec.cat/faraudo/results.asp?search=honta&optCriteria=0&optSearchType=1} (дата обращения: 14.04.2023).}.

\emph{Таблица }\emph{X}\emph{. Обозначение имущества бездетных крестьян в 86(109) статье «Барселонских обычаев»}

\begin{center}
\small
\setlength{\tabcolsep}{2pt}
\begin{longtable}{|p{\dimexpr0.94\textwidth/4\relax}|p{\dimexpr0.94\textwidth/4\relax}|p{\dimexpr0.94\textwidth/4\relax}|p{\dimexpr0.94\textwidth/4\relax}|}
\hline
\textbf{Латинский текст} & \textbf{Старокаталанский текст по изданию Ж. Бастардеса-и-Пареры} & \textbf{Старокаталанский текст по рукописи, опубликованной Ж. Гудиолем-и-Кунилем} & \textbf{Русский перевод} \\ \hline
\emph{De rebus et facultatibus de pagensibus exorquiis} & \emph{De les coses e de les riquees e de les exorquies dels pageses…} & \emph{De les coses et de les riquees dels pageses exorcs qui son morts…} & Об имуществе и состоянии умерших бездетными крестьян \\ \hline
\emph{quam deberent in simul habere filii} & \emph{…que deurien aver ensemps los fils…} & \emph{…ensems que agien los hereus…} & …которую совместно получили бы сыновья (кат2.--- наследники) \\ \hline
\end{longtable}
\end{center}

При этом слово «богатства» (кат. --- \emph{riquees}) передаёт латинское обозначение «имущество, состояние» (лат. --- \emph{facultates}) и употребляется не как характеристика имущественного достатка крестьянина, а как общее наименование принадлежавшего ему имущества. Такое соответствие показывает, что в старокаталанском переводе слово \emph{riquees} могло иметь более широкое значение, чем «богатство», и обозначать~совокупность имущества, переход которой после смерти крестьянина регулировался данной нормой. Поэтому в русском переводе его точнее передавать словами «имущество» или «состояние», сохраняя буквальное значение «богатства» только при анализе терминологии старокаталанского текста. Что особенно примечательно, в данном случае составитель перечисляет «богатства» и «экзоркии»\footnote{\textbf{Экзоркия} (кат. exorquia) --- право сеньора на выморочное имущество крестьянина, умершего без наследников и само выморочное имущество. См. Eixorquia, [exorquia] s.// Vocabulari de la llengua catalana medieval de Lluís Faraudo de Saint-Germain. URL~: \url{http://www.iec.cat/faraudo/results.asp?Word=Eixorquia,\%20[exorquia]&Id=3070441&search=eixor&optCriteria=0&optSearchType=1&ActualPage=1}(дата обращения: 14.04.2023); \emph{Варьяш~И.И., Астахов~М.А.} Арагонская Корона в XV~в. \emph{// }История Испании / под ред. \emph{В.А.~Ведюшкина}, \emph{Г.А.~Поповой}. М.: «Индрик», 2012. Т.~1. С.~346.} как ряд однородных членов\footnote{См. комментарии к тексту «Барселонских обычаев»: Usatges de Barcelona / ed. por Bastardas i Parera J. Barcelona, 1991. N. 73.}. Таким образом, последний термин к XIII~в. уже прочно закрепился за земельным наделом, принадлежавшим бездетным крестьянам, и был включён переписчиком в число крестьянского имущества. В данном случае локальный термин выполняет не вспомогательную, а собственно понятийную функцию, фиксируя такую форму имущественного изъятия, которая была значима именно в каталонской юридической практике и воспринималась как самостоятельный институт местного права. В случае с экзоркией перед нами фиксация правового института, который существовал задолго до составления «Барселонских обычаев», и потому сохранился в тексте в своей локальной форме.

Таким образом, термины, связанные с локальными правовыми реалиями, занимают в системе иноязычных включений особое место, поскольку обозначают институты и категории, которые укоренены в каталонской правовой культуре и документальной практике. В подобных случаях составители латинского текста «Барселонских обычаев» сохраняли наиболее точные и узнаваемые своими современниками обозначения. Поэтому такие слова, как экзовар, экзоркия и кугусия, закономерно вошли в латинский текст «Барселонских обычаев» и несли в себе специфическое правовое содержание. Их употребление показывает, что письменная латинская норма была вынуждена учитывать и закреплять те институты местного права, которые уже обладали собственной устойчивой терминологией. Тем самым латинский текст «Барселонских обычаев» выступает не только как средство юридической фиксации, но и как пространство, в котором локальный правовой опыт получал письменное и понятийное оформление.

\textbf{C}\textbf{§2.2.6}\textbf{.}\textbf{ Ландшафтная и хозяйственная лексика в латинском тексте обычая}

Не менее показательной является и та группа иноязычных включений, которая связана с обозначением ландшафта, землепользования и хозяйственной среды. В подобных случаях латинский текст обычая сталкивается с необходимостью назвать не отвлечённую правовую категорию, а конкретный объект пространства --- участок земли, водный источник, дорогу, угодье, постройку или элемент хозяйственной инфраструктуры. Именно поэтому здесь особенно заметна опора на местное словоупотребление: локальный термин сохраняется там, где он уже был прочно связан с определённой реальностью и служил наиболее точным её обозначением. Анализ этой лексики позволяет увидеть, каким образом в правовой текст входили представления о пространстве и хозяйственном укладе, значимые для повседневной практики обычного права.

В тексте «Барселонских обычаев» встречается ещё один арабизм, гораздо лучше вписанный в латинскую грамматическую систему. Его составители латинской версии данного свода законов также по определённым причинам, которых мы коснёмся далее, предпочли латинским аналогам. Рассмотрим 69(74) -ю статью «Барселонских обычаев»: \emph{«Повелеваем, чтобы асекия, по которой вода мельниц течёт к Барселоне, оставалась неприкосновенной во все времена…»}\footnote{Usatges de Barcelona. Us. 69 (74): «CECHIAM aque molinorum que fluit ad Barchinonam mandamus esse intactam omni tempore…$\rightarrow$ La ciquia de l'ayga dels molins qui corre a Barcelona, manam éser no tocada per tots temps…». P.~106--107.}\emph{.} В данном случае нас интересует термин «оросительный канал»\emph{ }(лат. ---\emph{ }\emph{cechia}; кат. ---\emph{ }\emph{ciquia}\emph{)}, обозначающий оросительные или мельничные каналы, то есть ирригационные сооружения, снабжённые водными колёсами. Он происходит от арабского \arabictext{ساقية} \emph{(saqya, sâqiya)}\footnote{\emph{Biosca i Bas~A.} Op. cit. P.~221.}. А.~Биоска-и-Бас отмечает, что к данному термину при его адаптации к латинскому тексту в валенсийских документах практически всегда добавляли слово «вода» (лат.\emph{ --- }\emph{aqua}), хотя в большинстве случаев это было избыточным\footnote{Ibid. P.~225.}. То же самое мы видим и в обеих версиях текста «Барселонских обычаев».

Отдельный исследовательский вопрос --- почему составитель свода законов предпочёл именно этот термин, когда в то же время существовали его аналоги латинского и греческого происхождения: «канал» (лат. --- \emph{cănālis}), «водоток, ручей» (лат. --- \emph{rīvus}) и др. С последними составитель, вероятнее всего, был знаком в силу знакомства с классическими текстами. Мы имеем дело с правовым оттенком термина «оросительный канал» (лат. --- \emph{cechia}; кат. ---\emph{ }\emph{ciquia}), поскольку на территории Новой Каталонии и Валенсии так обозначали общественные ирригационные сооружения, которые не могли находиться в собственности у частного лица\footnote{Ibid.}. Латинские же термины такого значения не имели. Это очень хорошо вписывается в контекст рассмотренной статьи «Барселонских обычаев». В данном случае мы имеем дело с вполне осознанным переключением кодов и использованием именно того термина, который отражал важные для легиста правовые категории, а именно невозможность передачи объекта в собственность или владение.

С этим же термином мы встречаемся в грамоте №~666 (1184~г.), посвящённой урегулированию земельного спора супругов де Пужоль с братьями Гильемом и Бернатом де Копонсу : «\emph{…и сказали, что мы, Пётр де Подиоло и моя жена Эсклармунда, должны иметь полностью и в целости, как наше собственное наследственное имущество, [землю] в этих границах: вниз, по направлению к мансу Гильельма де Сант-Кугата, до асекии, которая соединяется с другим нашим аллодом…}»\footnote{Justícia i resolució de conflictes…2024. №~666: «...et dixerunt quod nos Petrus de Podiolo et uxor mea Esclarmunda de istis terminis inferius versus mansum Guillelmi de Sancto Cucuphate usque ad cechiam que se continet cum illo alio alodio nostro...». P.~806.}.

Ещё одним проявлением мультилингвизма в латинском тексте «Барселонских обычаев» является включение в него ряда хозяйственно-географических терминов не автохтонного происхождения, а заимствованных из старопровансальского диалекта окситанского языка, что говорит об активных языковых контактах с этим регионом и использовании ряда терминов, которые точнее описывают специфику региона, чем их классические латинские эквиваленты. Здесь, как и в предыдущем случае, мы сталкиваемся с тем, что составитель почему-то не воспользовался латинской терминологией, которая существовала.

Ярким примером такого термина является слово «гаррига, заросли кустарника» (лат. --- \emph{garrice}; кат. --- \emph{gariges})\footnote{Usatges de Barcelona. Us. 68 (72): «STRATE et uie publice, aque currentes et fontes uiui, prata et paschue, silue et garrice et roche in hac patria fundate sunt de potestatibus...». P.~106.}, означающее гарригу, т.е. пустошь, покрытую травами и кустарником. И.С.~Филиппов описывал гарригу в Южной Франции следующим образом: \emph{«Речь идет о неровной каменистой местности, покрытой вереском, лавандой, тимьяном и другими ароматными травами, иногда и непролазным кустарником, хотя можно встретить и одиноко стоящие деревья, даже дубы»}\footnote{\emph{ Филиппов И.С. }Средиземноморская Франция в раннее средневековье… С.~154.}. По данным актового материала, такие пустоши были весьма распространены как в Южной Франции, так и в Каталонии и составляли неотъемлемую часть сельского пейзажа, поэтому именно такое обозначение было принято в старокаталанском языке.

Подводя итог, следует отметить, что ландшафтная и хозяйственная лексика позволяет особенно ясно увидеть практическую природу иноязычных включений в латинском тексте обычая. В этой сфере составитель обращается к местным терминам не ради стилистического разнообразия и не вследствие простой языковой инерции, а потому, что именно они точнее всего обозначали объекты, постоянно входившие в зону правового регулирования: каналы, угодья, пустоши, элементы освоенного пространства. Такие слова фиксируют не только языковые контакты Каталонии с арабским и окситанским миром, но и особенности самого правового мышления, ориентированного на конкретную материальную среду. Тем самым латинский текст «Барселонских обычаев» предстаёт не как отвлечённый свод норм, а как источник, тесно связанный с ландшафтом, хозяйственным укладом и практикой пользования пространством, внутри которых право реально действовало.

\textbf{Выводы к параграфу §2.2}\textbf{.}

В целом анализ материала показывает, что иноязычные элементы в латиноязычных сводах местного обычного права Каталонии XII--XIII~вв. не были случайными вкраплениями в латинском тексте и не сводились к единичным лексическим отклонениям. Напротив, они образовывали функционально значимую часть правового языка, через которую учёная правовая культура соприкасалась с локальной языковой, социальной и хозяйственной средой. Особенно ясно это видно на материале «Барселонских обычаев», где иноязычные элементы выступают как часть продуманной языковой стратегии составителя, работавшего на стыке книжной латинской традиции и живого местного словоупотребления.

Проведённый анализ позволил установить, что включение таких элементов в латинский правовой текст осуществлялось несколькими способами. В одних случаях составитель прямо отмечал, что перед ним слово, принадлежащее народному или крестьянскому употреблению; в других --- местный термин вводился в текст без специальных пояснений, как уже привычное и понятное обозначение. Таким образом, составители записей обычного права фиксировали наличие иной языковой среды и определяли степень проникновения местного слова внутрь латинского текста. Там, где требовалась дистанция по отношению к разговорной речи, использовалось специальное пояснение, маркирующее разговорное употребление слова, а там, где слово уже воспринималось как в достаточной мере значимое, оно включалось в норму непосредственно.

Не менее существенным оказалось и то, что иноязычные элементы в правовом тексте могли получать социальную характеристику. Указания на то, что то или иное слово употребляется крестьянами или «в народе», свидетельствуют о том, что составитель осмыслял языковую вариативность и в категориях социальной иерархии. Язык, на котором общались крестьяне, простолюдины и горожане, не был однородным фоном для правовой культуры, которая становилась пространством взаимодействия носителей различных регистров латинского и старокаталанского языков. Тем самым в латинском нормативном тексте выстраивалось определённое соотношение между книжным языком права и локальным словоупотреблением, одновременно признавалось его существование.

Особенности военной и технической лексики показали, что местные обозначения особенно устойчиво сохранялись там, где речь шла о конкретных формах вооружённого действия, осадной практике, охране замков и связанных с этим реалиях. В этой сфере составители латинских текстов нередко опирались на местные термины как на наиболее точные и практически понятные обозначения. При этом часть такой лексики вводилась без социальной маркировки, что позволяет связывать её с привычным словарём знати и её повседневной культуры, тогда как другая часть напрямую приписывалась крестьянской среде. Следовательно, и внутри одной тематической группы правовой текст отражал различие языковых регистров. Это особенно важно для понимания каталонской правовой культуры, формировавшейся в условиях пограничья, военной активности и тесной связи права с практикой вооружённого конфликта.

Особую группу составляют слова, которые становились непосредственным объектом правовой оценки. Оскорбительная лексика, конфессиональные обозначения и формулы словесного бесчестия важны тем, что здесь юридическое значение имело не только содержание высказывания, но и его конкретная словесная форма. Именно поэтому составитель стремился сохранять подобные слова в том виде, в каком они реально могли прозвучать в повседневной речевой практике. В этих случаях латинский текст фиксировал не отвлечённые категории оскорбления, а социально и конфессионально значимые формы речевого конфликта. Это показывает, что правовая письменность включала в себя не только язык описания нормы, но и язык самого правонарушения, превращая локальные словесные формы в предмет юридической квалификации и санкции.

Ещё яснее зависимость латинского текста от местной правовой практики проявляется в терминах, обозначающих специфические институты каталонского обычного права. Такие слова, как обозначения приданого, сеньориальных прав на имущество бездетного держателя или взысканий, связанных с супружеской неверностью, не могут рассматриваться как простые бытовые заимствования. За ними стоят вполне определённые местные правовые реалии, для которых традиционный латинский юридический язык не всегда предлагал равнозначный и точный эквивалент. Поэтому сохранение местного термина в латинском тексте было не признаком языковой небрежности, а способом точной фиксации особого института. В этом отношении иноязычные элементы свидетельствуют о том, что латинская письменная норма не только описывала местный обычай, но и признавала за ним собственную понятийную и терминологическую самостоятельность.

То же относится и к лексике, связанной с ландшафтом, землепользованием и хозяйственной жизнью. Именно здесь особенно заметно, насколько тесно право было связано с конкретной материальной средой. Каналы, пустоши, угодья, дороги, хозяйственные объекты и иные элементы освоенного пространства требовали точного называния, и латинский текст нередко опирался на местные обозначения как на наиболее адекватные. В эту сферу входили не только романские слова, но и арабские и окситанские заимствования, что отражает широкий круг языковых контактов Каталонии. Однако особенно важно другое: подобные слова включались в текст не ради этнографической точности, а потому, что только они позволяли передать ту совокупность хозяйственных, территориальных и нередко правовых оттенков, которая была существенна для самой нормы.

Сопоставление «Барселонских обычаев» с другими сводами и с актовым материалом показывает, что подобные явления не были исключительной особенностью одного памятника. Формулы, вводившие местное слово как общеупотребительное, а также различные способы включения местных обозначений в латинский документ встречаются и в грамотах XII~в., где служат средством точной идентификации лиц, мест и объектов. Это позволяет рассматривать «Барселонские обычаи» не как исключение, а как наиболее последовательное выражение более широкой практики. Различие между сводом права и деловым документом заключается не в наличии или отсутствии местной лексики, а в степени её обработки и осмысления: в кодифицированном тексте такая лексика получает более высокий уровень рефлексии и становится частью сознательной стратегии юридического письма.

Таким образом, иноязычные элементы в латиноязычных сводах местного обычного права Каталонии XII--XIII~вв. следует рассматривать как важный показатель самой правовой культуры этого региона. Они отражают взаимодействие латинской письменной традиции с местным языковым узусом, социальную стратификацию речевых регистров, включение в правовой текст практического опыта войны, хозяйства и землепользования, а также закрепление специфических институтов местного права в их собственных терминах. Всё это позволяет заключить, что латинский нормативный текст в Каталонии постоянно взаимодействовал с местной языковой традицией, варьируя степень её маркировки, адаптации и признания. В этом и проявилась одна из характерных черт каталонской правовой культуры рассматриваемого периода: способность соединять авторитет письменной латинской нормы с конкретной социальной и хозяйственной средой, в которой это право действовало.

\section{Латинские заимствования и цитаты в каталаноязычных сводах обычного права}

В каталаноязычных сводах обычного права латинские элементы сохранялись как значимые компоненты языка права. Речь идёт как о прямых заимствованиях отдельных терминов, так и о более или менее развернутых латинских цитатах, включённых в тексты на старокаталанском языке. Такие включения особенно важны для понимания того, каким образом формировался язык права в условиях сосуществования книжной латыни и развивающейся письменной каталанской нормы. Они позволяют увидеть, какие понятия ещё продолжали связываться с авторитетом латинской традиции, а какие уже свободно переходили в романскую правовую среду. Поэтому анализ латинских заимствований и цитат в данном разделе будет направлен не только на описание языковых фактов, но и на выявление тех механизмов, посредством которых латинский юридический фонд сохранял своё присутствие в каталаноязычном обычном праве.

\subsection{Латинизмы в старокаталанском переводе «Барселонских обычаев»}

Далее обратимся к тексту старокаталанского перевода «Барселонских обычаев», в данном случае нашей целью будет поиск латинского понятийного аппарата. Однако нас будут интересовать не только случаи использования латинизмов, но и те обстоятельства, в которых переводчик прибегает к их использованию и причины, заставившие его принять именно такое решение для каждого случая. Отдельным исследовательским вопросом в таком контексте становится и то, в какой степени переводчик считал необходимым адаптировать эти латинские понятия к формирующимся правилам старокаталанской графики и орфографии.

По нашим наблюдениям, большинство латинских терминов в переводе обычаев на старокаталанском языке составляют термины, связанные с юриспруденцией в целом и процессуальным правом в частности. Практически без изменений сохранились, например, такие терминологические соответствия: «правосудие» (лат. --- \emph{iusticia}; кат. --- \emph{justicía})\emph{\textsuperscript{ }}\footnote{Usatges de Barcelona. Us. 60 (64), 70 (76), 73 (93-94-95), 77 (103), 79 (104), 80 (105), 87 (110), 95 (118), 98 (121), 103 (124).}, «возмещение» (лат. --- \emph{compositio}; кат. --- \emph{composició})\emph{\textsuperscript{ }}\footnote{Ibid. 82 (101), 96 (119), 102 (123.3).}, «оскорбление» (лат. --- \emph{contumeliam}; кат. --- \emph{honta})\emph{\textsuperscript{ }}\footnote{Ibid. 12 (15). Сомнительное чтение, по мнению Ж. Бастардеса-и-Пареры, в статье №~5: см. комментарии к тексту «Барселонских обычаев»: Usatges de Barcelona / ed. por Bastardas i Parera J. Barcelona, 1991. P.~55. N. 13.}, «неисполнение» (лат. --- \emph{fallimentum}; кат. --- \emph{faliment})\emph{\textsuperscript{ }}\footnote{Ibid. 30 (34), 64 (68).}, «притязание, обвинение» (лат. --- \emph{calumpnia}; кат. --- \emph{calompnia})\textsuperscript{ }\footnote{Ibid. 111 (132). Этот латинский термин сам по себе очень многозначен, это могло, по нашему мнению, вызвать определённые трудности в переводе на старокаталанский отдельных его значений. Сходное явление фиксируется и для португальского права: см. \emph{Варьяш~О.И. }Средневековое право: Опыт пристального чтения некоторых памятников... С.~196.} и др.

В случаях с терминами «правосудие» (кат. --- \emph{justícia}) и «компенсация» (кат. --- \emph{composició}) важно учитывать наличие в старокаталанском языке близких по значению аналогов. Ими являются соответственно «право» (кат. --- \emph{dret})\textsuperscript{ }\footnote{Usatges de Barcelona. Us. 77 (103), 79 (104), 80 (105), 98 (121).}, а также, на наш взгляд, более точное по смыслу «правосудие» (кат. --- \emph{dretura})\textsuperscript{ }\footnote{Ibid. 62 (66.2).} и «возмещение» (кат. --- \emph{emena})\textsuperscript{ }\footnote{Ibid. 4 (4.2--5), 5 (6), 10 (12), 12 (15), 14 (17), 82 (101).}. К этим словам переводчики «Барселонских обычаев» прибегали достаточно часто. Семантические поля терминов «правосудие» и «право», несмотря на внешнюю схожесть и смысловую близость, довольно сильно отличаются. Для переводчика на старокаталанский язык чрезвычайно важны категории результата, но не процесса, поэтому мы можем говорить и о том, что право для него является в каком-то смысле продуктом, результатом правосудия как процесса, и именно поэтому в значении статичной категории он предпочитает использовать термин «право».

Так, в 62(66.2) статье «Барселонских обычаев» латинская формула\emph{ «верность, правосудие, мир и истина государя»} передана как \emph{«верность, справедливость, мир и истина государя»}\footnote{Ibid. Us. 62 (66.2). P.~100--101.}. В данном контексте правосудие представлено не как совершаемое действие, а как одно из свойств государя, посредством которых осуществляется управление королевством. Поэтому латинскому слову «правосудие» (лат. --- \emph{iusticia}) соответствует каталанское слово «справедливость, правота» (кат. --- \emph{dretura}).

Иная ситуация представлена в 73(93-94-95) статье, где латинская формула \emph{«осуществлять правосудие над преступниками»} передана как \emph{«вершить правосудие над преступниками»}\footnote{Ibid. Us. 73 (93-94-95): «...iusticiam facere de malefactoribus...$\rightarrow$...fer justícia dels malsfeytors...». P.~112--113.}. Далее следует перечисление конкретных карательных действий, поэтому слово «правосудие» (кат. --- \emph{justícia}) обозначает здесь осуществление судебно-карательных полномочий власти.

Особенно показателен перевод статьи №~80 (105). Латинская формула «\emph{призвать его к осуществлению правосудия» }передана как \emph{«призвать его предоставить удовлетворение по праву»}\footnote{Ibid. Us. 80 (105): «…ad iusticiam faciendam eum uocauerit…$\rightarrow$...a dret a fer l’apelara...». P.118--119.}. В данном случае речь идёт о предоставлении жалобщику причитающегося ему по праву удовлетворения. Однако ниже латинскому выражению \emph{«предстанет перед правосудием»} соответствует формула «явится к правосудию», то есть предстанет перед судом и подчинится судебному разбирательству\footnote{Ibid. Us. 80 (105): «… ad iusticiam uenerit …$\rightarrow$... a justícia vendrà...». P.~118--119.}. Следовательно, слово «право» (кат. --- \emph{dret}) обозначает прежде всего должный результат разрешения притязания, тогда как слово «правосудие» (кат. --- \emph{justícia}) связано с самой судебной процедурой и осуществлением властных полномочий. Поэтому речь идёт о контекстуальном разграничении семантических полей этих терминов: правосудие понималось как деятельность, а право --- как удовлетворение, которое должно быть предоставлено или получено.

Контекстуальное разграничение близких по значению терминов прослеживается и в употреблении слов «возмещение» (кат. --- \emph{emena}) и «компенсация» (кат. --- \emph{composició}). В статье №~82 (101) латинское обозначение «компенсация» (лат. --- \emph{compositio}), вынесенное в заголовок нормы, передано более общим словом «возмещение» (кат. --- \emph{emena})\footnote{Ibid. Us. 82 (101): «De composicione omnium hominum qui sunt interfecti $\rightarrow$ De la emena dels homens qui son morts». P.~120--121.}, в основной части статьи при обозначении причитающейся родственникам убитого выплаты сохранён термин «компенсация» (кат. --- \emph{composició}). Это позволяет предположить, что первый термин имел более общее значение устранения последствий причинённого ущерба, тогда как слово «компенсация» обозначало конкретную выплату за убийство, присуждаемую на основании закона или местного обычая.

В отдельных случаях переводчик не мог в полной мере понять, что именно имел в виду составитель «Барселонских обычаев». Сравним перевод латинской и старокаталанской версий статьи №~46 (49):

\emph{Таблица }\emph{X}\emph{.}\emph{ }\emph{Передача формулы клятвы в латинской и старокаталанской версиях статьи №~46 (49) «Барселонских обычаев»}

\begin{center}
\small
\setlength{\tabcolsep}{2pt}
\begin{longtable}{|p{\dimexpr0.94\textwidth/2\relax}|p{\dimexpr0.94\textwidth/2\relax}|}
\hline
\textbf{Латинская версия} & \textbf{Старокаталанская версия} \\ \hline
\emph{«<...> et ille qui iurauerit, in omni sacramento debet mittere suo «sciente», excepto in bauzia et in tradicione, et “per Deum et hec sancta”»} & \emph{«<...> }\emph{e}\emph{ }\emph{aquel}\emph{ }\emph{qui}\emph{ }\emph{jurar}\emph{à }\emph{dir}\emph{ }\emph{veritat}\emph{ }\emph{si}\emph{ }\emph{“}\emph{Deus}\emph{ }\emph{li}\emph{ }\emph{ajut}\emph{ }\emph{e}\emph{ }\emph{aquests}\emph{ }\emph{sens}\emph{”»}\footnote{Usatges de Barcelona. Us. 46 (49). P.~86--87.}\emph{.} \\ \hline
\emph{«<…> и тот, кто клянётся, в каждой клятве должен призвать в “свидетели”, за исключением случаев лжи и предательства, “Бога и его святых”»} & \emph{«<…> и тот, кто клянётся, пусть правду говорит: “Если Бог поможет ему со этими святыми”»}\emph{.} \\ \hline
\end{longtable}
\end{center}

По замечанию Х.~Бастардеса-и-Пареры, переводчик допустил ряд ошибок и неточностей при передаче смысла оригинальной статьи, тем самым продемонстрировав непонимание текста данного правового установления\footnote{См. комментарии к тексту «Барселонских обычаев»: Usatges de Barcelona / ed. por Bastardas i Parera J. Barcelona, 1991. P.~87. N. 34--35.}. В старокаталанском переводе опущены указание на роль Бога в качестве «свидетеля» и фрагмент о том, в каких случаях данная форма клятвы не может применяться. Но они не просто теряются, их заменяют совсем другие термины, и, более того, новые смыслы, связанные с идеей установления истины и Божьей помощью приносящему клятву.

Особенно нам в этом фрагменте интересно написание слова \emph{“Déu}\emph{s}\emph{”.} По замечанию каталанского исследователя Дж.~Бругера, это переходная форма встречается в «Хронике Жауме I» (XIII~в.) как рудимент падежной латинской системы, поскольку во всех значениях, совпадающих по значению с косвенными падежами, используется форма \emph{“Déu”}\footnote{\emph{Bruguera~J.} Notes al vocabulari de la Crònica de Jaume I // Actes del IV Col·loqui internacional de Llengua i Literatura Catalanes. Basilea, 22--27 de mars 1976. Barcelona. 1977, P.~87--88.}. Поскольку перевод «Барселонских обычаев» также относится к XIII~в., то можно предположить, что в данном случае мы имеем дело с тем же явлением, при этом переводчик заменяет винительный падеж латинского оригинала на именительный. Вероятно, можно говорить об отношении к религиозной и шире сакральной терминологии, в которой инновации лингвистического рода происходят медленнее в силу уважительного и особо трепетного отношения к ней носителей языка.

\subsection{Латинские вставки в «Обычаях Тортосы»}

Латинские вставки в тексте «Обычаев Тортосы» представляли собой как отдельные термины или словосочетания, так и установления, которые синтаксически оказывались вписаны в предложение на старокаталанском языке.

Количество таких заимствований различалось: в рукописи А таких случаев фиксируется 510, а в более поздних рукописях B и С -- около 200. Такое соотношение можно объяснить, во многом, обратившись к сохранившимся в рукописи A маргиналиям, которые в большинстве случаев указывают на то, какие положения были приняты «арбитрами» при составлении второй редакции памятника, а какие -- отклонены. Однако исчезновение некоторых латинских установлений из текста «Обычаев» не сопровождалось никакими пометами в маргиналиях.

Среди выявленных заимствований в первую очередь стоит выделить самый многочисленный класс -- служебные словосочетания и термины, в который помимо отметок о начале книги или рубрики «Обычаев Тортосы» («Конец первой книги. Начало второй книги» и т.~д.) вошли некоторые вводные и союзные слова\footnote{Costums de Tortosa. P.~77, 131, 183, 264, 298, 340, 372, 408.}. Наиболее частым вводным словом, открывающим то или иное правовое установление, является термин “Item”: в рукописи А он встречается 80 раз, а в поздних рукописях -- 31. Столь частое использование данного маркера нового установления является типичным не только для законодательных текстов (в подражание текстам реципированного римского права), но и для актового материала на территории Каталонии, что подтверждают наблюдения М.~Зиммермана\footnote{Zimmerman M. Écrire et lire… T.~1. P.~286.}.

Снижение частоты употребления \emph{“Item”} в поздних рукописях связано с правкой арбитров, которые во многих случаях исключили это обозначение для облегчения цитирования\footnote{Costums de Tortosa. P.~35.}. Довольно частотным является вводное словосочетание «в силу самого закона»\emph{ }(лат. ---\emph{ }\emph{ipso iure}), использованное легистами в качестве наречия для обозначения правовых последствий того или иного действия или поступка. Сохранение данных вводных словосочетаний в отдельных случаях можно объяснить тем, что и сами арбитры при редакции текста сочли их его типичной частью, если это не противоречило их основной цели.

Латинский язык использовался составителями для того, чтобы зафиксировать формулы, применявшиеся при составлении документов\footnote{Ibid. 6.4.19; 6.4.25--26; 9.9.7--8. P.~311, 313--314, 432--433.}. Например, в тексте обычая 6.4.26 предлагается следующая стандартная форма для подписи свидетеля в завещании: \emph{«Я такой-то по просьбе такого-то подписываюсь свидетелем и подпись свою ставлю»}\footnote{Ibid. 6.4.26: «“ego talis rogatus a tali me subscribo pro teste et signum meum appono”». P.~314.}. Акт, составленный иначе, согласно тексту данного обычая, мог быть признан недействительным. В большинстве случаев (11 из 12) эти формулы сохранились и в поздних рукописях; арбитры решили их не исключать.

Текст статьи 1.4.9 был исключён в ходе редакции «Обычаев Тортосы». Он предусматривал специальную форму заместительной подписи для должностных лиц и судей на тот случай, если они не умели писать. Тот факт, что Пётр Тамарит и Пётр Эгидий как нотарии, так или иначе имевшие дело с практиками составления актов, считали важным оговорить такую ситуацию, свидетельствует о том, что она была весьма вероятной, однако по каким-то причинам во вторую редакцию «Обычаев Тортосы» этот фрагмент не вошёл. По нашему мнению, таких причин могло быть несколько, несмотря на то, что маргиналия в первой редакции «Обычаев Тортосы» о них умалчивает.

Исключения фрагментов текстов можно наблюдать в грамотах на «фаршированной латыни», происходящих из Южной Франции и представляющих собой клятвы верности сеньору. В них отдельные глаголы и даже целые предложения записаны на окситанском и старопровансальском языках в то время, как весь основной текст грамоты, включая формуляр, на латыни\footnote{\emph{Belmon J., Vielliard F.} Latin farci et occitan dans les actes du XIe siècle // Bibliothèque de l’École des chartes. Paris, 1997. Vol.~155, no. 1. P.~149.}.

Следующая категория заимствований, выделенная нами -- это отдельные термины, также неоднократно появляющиеся в названиях рубрик «Обычаев Тортосы». Отдельные термины приводятся в своей латинской неизменной форме, например, \emph{“}\emph{Rubrica}\emph{ }\emph{de}\emph{ }\emph{usu}\emph{ }\emph{fructu}\emph{…”}\footnote{Costums de Tortosa. 3.9.1. P.~154.}, \emph{“}\emph{Rubrica}\emph{ }\emph{de}\emph{ }\emph{negociis}\emph{ }\emph{gestis}”\footnote{Ibid. 2.11.1. P.~107.} и т. д. Латинская терминология в заголовках рубрик встречается преимущественно в тексте 1-й редакции «Обычаев Тортосы». В таких случаях после латинского обозначения рубрики обычно следовала краткая дефиниция предмета рубрики на старокаталанском языке: \emph{«Рубрика об узуфрукте}\footnote{\textbf{Узуфрукт} (лат. usufructus)--- вещное право лично пользоваться чужой непотребляемой вещью и извлекать из неё естественные плоды, не изменяя сущности самой вещи. См. Usufructus // \emph{Бартошек М. }Римское право: (Понятия, термины, определения) / пер. с чешк. М.: Юридическая литература, 1989. С.~322.}\emph{, то есть о тех [лицах--А.К.], которые имеют право получать доходы с какой-либо вещи и не обладают правом собственности на неё»}\footnote{Costums de Tortosa. 3.9.1: «Rúbrica de usu fructu, so és d'aquels que an dret en reebre fru<it> d' aquela cosa e no an dret en la proprietat». P.~154.}. В данном случае мы видим, как при попытке объяснить определённый латинский термин на старокаталанском языке изменялось и его содержание: в данном случае фокус переместился с абстрактного права пользоваться чем-либо на людей, которые таким правом могли обладать.

Подобное построение предложения представляет собой не неосознанное переключение языков в тексте легиста-интеллектуала, а его осознанный компромисс с адресатом данного законодательного текста, его предполагаемым пользователем, для которого отдельные термины римского права требовали специальных пояснений, причём на родном для него языке. Подобные синтаксические и смысловые структуры были не единственным вариантом придания понятийному аппарату римского права более «конвенциональных» форм. Так, К.~Дуарте-и-Монтсеррат для рукописи А фиксирует рост числа латинских заимствований в собственно этимологическом смысле по отношению к более ранним текстам, многие из которых были зафиксированы им впервые в документах на старокаталанском языке\footnote{\emph{Duarte i Montserrat C.} La llengua del manuscrit de 1272 del Llibre de les Costums de Tortosa // Revista de llengua i dret. 1991. №~16. P.~LXX.} Это говорит об активном освоении и в определённой мере «присвоении» латинского юридического вокабулярия.

С другой стороны, в тексте «Обычаев Тортосы» существовали и различного рода дефиниции, записанные полностью на латыни: \emph{«Прокуратор -- это тот, кто безвозмездно управляет каким-либо делом по поручению господина»}\footnote{Costums de Tortosa. 2.9.1: «Procurator est qui gratuito aliena negotia mandato domini administrat». P.~99.}. На наличие сходного определения термина «прокуратор» в «Обычаях Миравета» указал испанский исследователь истории права Ф.~де~Арвизу-и-Галаррага\footnote{\emph{de Arvizu y Galarraga F}. Fianzas procesales en la documentación altomedieval // Anuario de Historia del Derecho Español. 2022. Vol.~92.P.~54--55.}. Уже в тексте другой статьи встречаем следующее определение страха: \emph{«Страх -- это смятение разума по причине настоящей или будущей опасности»}\footnote{Costums de Tortosa. 2.12.1: «Metus est instantis vel futuri periculi causa mentis trepidacio». P.~110.}. По наблюдениям А.~Иглесиаса-Ферейроса и Дж.~Массипа-и-Фонойосы большинство этих дефиниций происходили из компилятивного сочинения известного в XIII~в. болонского легиста Азо Порциуса “Summa codicis”, которое представляло собой сокращённое обозрение Кодекса Юстиниана и Институций, или из сборника устоявшихся римских правовых принципов, получившего название “Generalia” (официально) или “Brocada” (в профессиональном юридическом лексиконе)\footnote{Ibid. 2.12.1. P.~110; \emph{Kuttner S.} Gratian and the Schools of Law, 1140--1234. London; New York: Routledge, 2018. P.~251.}. О знакомстве Петра Тамарита и Петра Эгидия с компиляциями Азо Порциуса свидетельствует эксплицитное упоминание этого болонского юриста при одной из дефиниций в рукописи A: \emph{«}\emph{…или, }\textbf{\emph{согласно Азо}}\emph{, арбитраж определяется следующим образом: арбитраж --- это дело, добровольно переданное для вынесения }\emph{решения при участии определённых лиц, а именно арбитра, истца и ответчика}\emph{»}\footnote{Costums de Tortosa. 2.16.1: «…vel secundum Azonem diffinitur arbitrium hoc modo: arbitrium est causa in dicendo iure voluntarie posita cum certarum personarum interposicione, scilicet arbitri, actoris et rey». P.~120.}. Обратимся к принципам подбора подобных дефиниций, которыми руководствовались тортосские нотарии при составлении «Обычаев Тортосы».

Во-первых, среди дефиниций практически все имеют стандартную форму: термин, форма глагола “esse” (лат. «быть») и собственно определение. Во-вторых, такие определения довольно лаконичны, в отличие от рассмотренных выше дефиниций и пояснений на старокаталанском языке. Ещё одним отличием таких определений от рассмотренных выше случаев использования латинского термина и объяснения на старокаталанском является то, что большинство латинских дефиниций из рукописи А были исключены арбитрами как «излишние» или носящие дидактический характер, на что указывает издатель источника\footnote{Ibid. P.~110.}. Однако в тех случаях, когда дефиниция такого рода составляла само правовое установление, арбитры допустили её использование: \emph{«Обвинять или защищать [в суде--А.К.] -- это ходатайство своё или друга своего по праву перед тем, кто осуществляет правосудие, излагать, или возражать против ходатайства других [лиц--А.К.]»)}\footnote{Ibid. 2.7.1: «Postulare sive advocare est, desiderium suum vel amici sui in iure apud eum, qui iurisdiccioni preest, exponere, vel alterius desideria contradicere». P.~96.}. Подобное использование латиноязычного объяснения в качестве нормы можно видеть и в обычае 4.17.3: \emph{«}\emph{Различие между займом, депозитом и ссудой состоит в следующем. }\textbf{\emph{Займом}}\emph{ является то, что определяется весом, числом или мерой и потому становится собственностью получателя. }\textbf{\emph{Депозитом}}\emph{ является то, что не становится собственностью получателя и чем он не должен }\emph{пользоваться. }\textbf{\emph{Ссудой}}\emph{ является то, что не становится собственностью получателя, но чем он может пользоваться}\emph{»}\footnote{Ibid. 4.17.3: «Differentia est inter mutuum, depositum et comodatum: Mutuum est quod consistit in pondere, numero et mensura, et sic accipientis. Depositum est quod non sit accipientis nec debet uti eo, si consignatum est.Comodatum est quod non sit accipientis, set potest uti eo». P.~226.}.

Собственно, правовые установления, записанные на латинском языке в «Обычаях Тортосы» также присутствуют: назовём рубрики 22--23, 29 девятой книги данного компендиума. Установления рубрик 22--23 в большинстве своём лаконичны, просты для понимания и имеют много общего со случаями прямого цитирования в латинских дефинициях. Издатель «Обычаев Тортосы» в примечаниях к данным рубрикам отмечал, что именно в их тексте прослеживаются параллели с «Фуэро Валенсии»\footnote{Ibid. P.~479.}. В данном случае говорить о текстуальных связах «Обычаев Тортосы» и «Фуэро Валенсии» трудно, поскольку текст установлений, несмотря на незначительные различия, вероятно, является вторичным заимствованием 16-го и 17-го титулов 50-й книги Дигест Юстиниана из одного третьего источника.

Далее, по тексту 29 рубрики 9-й книги «Обычаев Тортосы» мы можем судить, что тортосские легисты пользовались именно латинской версией «Барселонских обычаев». В частности, в этой рубрике в рукописи А и поздних рукописях приведён полный латинский текст обычаев 4--11, 13--15, 17--22, 58, 129--131\footnote{Ibid. 9.29.1--14. P.~520--524.}. Текст рубрики дополняется коротким добавлением на старокаталанском языке следующего содержания: \emph{«В этих установлениях содержится обычай Тортосы, поэтому сеньория по этим обычаем никакой компенсации требовать не может, за исключением пятой части штрафа, назначенного по наказанию, поэтому на территории сеньории Тортосы никакое преступление публичное или частное, не может судиться без жалобы или обвинения, если только пятая часть }\emph{любой жалобы не может быть удовлетворена посредством денег. Все же суммы пусть исчисляются в монете, имеющей хождение в Тортосе, и ни в какой иной»}\footnote{Ibid. 9.29.15: «En aquests usàgies se contén la custum de Tortosa, enaxí que la seynoria, nulla pena per aquests usàgies demanar no pot, sinó la quinta part d’aitant com serà condemnat lo demanat; per so com la seynoria en Tortosa, ne en sos térmens, de nul crim públic o privat, no a ne pot ne deu fer contra alcun demanda ne accusatió, sinó tant solament lo quint de les condemnacions en què alcun serà condemnat de so on clams seran fets peccuniàriament. E els sous són entesos de la moneda currible en Tortosa, e no d’altra moneda». P.~524.}. Из текста этой нормы следует, что «Барселонские обычаи», которые считались основным источником права в Тортосе с 1149~г., легисты вынуждены были адаптировать к локальным реалиям и сопровождать пояснениями на старокаталанском языке.

\subsection{Латинская цитата как способ обращения к авторитетному тексту}

Помимо указанных рубрик, стоит упомянуть и о том, что в тексте «Обычаев Тортосы» существовали отдельные установления на латинском языке, помещённые внутрь рубрики, основной текст которой был старокаталанским (например, статья 2.18.9)\footnote{Ibid. 2.18.9: «Ius iurandi contempta religio satis Deum habet ultorem nam sufficit pena periurii quam a Deo expectat». P.~129.}. Большинство подобных заимствований были либо признаны избыточными и изъяты арбитрами\footnote{Ibid. 2.18.9: «Terciam decimam que incipit: “Si alcun christià et cetera”, repprobant quia deperit veritas et iusticia impeditur». P.~129.}, либо были переведены на старокаталанский язык, как в статье 6.6.1\footnote{Ibid. 6.6.1: «Idem est in extraneo herede quod de filio» $\rightarrow$ «E aquel dret metex és en l’ereu estrayn, que és e·l fiyl». P.~320.}.

Гораздо чаще встречается ситуация, когда тот или иной обычай на старокаталанском языке лишь дополняется кратким латинским изречением дидактического или морализаторского характера. Чаще всего в таких случаях латинская фраза не имела характера правового установления и лишь завершала обычай: \emph{«Поскольку обязанность мужчины -- защищать других, поскольку он считается выше }\emph{женского пола»}\footnote{Ibid.:«Quia virile officium est suscipere alienam defensionem e ultra sexum muliebre esse constant». P.~102.}. Подобные сентенции, также, вероятно происходившие из юридических трактатов, циркулировавших в среде легистов, из поздних рукописей в большинстве случаев исключались. Однако важно отметить, что в этом случае арбитры отказывались не только от латинских сентенций, но и от самих обычаев на старокаталанском языке, в которые эти латинские сентенции были включены. Происходило это тогда, когда, по мнению арбитров, обычаи не соответствовали церковным установлениям, о чём свидетельствуют пометки в маргиналиях рукописи А, причём фиксировалось, что «греховные» установления должны быть заменены принципами общего римского права\footnote{Ibid. 1.5.2: «…repprobant quia continet peccatum et loco eius servetur Ius commune». P.~48.}. Всё это свидетельствует о том, что даже в поздних рукописях влияние учёной традиции римского права продолжало сохранять свой авторитет при исключении отдельных положений, которые представлялись сомнительными церковным иерархам, которые были знакомы с римской правовой традицией прежде всего по каноническому праву.

Выразить отдельные религиозные термины и установления легистам было гораздо легче на латинском языке, ведь именно с латинским вариантом они чаще всего встречались в церковных текстах или при составлении тех или иных документов, с которыми грамотный человек встречался в течение своей жизни. Пример такой вставки, органично вписавшейся в предложение на старокаталанском, можно встретить в статье 2.8.6: \emph{«Поскольку, как говорит Павел, браки заключаются во Христе» }(1 Кор. 7:39)\footnote{Ibid. 2.8.6: «Per so com son Paul diga nubat in Xristo». P.~98.}. Скорее всего, в данном случае мы имеем дело с не совсем точной ссылкой на первое послание апостола Павла к Коринфянам.

Церковные иерархи усомнились в необходимости данного пассажа в тексте законодательного характера, как считал Ж.~Массип-и-Фонойоса\footnote{Ibid. P.~98.}. При этом тенденция к устранению подобных фрагментов религиозного содержания напрямую не коррелирует с их языковой принадлежностью, поскольку из обычая 6.1.13 был исключён по тем же причинам абзац на старокаталанском языке об искуплении грехов Иисусом Христом\footnote{Ibid. 6.1.13: «Per so con dura cosa e lèygia seria que aquels qui reemutz són per la gràcia de Jesuxrist e per la sua preciosa sanc, fossen detengutz e preses en preson ne en vincle de jueus ne en neguna servitut». P.~98.}. Говоря о заимствованиях религиозного характера, также необходимо отметить, что даже в старокаталанском тексте терминология, связанная с христианской верой и её правовыми аспектами, предстаёт в своеобразной форме. В «Обычаях Тортосы» фиксируется типичная для XIII~в. двойственность в морфологии и графике слова \emph{“Deus”} (лат. «Бог»); она выражается в том, что в отдельных случаях у слова, употребляющегося внутри старокаталанского текста, фиксируются отголоски латинской падежной системы. Это явление неоднократно фиксировалось исследователями-филологами и встречается в других источниках этой эпохи, например, в переводе на старокаталанский язык «Барселонских обычаев», «Хронике Жауме I»\footnote{\emph{Duarte i Montserrat C. }La llengua del manuscrit de 1272 del Llibre de les Costums de Tortosa // Revista de llengua i dret. 1991. №~16. P.~LV; №~3, P.~87--88.}.

\textbf{Выводы к параграфу §2.3}\textbf{.}

Итак, в текстах обеих редакций «Обычаев Тортосы» нами было выделено 7 категорий латиноязычных элементов: служебные и вводные слова, нормативные варианты текстов подписей сторон для договоров, тексты печатей для публичной власти, латинская терминология в названиях правовых установлений, дефиниции на латинском языке, сентенции морализаторского содержания и религиозная терминология.

Анализируя обстоятельства использования латиноязычных элементов в тексте «Обычаев Тортосы», мы пришли к следующим выводам.

Во-первых, в первоначальной редакции «Обычаев Тортосы» гораздо больше латинских заимствований. Не менее важно и то, что практически половина латинских заимствований в результате второй редакции исчезла из текста памятника, показывает, что многие из них носили вспомогательный характер и были призваны дополнить изначальный старокаталанский текст.

Во-вторых, сохранение значительного числа латинских заимствований свидетельствует о том, что легисты, носители старокаталанского языка, воспринимали латинские термины, словосочетания и фразы как органичный элемент языка. Это и есть основное проявление мультилингвизма в различных ситуациях. Использование и латинского, и старокаталанского языков было вполне нормальным и естественным даже в рамках одного предложения. Об этом свидетельствуют результаты и количественного анализа подобных заимствований.

В-третьих, легисты-интеллектуалы, наслышанные о римском праве, знакомые с ним по сборникам и обзорам Corpus Iuris Civilis, стремились по возможности сделать кодекс, который они составляли, наиболее приближенным к римскому праву, которое казалось им высшим универсальным достижением правовой науки, при этом сохранив ряд традиционно легитимирующих элементов на латинском языке.

В-четвёртых, церковные иерархи, редактировавшие «Обычаи Тортосы» спустя 5 лет, руководствовались другой целью, а именно сделать и без того объёмный текст обычаев более понятным, они были в меньшей степени знакомы с имевшими тогда хождение сборниками норм римского права, что довольно чётко фиксируется в маргиналиях к рукописи A. Отчасти этим можно объяснить сокращение количества латинских цитирований во второй редакции кодекса.

Однако с традициями канонического права они были знакомы довольно хорошо, потому отдельные нормы римского права не были ими отвергнуты.

\section{Традиция перевода правовых текстов в Каталонии XII--XIII~вв.}

Сопоставление двух версий одного и того же правового памятника позволяет проследить, каким образом при переходе от латинского текста к тексту на старокаталанском языке менялись способы выражения норм, структура и состав понятийного аппарата. Речь, следовательно, идёт не только о переводе в узком филологическом смысле, но и о переоформлении правового содержания в иной коммуникативной среде.

Для истории права это особенно важно, поскольку переводной текст фиксирует изменение представлений о том, как именно должна была быть выражена и воспринята правовая норма. В одних случаях в переводе сохраняли латинскую терминологию почти без изменений, в других --- заменяли её более широкими или, напротив, более конкретными соответствиями, а в третьих --- перестраивали саму синтаксическую организацию статьи. Тем самым переводы «Барселонских обычаев» на старокаталанский язык могут считаться самостоятельными источниками по истории правовой культуры Каталонии XII--XIII~вв.

Особый интерес для нас представляет вопрос, какие элементы латинской юридической традиции сохранялись в переводе, какие подвергались упрощению или переосмыслению, а в каких случаях переводчик позволял себе более свободную интерпретацию нормы. Анализ этих особенностей даёт возможность увидеть, каким образом текст правового памятника приспосабливался к новой аудитории и к меняющимся условиям правового общения. В дальнейшем внимание будет сосредоточено на трёх основных аспектах: семантической переработке терминов, сохранении устойчивых терминологических соответствий и синтаксико-грамматической трансформации латинского оригинала. Такой подход позволяет рассматривать переводы «Барселонских обычаев» как важное свидетельство не только языковой, но и правовой эволюции Каталонии XII--XIII~вв.

\subsection{Выбор терминологических соответствий при переводе с латинского на старокаталанский язык}

Сопоставление латинского текста «Барселонских обычаев» со старокаталанскими переводами показывает, что одна из важнейших особенностей перевода связана с переработкой профессионально-правовой терминологии. Наиболее показательные случаи касаются таких замен, при которых старокаталанский эквивалент оказывался либо семантически шире латинского термина, либо, напротив, уже и конкретнее его. Речь идёт, таким образом, не только о переводе в собственном смысле слова, но и о смысловой адаптации правового текста к иной языковой среде. Именно на этом уровне особенно хорошо заметно, что переводчик стремился не столько воспроизвести латинскую терминологическую систему, сколько сделать норму более прозрачной для круга её читателей.

В результате часть латинских различий в переводе сглаживалась, тогда как другие, напротив, приобретали более определённое и прикладное выражение. Для описания этих явлений далее будут использоваться категории смыслового расширения (генерализация) и смыслового сужения переводного соответствия (конкретизация)\footnote{О генерализации и конкретизации как видах лексико-семантических переводческих трансформаций см.: \emph{Комиссаров В.Н.} Современное переводоведение: учеб. пособие. М.: Р.~Валент, 2011. С.~161--162; \emph{Бархударов Л.С. }Язык и перевод: Вопросы общей и частной теории перевода. М.: ЛЕНАНД, 2016. С.~209--213.}.

Зачастую мы видим замены ряда латинских глаголов старокаталанским глаголом «делать» (кат\emph{. --- }\emph{fer}), причём им заменяются самые различные по значению глаголы «делать» (лат. --- \emph{facere}), «приводить в порядок» (лат. --- \emph{condirigere})\footnote{Usatges de Barcelona. Us. 68 (72--73). P.~106--107.}, «исправлять» (лат. --- \emph{redirigere})\footnote{Ibid. Us. 100(123.1). P.~136--137.} и «восстанавливать, возвращать» (лат. --- \emph{restituere})\footnote{Ibid. Us. 62(66.2), 79(104), 106(127). P.~100--101, 118--119, 140--141.}. Что особенно интересно, один из этих глаголов --- «приводить в порядок» (лат. --- \emph{condirigere}), как писал в своей статье Ж.~Бастардес-и-Парера, попал в язык «Барселонских обычаев» из местной разновидности народной латыни\footnote{\emph{Bastardas i Parera J.} El català pre-literari // Actes del Quart Col·loqui Internacional de Llengua i Literatura Catalanes (Basilea, 22--27 de març de 1976) / a cura de G. Colón. Barcelona: Publicacions de l’Abadia de Montserrat, 1977. P.~59.}, использование в нём приставки со значением действия должно было, по мнению исследователя, скомпенсировать продолжавшуюся утрату падежных окончаний. Конечно, в отношении таких частотных глаголов, как «давать» (кат. --- \emph{donar}), «делать» (кат. --- \emph{fer}) и «находиться» (кат. --- \emph{estar}), это постоянное явление\footnote{Usatges de Barcelona. Us. 59(62), 71(91), 80(105), 84(107).} ввиду их частого использования в речи. В таких случаях переводчик решал отойти от сложной юридической лексики в сторону наиболее базовых значений, чтобы сохранить «смысловое ядро» той или иной правовой нормы.

Рассмотрим латинский текст 100 (123.1)-й статьи обычаев и его старокаталанский перевод: \emph{«Христианам запрещается продавать оружие сарацинам без согласия государя; если же они это сделают, }\emph{они обязаны возвратить проданное оружие, хотя бы это и было для них обременительно; в противном случае они должны }\emph{уплатить}\emph{ (лат.}\emph{~}\emph{persoluant}\emph{ --- кат. }\emph{donen}\emph{) власти сто унций золота»}\footnote{Ibid. Us. 100(123.1): «CHRISTIANI non uendant arma sarracenis, nisi ex consensu principis; quod si fecerint arma que uendiderint recuperent, quamuis eis graue sit, et nisi hoc fecerint, centum uncias auri potestati persoluant  $\rightarrow$ [C]restians no venen armes a ssarraïns, sinó ab lo consentiment del príncep; que si ho fan, les armes que an venudes, cobren, ja sie ço que greu sia; e si açò no fan, .С. onçes d’or a la postat». P.~136--137.}\emph{.}

\emph{Таблица X. Передача глагола «уплатить, исполнить обязательство» (лат. persolvere) в старокаталанских версиях «Барселонских обычаев»}

\begin{center}
\small
\setlength{\tabcolsep}{2pt}
\begin{longtable}{|p{\dimexpr0.94\textwidth/5\relax}|p{\dimexpr0.94\textwidth/5\relax}|p{\dimexpr0.94\textwidth/5\relax}|p{\dimexpr0.94\textwidth/5\relax}|p{\dimexpr0.94\textwidth/5\relax}|}
\hline
\textbf{Статья и контекст} & \textbf{Латинский текст по изданию Ж. Бастардеса-и-Пареры} & \textbf{Старокаталанская версия по изданию Ж. Бастардеса-и-Пареры} & \textbf{Старокаталанская версия по изданию Ж. Гудиоля-и-Куниля} & \textbf{Русский перевод} \\ \hline
100 (123.1) & \emph{potestati persoluant} & \emph{donen a la postat} & \emph{donen a la postat} & Пусть дадут власти (кат. \emph{donen}). \\ \hline
113 (134) & \emph{de suo debitum persoluerit} & \emph{de son deute pagarà} & \emph{paga lo deute del seu} & Уплатит долг из собственного имущества \\ \hline
113 (134) & \emph{in duplo ei persoluere cogatur} & \emph{forçat sia que pac en doble}\footnote{Usatges de Barcelona. Us. 113 (134). P.~146--147.} & \emph{en doble li ho reta}\footnote{\emph{Gudiol i Cunill J. }Traducció dels Usatges, les més antigues constitucions de Catalunya y les Costumes de Pere Albert // Anuari de l’Institut d’Estudis Catalans. 1907. Vol.~1. P.~304.} & «Должен возместить ему в двойном размере (кат1.--- \emph{п}\emph{усть уплатит}) \\ \hline
\end{longtable}
\end{center}

Сопоставление показало, что значение латинского глагола \emph{«уплатить, полностью исполнить денежное обязательство» }(лат. --- \emph{persolvere}) определялось сочетанием с обозначением предмета платежа. В~статье 100(123.1) речь шла об уплате установленной нормой денежной суммы --- ста унций золота в пользу власти. В статье 113(134) тот же латинский глагол сначала обозначал погашение долга поручителем из собственного имущества, а затем --- возмещение в двойном размере ущерба, который поручитель должен был понести вследствие неисполнения обязательства основным должником.

Старокаталанский текст не сохранил для этих контекстов единого глагольного соответствия. В 100(123.1) статье переводчик использовал глагол «дать»\emph{ }(кат. --- \emph{donar}), который в сочетании с указанием суммы и её получателя приобрёл значение «уплатить». В 113(134) статье он употребил глагол «платить» (кат. --- \emph{pagar}): сначала применительно к долгу, затем --- к возмещению ущерба\footnote{Pagar // Vocabulari de la llengua catalana medieval de Lluís Faraudo de Saint-Germain. URL~: \url{http://www.iec.cat/faraudo/results.asp?Word=Pagar&Id=3084898&search=pagar&optCriteria=0&optSearchType=1&ActualPage=1} (дата обращения: 14.04.2023).}. Таким образом, значение платежа в обоих языковых вариантах определялось всей конструкцией: глаголом, обозначением суммы или обязательства, получателем и указанием на размер возмещения. В рукописи, опубликованной Ж.~Гудиолем-и-Кунилем, использован другой старокаталанский глагол --- «возместить» (кат. --- \emph{retre}).

Различие между старокаталанскими глаголами в данном случае нельзя уверенно объяснить стремлением переводчика заменить сложную правовую лексику общеупотребительной. Напротив, перевод демонстрирует контекстуальный выбор между двумя глаголами, поэтому речь идёт об отсутствии однозначного соответствия между латинской и старокаталанской лексикой и о распределении значений одного латинского глагола между несколькими переводными эквивалентами.

Встречаются и другие случаи использования генерализации при переводе «Барселонских обычаев» применительно не только к глаголам. Так, в 92(115)-й статье «Барселонских обычаев», где речь идёт об опеке над детьми, оставшимися без попечения родителей, мы встречаем фразу: \emph{«…Если же речь идёт о девушке }(лат.~\emph{puella}---кат. \emph{fembra}),\emph{ то её следует выдать замуж по одобрению и совету достойных }\emph{ }\emph{людей…}»\footnote{Usatges de Barcelona / ed. J. Bastardas i Parera. Barcelona: Fundació Noguera, 1991. Us. 92 (115): «…Sin autem et puella est, det ei maritum laude et consilio proborum hominum ...---… Però si és fembra dó a ela marit per laor e per conseyl de pròmens hominum...». P.~130--131.}.

\emph{Таблица }\emph{X}\emph{. Соответствия латинского слова puella в старокаталанских версиях «Барселонских обычаев»}

\begin{center}
\small
\setlength{\tabcolsep}{2pt}
\begin{longtable}{|p{\dimexpr0.94\textwidth/5\relax}|p{\dimexpr0.94\textwidth/5\relax}|p{\dimexpr0.94\textwidth/5\relax}|p{\dimexpr0.94\textwidth/5\relax}|p{\dimexpr0.94\textwidth/5\relax}|}
\hline
\textbf{Статья} & \textbf{Латинский текст} & \textbf{Старокаталанская версия по Бастардесу} & \textbf{Старокаталанская версия по Гудиолю} & \textbf{Русский перевод} \\ \hline
\textbf{85 (108)},\par \textbf{Usatge 104} & \emph{Si quis violenter virginem corruperit, aut ducat eam in uxorem… }\par \emph{Si non virginem quis violenter adulteraverit et impregnaverit, similiter} & \emph{Si negun volentament corromprà }\textbf{\emph{verge}}\emph{, o la prena }\textbf{\emph{per muller}}\emph{… }\par \emph{Si algú violentament adulterarà la }\textbf{\emph{fembra qui no és verge}}\emph{ e la emprenyarà, axí mateix} & \emph{Si negu força }\textbf{\emph{puncela}}\emph{, ho la prena }\textbf{\emph{per muyle}}\emph{…} Далее в той же норме употребляется \emph{verge}: \emph{no verges / no uergen} & Если кто силой лишит девственности девушку, пусть либо возьмёт её в жёны... Если кто силой вступит в связь с женщиной, не являющейся девственницей, и сделает её беременной, пусть поступит так же» \\ \hline
\textbf{92 (115)} \textbf{Usatge 110} & \emph{Sin autem et }\textbf{\emph{puella}}\emph{ est…} & \emph{Però si és }\textbf{\emph{fembra}}\emph{…} & \emph{E si es }\textbf{\emph{macipa}}\emph{…} & Если же подопечная --- девушка… \\ \hline
\end{longtable}
\end{center}

Мы видим, что латинское существительное «девочка, девушка» (лат. --- \emph{puella}) переводчик передаёт при помощи более общего термина «женщина, самка» (кат. --- \emph{fembra}), который обозначает любую женщину вне зависимости от её возраста или семейного положения. Это можно проследить и по тексту «Барселонских обычаев»: так, например, в 19(22), 85(108) статьях.

Латинские существительные «девушка, дева» (лат. --- \emph{virgo}) и «девочка, девушка» (лат. --- \emph{puella}) не получают в старокаталанских версиях единообразных соответствий. В версии, опубликованной Ж.~Бастардесом-и-Парерой, латинское обозначение девы (лат. --- \emph{virgo}) передано словом «дева» (кат. --- \emph{verge}). Более общее обозначение женщины (кат. --- \emph{fembra}) появляется во второй части статьи в сочетании «женщина, которая не является девой» (кат. --- \emph{fembra qui no és verge}). В данном случае слово «женщина» (кат. --- \emph{fembra}) указывает на принадлежность лица к женскому полу, тогда как слово «дева» (кат. --- \emph{verge}) характеризует состояние девственности.

Ещё отчётливее различие между рукописными версиями проявляется в статье об опеке. Латинское обозначение девочки или девушки (лат. --- \emph{puella}), семантика которого включает значения «девочка», «девушка» и «молодая женщина», в версии Ж.~Бастардеса-и-Пареры передано словом «женщина» (кат. --- \emph{fembra}). Последнее обозначает лицо женского пола, но само по себе не выражает возраст подопечной. В данном случае перевод сопровождается генерализацией: возрастной компонент латинского слова в старокаталанском соответствии непосредственно не воспроизводится.

В версии, опубликованной Ж.~Гудиолем-и-Кунилем, употреблено слово «девушка» (кат. --- \emph{macipa}). Эта лексема обладала несколькими значениями: она могла обозначать служанку или ученицу, а также молодую девушку. Контекст нормы --- несовершеннолетие подопечной, её нахождение под опекой и обязанность опекуна устроить её брак --- актуализирует значение молодой девушки. По сравнению со словом «женщина» (кат. --- \emph{fembra}) обозначение «девушка» (кат. --- \emph{macipa}) полнее сохраняет возрастной компонент латинского слова (лат. --- \emph{puella}). Одновременно оно могло передавать зависимое положение молодой женщины, находившейся под чужим попечением.

Таким образом, в старокаталанских версиях использовалось несколько обозначений девушки или молодой женщины: «дева» (кат. --- \emph{verge}), «девица» (кат. --- \emph{puncela}), «девушка» (кат. --- \emph{macipa}), а также более общее слово «женщина» (кат. --- \emph{fembra}). Ни одно из них не совпадало во всём объёме значений с латинскими обозначениями девы и девушки (лат. --- \emph{virgo}, \emph{puella}). Различия между рукописями свидетельствуют о вариативности переводческих решений и о несовпадении семантических границ латинских и старокаталанских лексем.

Так, например, в латинском тексте 119(77)-й статьи «Барселонских обычаев» мы встречаем термин \emph{“}\emph{g}\emph{/}\emph{ienitores}\emph{”} (лат. --- предки), в то время как в переводе имеем дело с термином \emph{“}\emph{pares}\emph{”} (кат. --- родители): \emph{«Вместе с тем названные выше предки (лат. }\emph{ienitores}\emph{---кат. }\emph{pares}\emph{) вправе лишить наследства своих детей, а также внуков и внучек…»}\footnote{Ibid. Us. 119 (77): «EXHEREDARE autem possunt \textbf{predicti ienitores} filios suos, uel nepotes siue neptes...---[D]eseretar poden los \textbf{pares} los seus fils o eis nets o les netes...». P.~154--155.}\emph{.} В данном случае конкретизация, по нашему мнению, не полностью оправдана, поскольку за ней в обеих версиях текста следует перечисление близких родственников: отец, мать, дед и бабка. Латинский термин в данном случае представляется более точным. Однако, вероятно, его полного эквивалента в старокаталанском языке не существовало, а слово «родители, предки» (кат. --- \emph{pares}) могло относиться к любым родственникам по прямой восходящей линии. Вместе с тем в значении «родители» слово \emph{pares} употребляется в ряде обычаев как эквивалент латинского слова «родители» (лат. --- \emph{patres})\footnote{Ibid. Us. 92 (115), 117 (138). P.~128--129, 138--139.}.

Сходную ситуацию, в которой необходимость применения конкретизации была вызвана контекстом и жанром «Барселонских обычаев», мы видим и при анализе текста 64(68)-й статьи: \emph{«Если же }\emph{государь по какой-либо причине (лат. casu---кат. aventuras) будет осаждён…»}\footnote{Ibid. Us. 64 (68): «PRINCEPS NAMQVE si quolibet casu obsessus fuerit…$\rightarrow$[D]e príncep, però, si aquest, qualques aventuras...». P.~102--103.}\emph{.} В данном случае мы видим конкретизацию, поскольку в старокаталанском дериват латинского \emph{“}\emph{casus}\emph{” --- “}\emph{cas}\emph{”} \emph{---} обозначал не только любое случайное событие, но и юридический факт правонарушения\footnote{Cas// Vocabulari de la llengua catalana medieval de Lluís Faraudo de Saint-Germain. URL: \url{http://www.iec.cat/faraudo/results.asp} (дата обращения: 14.04.2023).}, что могло ввести в заблуждение человека, постоянно сталкивавшегося с его последним значением, поэтому переводчик решает использовать более точный термин \emph{“}\emph{aventura}\emph{”}\footnote{Aventura s.// Vocabulari de la llengua catalana medieval de Lluís Faraudo de Saint-Germain. URL: \url{http://www.iec.cat/faraudo/results.asp?Word=Aventura&Id=3060252&search=aventura&optCriteria=0&optSearchType=1&ActualPage=1} (дата обращения: 14.04.2023).} в значении «случайное событие», прибегая ещё и к грамматической трансформации: использует множественное число существительного в переводе, чтобы показать, что возможно множество ситуаций с подобными последствиями.

Так в 68(72--73)-й статье «Барселонских обычаев» латинский термин \emph{“strata”} (лат. --- мощёная дорога) заменяет старокаталанское слово \emph{“}\emph{entrada}\emph{”} (кат. --- вход в город, ворота, проход)\footnote{Entrada s.// Vocabulari de la llengua catalana medieval de Lluís Faraudo de Saint-Germain. URL:~\url{http://www.iec.cat/faraudo/results.asp?Word=Entrada&Id=3071720&search=entrada&optCriteria=0&optSearchType=1&ActualPage=1} (дата обращения: 14.04.2023).}: \emph{«Мощёные дороги (лат. }\emph{strat}\emph{e}\emph{---кат.}\emph{ }\emph{entrades}\emph{) }\emph{и публичные пути, текущие воды и живые источники, луга и пастбища, леса и кустарники, а также скалы, находящиеся в этой земле, принадлежат властям…»}\footnote{Usatges de Barcelona. Us. 68(72--73)~: «STRATE et uie publice, aque currentes et fontes uiui, prata et paschue, silue et garrice et roche in hac patria fundate sunt de potestatibus...$\rightarrow$[L]es entrades e les carrercs publiques, les aygues correns e les fons vives, los prats e les pastures, les selves e les gariges, e les roques en aquesta terra fondades...». P.~106--107.}\emph{. }

\emph{Т}\emph{аблица X. Обозначение дорог в латинском тексте и старокаталанских версиях статьи 68 (72--73) «Барселонских обычаев»}

\begin{center}
\small
\setlength{\tabcolsep}{2pt}
\begin{longtable}{|p{\dimexpr0.94\textwidth/3\relax}|p{\dimexpr0.94\textwidth/3\relax}|p{\dimexpr0.94\textwidth/3\relax}|}
\hline
\textbf{Текстовая версия} & \textbf{Оригинальный текст} & \textbf{Русский перевод} \\ \hline
Латинский текст по изданию Ж.~Бастардеса-и-Пареры & лат. \textbf{\emph{stratae}}\emph{ et }\textbf{\emph{viae publicae}}\emph{, aquae currentes et fontes vivi} & \emph{«Мощёные дороги и публичные пути, проточные воды и живые источники…»} \\ \hline
Старокаталанский текст рукописи XIII~в. из ACA, опубликованный Ж.~Бастардесом-и-Парерой & кат. \emph{les }\textbf{\emph{entrades}}\emph{ e les }\textbf{\emph{carreres públiques}}\emph{, les aygues correns e les fons vives} & \emph{«Входы и публичные дороги, проточные воды и живые источники…»} \\ \hline
Старокаталанский текст рукописи, опубликованной Ж.~Гудиолем-и-Кунилем & кат. \textbf{\emph{les strades}}\emph{ e }\textbf{\emph{les vies publiques}}\emph{, les aygues corrents e les fons viues} & \emph{«Мощёные дороги и публичные пути, проточные воды и живые источники…»} \\ \hline
\end{longtable}
\end{center}

В издании Ж.~Бастардеса-и-Пареры латинскому сочетанию \emph{«мощёные дороги и публичные пути»} (лат. --- \emph{stratae et viae publicae}) соответствует старокаталанское \emph{«входы или входы и публичные дороги»} (кат. --- \emph{entrades e carreres públiques}). Таким образом, изменён не только отдельный термин, но и вся парная формула.

Латинское «мощёные дороги» (лат. --- \emph{stratae}) обозначает один из видов дорог и характеризует их по способу устройства. Старокаталанское «входы» или «проходы» (кат. --- \emph{entrades}) имеет иное пространственное значение. Контекст нормы не позволяет точнее определить, какие именно объекты обозначены этим словом: перечисление показывает лишь, что они, наряду с публичными дорогами, водами, пастбищами и лесами, предназначались для общего пользования. Поэтому соответствие «мощёные дороги» (лат. --- \emph{stratae}) --- «входы» (кат. --- \emph{entrades}) нельзя признать точным терминологическим соответствием.

В другой рукописи «Барселонских обычаев» XIII~в. на старокаталанском языке, опубликованной Ж.~Гудиолем-и-Кунилем, латинское сочетание передано почти дословно, с помощью этимологически родственных старокаталанских слов. Следовательно, в старокаталанском языке существовал прямой эквивалент латинского \emph{stratae}.

Таким образом, соответствие \emph{«мощёные дороги»} (лат. --- \emph{stratae}) --- \emph{«входы»} (кат. --- \emph{entrades}) не характеризует старокаталанские переводы «Барселонских обычаев» в целом. Оно засвидетельствовано в тексте, положенном в основу издания Ж.~Бастардеса-и-Пареры, тогда как рукопись, опубликованная Ж.~Гудиолем-и-Кунилем, сохранила более близкое к латинскому чтение. Этот случай следует рассматривать как проявление вариативности старокаталанской рукописной традиции. По двум версиям нельзя определить, присутствовало ли слово «входы, проходы» (кат. --- \emph{entrades}) в первоначальном переводе или появилось при последующей переработке текста.

Мы можем также заметить модуляцию в 28--29 статьях\footnote{Usatges de Barcelona. Us. 28 (32), 29 (33).}, где устойчивая формула \emph{«без согласия сеньора» }(лат. --- \emph{sine}\emph{ }\emph{consensus}\emph{ }\emph{senioris}) переводится как \emph{«без воли сеньора»} (кат. --- \emph{sens}\emph{ }\emph{volentat}\emph{ }\emph{de}\emph{ }\emph{senyor}). Нейтральный термин \emph{“}\emph{consensus}\emph{”}\textsuperscript{ }в старокаталанском заменяется на существительное \emph{“}\emph{volentat}\emph{”}\footnote{Volentat // Vocabulari de la llengua catalana medieval de Lluís Faraudo de Saint-Germain.~URL:~\url{http://www.iec.cat/faraudo/results.asp?Word=Volentat&Id=3097376&search=volentat&optCriteria=0&optSearchType=1&ActualPage=1} (дата обращения: 14.04.2023).}, которое имеет более активную коннотацию собственно желания и волеизъявления. В данном случае мы наблюдаем изменение именно смыслового наполнения как одного используемого термина, так и в целом правовой нормы.

При анализе статей №~5 (6) и 103 (124) также сталкиваемся с похожим приёмом перевода\footnote{Usatges de Barcelona. Us. 5 (6), 103 (124). P.}: \emph{«<…> а именно: что должны они иметь придворную курию и большую компанию <…>»}\footnote{Ibid. Us. 103 (124): «...miles uero qui habuerit duos milites ad homines locatos de suo honore et tenuerit unum de familia sua $\rightarrow$…cavaler, però, qui auri .ii. cavalers ha hòmens loçats de sa honor e·n tendrà un de sa companya». P.}. Латинский термин \emph{“}\emph{familia}\emph{” }(лат. --- семья) переводится на старокаталанский как \emph{“}\emph{companya}\emph{”} (кат. --- группа людей, живущих вместе и ведущих совместное хозяйство)\footnote{Сompanya // Vocabulari de la llengua catalana medieval de Lluís Faraudo de Saint-Germain.~URL:~\url{http://www.iec.cat/faraudo/results.asp?Word=Companya&Id=3065400&search=companya&optCriteria=0&optSearchType=1&ActualPage=1}~(дата обращения: 14.04.2023).}. Таким образом, мы видим, что в данном случае для переводчика важны не родственные связи и отношения между представителями знати, а скорее то, что они происходят из одного рода, то есть трактовка данного понятия сильно расширяется.

Для анализа лексики родства особый интерес представляет 122(80) статья «Барселонских обычаев», в латинской и каталанской версиях которой употребляется несколько обозначений родственников: \emph{«}\emph{…И если кто пожелает повредить принцепсу, то пусть будет наказан и осуждён на все времена вместе со своими потомками [}кат.\emph{---и со всем своим линьяжем]…}\emph{»}\footnote{Usatges de Barcelona. Us. 122 (80): «…et qui principem uult dampnare punitus et dampnatus sit omni tempore, ipse et cuncta sua proienies…---…qui·l príncep vol dampnar, punit e dampnat sia el per tots temps e tot lo seu linatge…». P.~156--157.}\emph{.} В латинской версии использовано существительное\emph{ }\emph{«}потомки\emph{» }(лат. ---\emph{ }\emph{proienies}), в то же время в данном случае каталанская версия предлагает другой термин южно-французского происхождения «линьяж»\emph{ }(кат.~--- \emph{linatge}). Термин «линьяж» имел определённую сферу применения, поскольку чаще всего использовался среди знати для обозначения широкого круга родственников, ведущих происхождение от общего предка. Существование линьяжей в средневековой Каталонии XI---XIII~вв. доказано многочисленными исследованиями с использованием актового материала; многие медиевисты и историки права, работавшие в области истории Каталонии в обозначенный период, так или иначе касались проблематики взаимодействия линьяжей\footnote{Детальные библиография и историографический обзор приведён в работе М.~Ауреля: \emph{Aurell~M}\emph{.} Les noces du comte… P.~25--29.}. Скорее всего, к XIII~в. переводчик мог наблюдать конфликты между графской властью и знатными семействами, и потому использовал модуляцию при интерпретации 122 (80) статьи, контекст которой как раз располагал к тому, чтобы говорить об участии знати в политической жизни принципата и регулировании последней.

При сравнении латинской и старокаталанской версий «Барселонских обычаев» мы можем наблюдать модуляцию и при переводе отдельных юридических терминов. Например, в 60(64) статье «Барселонских обычаев»: \emph{«}\emph{…Постановляем [кат. --- судим] и приказываем, чтобы все принцепсы, которые унаследуют этот приципат [кат. --- которые после нас в этот принципат придут] …}\emph{»}\footnote{Usatges de Barcelona. Us. 103(124): «…decernimus atque mandamus ut omnes principes qui in hoc principatu nobis sunt successuri… --- …jutgam e manam que tots los prínceps que en aquest principat són ni aprés nos vendran…». P.~138--139.}\emph{.} Латинский глагол\emph{ }\emph{«}постановлять» (лат. ---\emph{ }\emph{decernere}) заменён на старокаталанскую форму глагола «судить» (кат. --- \emph{jutgar})\footnote{Jutjar // Vocabulari de la llengua catalana medieval de Lluís Faraudo de Saint-Germain.~URL:~\url{http://www.iec.cat/faraudo/results.asp?Word=Jutjar&Id=3079277&search=jut&optCriteria=0&optSearchType=1&ActualPage=1} (дата обращения: 14.04.2023).}. Значение последнего из них напрямую связано с процессуальными действиями, в отличие от обобщённого законодательного значения у глагола \emph{decernere}. При этом фрагмент ничего не сообщает ни о причинах составления «Барселонских обычаев», ни о процедуре их обнародования или формальной промульгации, поэтому связывать рассматриваемое переводческое решение с этими обстоятельствами оснований нет.

Старокаталанский термин в данном случае уже по значению, чем латинский оригинал. Такая замена, по-видимому, связана не только с особенностями перевода, но и с характером самого латинского текста, в котором местами заметно стремление к архаизирующей юридической стилизации. Поэтому конкретизацию следует рассматривать, с одной стороны, как попытку переводчика сделать норму более определённой и понятной для практического восприятия, а с другой --- как косвенное свидетельство особенностей латинской редакции памятника. По всей видимости, её составители стремились придать тексту форму, присущую авторитетному юридическому своду. Архаизация при этом была лишь средством придания памятнику большей древности и, следовательно, большей легитимности.

Таким образом, сопоставление старокаталанских версий «Барселонских обычаев», опубликованных Ж. Гудиолем-и-Кунилем и Ж. Бастардесом-и-Парерой, с латинским оригиналом позволяет выявить несколько способов смысловой переработки исходного текста --- генерализацию, конкретизацию и модуляцию. Сравнение двух переводов между собой показывает, что выбор соответствия определялся не только значением отдельного латинского слова, но и интерпретацией всей правовой нормы. В одних случаях обе старокаталанские версии одинаково расширяли или сужали значение латинского термина, что свидетельствует об устойчивости определённого решения в переводческой традиции. В других -- они по-разному передавали один и тот же элемент оригинала: один перевод на старокаталанский язык по рукописи, опубликованной Ж.~Гудиолем-и-Кунилем, точнее сохранял объём значения латинского понятия, тогда как перевод по рукописи К, опубликованной Ж.~Бастардесом-и-Парерой, демонстрировал склонность к активному использованию переводческих трансформаций либо передавал норму через иную смысловую перспективу. Тем самым расхождения между версиями отражают не только лексическую вариативность, но и различия в интерпретации содержания правового установления. Усложнённая, а в отдельных случаях архаизированная лексика латинского текста могла расширять пространство возможных переводческих решений, однако сама по себе не предопределяла выбор генерализации, конкретизации или модуляции.

\subsection{Устойчивые терминологические соответствия в латинской и старокаталанской версиях «Барселонских обычаев»}

Наряду со случаями семантического расширения и сужения терминов, в старокаталанском переводе «Барселонских обычаев» отчётливо прослеживается и иная тенденция --- сохранение устойчивых соответствий между латинским и романским текстом. Речь идёт прежде всего о тех юридических понятиях, которые уже обладали достаточной степенью закреплённости в письменной и практической правовой речи и потому не требовали существенной смысловой переработки при переводе. Именно эта группа примеров особенно важна для характеристики старокаталанской версии памятника, поскольку она показывает, что перевод не сводился к упрощению или свободной адаптации латинского оригинала. Напротив, в ряде случаев переводчик последовательно сохранял основные элементы латинского понятийного аппарата, воспроизводя их в форме синтаксически частично адаптированных латинизмов. Анализ таких устойчивых соответствий позволяет точнее определить, какие элементы юридического языка сохраняли авторитет и функциональную значимость при переходе из латинской версии текста в старокаталанскую.

Отдельную группу составляют латинское \emph{mors} и старокаталанское \emph{mort}, которые прямо соответствуют друг другу и употребляются как в статьях о возмещении за убийство, так и в нормах о наследовании. В данном случае речь идёт не о генерализации, а о прямом лексическом и этимологическом соответствии. Конкретное правовое значение слова --- смерть человека, размер компенсации за убийство либо момент открытия завещания --- определяется содержанием соответствующей статьи.

В старокаталанской версии, опубликованной Ж. Бастардесом-и-Парерой, одна форма «мёртвый, умерший» (кат. \emph{mort}) передаёт два различавшихся в латинском тексте обозначения:

-в статьях 10 (12), 11 (13), 76 (102) и 82 (101) она соответствует обозначению лица, погибшего вследствие насильственного действия, --- «убитый» (лат. \emph{interfectus});

-в статье 81 (100) той же формой передано нейтральное обозначение умершего --- «покойный» (лат. \emph{defunctus}).

\emph{Таблица X. Передача обозначений убитого и умершего в старокаталанских версиях «Барселонских обычаев»}

\begin{center}
\small
\setlength{\tabcolsep}{2pt}
\begin{longtable}{|p{\dimexpr0.94\textwidth/5\relax}|p{\dimexpr0.94\textwidth/5\relax}|p{\dimexpr0.94\textwidth/5\relax}|p{\dimexpr0.94\textwidth/5\relax}|p{\dimexpr0.94\textwidth/5\relax}|}
\hline
\textbf{Статья} & \textbf{Латинский текст по изданию Ж. Бастардеса-и-Пареры} & \textbf{Старокаталанский текст по изданию Ж. Бастардеса-и-Пареры} & \textbf{Рукопись по изданию Ж. Гудиоля-и-Куниля} & \textbf{Русский перевод} \\ \hline
10(12) & \emph{Baiulus }\textbf{\emph{interfectus}}\emph{ vel debilitatus sive cesus vel captus…} & \emph{Batle }\textbf{\emph{mort }}\emph{ho debilitat ho trencat ho pres…}\footnote{Usatges de Barcelona. Us. 10(12). P.~60--61.} & Статья отсутствует & «Байл, убитый (кат. ---  мёртвый байл) либо искалеченный, избитый или захваченный…» \\ \hline
11 (13) & \emph{Rusticus }\textbf{\emph{interfectus}}\emph{ seu alius homo…} & \emph{Pagès }\textbf{\emph{mort}}\emph{ ho altre hom…}\footnote{Ibid. Us. 11(13). P.~60--61.} & Статья отсутствует & «Убитый крестьянин (кат. --- мёртвый крестьянин) или другой человек…» \\ \hline
81 (100) & \emph{…veniat in manu proximorum }\textbf{\emph{defuncti}}\emph{…} & \emph{…vinge en la mà dels proïsmes }\textbf{\emph{del mort}}\emph{…}\footnote{Ibid. Us. 81(100). P.~120--121.} & \emph{…venga en man dels prohismes }\textbf{\emph{del hom mort}}\emph{…}\footnote{\emph{Gudiol i Cunill J. }Traducció dels Usatges. Us. 95.\textbf{ }P.~298.} & «…пусть окажется во власти родственников покойного…» \\ \hline
82 (101) & \emph{De composicione omnium hominum qui sunt }\textbf{\emph{interfecti}}\emph{…} & \emph{De la emena dels homens qui }\textbf{\emph{son morts}}\emph{…}\footnote{Usatges de Barcelona. Us. 82(101). P.~120--121.} & \emph{E la composicion de tots homes qui }\textbf{\emph{seran morts}}\emph{…}\footnote{\emph{Gudiol i Cunill J. }Traducció dels Usatges. Us. 96.\textbf{ }P.~298.} & «О возмещении за всех убитых (кат. ---  за людей, которые мертвы / будут мертвы …» \\ \hline
76 (102) & \emph{…quando }\textbf{\emph{interfecti}}\emph{ fuerint…} & \emph{…can }\textbf{\emph{morts}}\emph{ siran…}\footnote{Usatges de Barcelona. Us. 76(102). P.~114--115.} & \emph{…en la esmena dels }\textbf{\emph{ocises}}\emph{…}\footnote{\emph{Gudiol i Cunill J. }Traducció dels Usatges. Us. 97.\textbf{ }P.~298.} & «…когда они будут убиты…»; кат1. -- «когда они будут мертвы»; кат2. -- «при возмещении за убитых» \\ \hline
\end{longtable}
\end{center}

На уровне отдельной лексемы различие между насильственной смертью и смертью как состоянием человека после прекращения жизни оказывается нейтрализовано. В статьях 10(12) и 11(13) речь идёт о размере возмещения за убитого человека определённого статуса; статья 82(101) непосредственно посвящена кровной мести. В статье 81(100) значение слова также устанавливается из упоминания убийства и родственников умершего. Таким образом, сведения о характере смерти передаются всем содержанием нормы.

Вместе с тем сопоставление с рукописью, опубликованной Ж.~Гудиолем-и-Кунилем, не позволяет считать такую нейтрализацию последовательной особенностью старокаталанских переводов «Барселонских обычаев» в целом. В последующих статьях (81(100) и 82(101)) также употреблено «умерший, мёртвый» (кат. ---\emph{mort}), но в соответствующем тексте статьи 76(102) появился более узкий термин «убитые» (кат. --- \emph{ocises}), близкий латинскому \emph{“}\emph{interfecti}\emph{”}\emph{.} Следовательно, старокаталанская рукописная традиция располагала средствами для сохранения семантики насильственной смерти.

Этот пример поэтому свидетельствует не об общем отказе от различений латинской уголовно-правовой лексики, а о лексической нейтрализации в отдельных случаях. Причём правовое различие перенеслось с уровня отдельного слова на уровень контекста всей нормы.

\emph{Та}\emph{блица }\emph{X}\emph{. }\emph{Передача латинских глаголов laedere и vulnerare старокаталанским nafrar}

\begin{center}
\small
\setlength{\tabcolsep}{2pt}
\begin{longtable}{|p{\dimexpr0.94\textwidth/5\relax}|p{\dimexpr0.94\textwidth/5\relax}|p{\dimexpr0.94\textwidth/5\relax}|p{\dimexpr0.94\textwidth/5\relax}|p{\dimexpr0.94\textwidth/5\relax}|}
\hline
\textbf{Статья} & \textbf{Латинский текст } & \textbf{Старокаталанский текст по изданию Ж. Бастардеса-и-Пареры} & \textbf{Рукопись по изданию Ж. Гудиоля-и-Куниля} & \textbf{Русский перевод} \\ \hline
\textbf{55 (58)}. Us. 58--59. & \emph{…si eum in aliquo leserit, emendet ei malum quod fecerit… Et si armata manu aliquem requisierit, }\emph{si non vulneraverit…} & \emph{…si él en re l nafrarà, emèn-li lo mal que li aurà feyt… E si ab mà armada lo requerrà e no·l nafrarà…} & \emph{…si lin fara negun mal, fassa li esmena da quel mal que fet aura… E si alcu requerra altre ab }\emph{armes, si no Ion naffra…} & …если он причинит ему какой-либо вред (кат1 ---ранит его) --- пусть возместит причинённый ущерб… Если же кто нападёт на другого с оружием, но не ранит его… \\ \hline
\textbf{59 (62)}. Us. 63 & \emph{…et si quis illos requisierit, cederit, vulneraverit vel deshonoraverit, aut abstulerit eis aliquid de eorum rebus…} & \emph{…e si algú aquels requerrà, ferrà ho nafrarà ho desonrarà en res, o·ls tolrà res de lurs coses…}\footnote{} & \emph{E si alcun hom requerra ho fara naifra, o encara aquels en alguna cosa desonra, ho les lurs coses lur tol…}\footnote{} & …если кто нападёт на них, ударит, ранит или обесчестит их либо отнимет что-либо из принадлежащего им… \\ \hline
\end{longtable}
\end{center}

Тексты статей 55 (58) и 59 (62) показывают, что в версии, опубликованной Ж. Бастардесом-и-Парерой, слово «ранить» (кат. --- \emph{nafrar}) передаёт как специальное обозначение нанесения раны (лат. --- \emph{vulnerare}), так и более широкий глагол «повреждать, причинять вред» (лат. --- \emph{laedere}). В первом случае старокаталанское слово достаточно близко латинскому по значению; во втором оно конкретизирует характер причинённого вреда, представляя его именно как телесное повреждение.

Однако сравнение с рукописью из Архива Викского епископства не позволяет считать соответствие «причинять вред» (лат. --- \emph{laedere}) --- «ранить» (кат. --- \emph{nafrar}) общим для всей старокаталанской рукописной традиции. В статье 55 (58) версия Ж. Гудиоля-и-Куниля передаёт латинский глагол описательно: «причинит ему какой-либо вред» (кат. --- \emph{fara negun mal}). В то же время более узкое «ранить» (лат. --- \emph{vulnerare}) переводится формой того же старокаталанского глагола --- \emph{naffra}. Следовательно, Викская рукопись сохраняет различие между общим причинением вреда и нанесением раны более последовательно, чем версия Ж. Бастардеса-и-Пареры.

Иная переработка представлена в статье 59 (62). В версии Ж. Бастардеса-и-Пареры латинские действия «ударить» (лат. --- \emph{cedere}) и «ранить» (лат. --- \emph{vulnerare}) переданы раздельно: «ударит или ранит» (кат. --- \emph{ferrà ho nafrarà}). В рукописи из архива Викского епископства, опубликованной Ж.~Гудиолем-и-Кунилем, оба обозначения сведены к конструкции «нанесёт рану» (кат. --- \emph{fara naifra}). Здесь нейтрализуется уже не различие между общим и частным обозначениями вреда, а граница между самим ударом и его возможным результатом --- раной.

Таким образом, слово «ранить» (кат. --- \emph{nafrar}) выступает устойчивым центром обозначения телесного повреждения, но характер его соотношения с латинскими глаголами зависит от рукописной версии. В одном случае оно конкретизирует широкое значение, в другом точно соответствует обозначению ранения, а в Викской рукописи может заменяться описательной конструкцией или конструкцией с существительным «рана» (кат. --- \emph{naifra}). Репрезентативность этих примеров состоит не в установлении неизменной пары для всех списков, а в повторяемости самого способа перераспределения значений между глаголом и контекстом нормы.

\emph{Таблица }\emph{X}\emph{. Передача латинских глаголов caedere и percutere старокаталанским ferrer}

\begin{center}
\small
\setlength{\tabcolsep}{2pt}
\begin{longtable}{|p{\dimexpr0.94\textwidth/5\relax}|p{\dimexpr0.94\textwidth/5\relax}|p{\dimexpr0.94\textwidth/5\relax}|p{\dimexpr0.94\textwidth/5\relax}|p{\dimexpr0.94\textwidth/5\relax}|}
\hline
\textbf{Статья} & \textbf{Латинский текст по изданию Ж. Бастардеса-и-Пареры} & \textbf{Старокаталанский текст по изданию Ж. Бастардеса-и-Пареры} & \textbf{Рукопись по изданию Ж.}\textbf{~}\textbf{Гудиоля-и-Куниля} & \textbf{Русский перевод} \\ \hline
\textbf{5(6)} & \emph{Si quis se miserit in aguayt et consideratamente requisierit militem et cum fuste }\textbf{\emph{cederit}}\emph{ }\emph{eum… }\emph{Si autem aliter quis, quolibet ictu indignans, }\textbf{\emph{cederit }}\emph{militem cum pugno vel calce, cum petra vel fuste…} & \emph{Si algú se metrà en aguayt e acordadament escometrà cavaler e·l }\textbf{\emph{ferrà}}\emph{ ab fust… E si en altra guisa, indignant de }\emph{qualque colp en baralla, }\textbf{\emph{ferrà }}\emph{cavaler ab puny, palmada, o ab pedra, o ab peu, o ab fust…}\footnote{Usatges de Barcelona. Us. 5(6). P.~54--55.} & Статья отсутствует & «Если кто устроит засаду, преднамеренно нападёт на рыцаря и ударит его палкой… Если же кто во время ссоры ударит рыцаря кулаком, ладонью, ногой, камнем или палкой…» \\ \hline
\textbf{13(14)} & \emph{Si quis aliquem }\textbf{\emph{percusserit}}\emph{ in facie, pro alapa dentur V solidos; pro pugno vel calce sive cum petra vel fuste, decem solidos…} & \emph{Si algú }\textbf{\emph{fer}}\emph{ altre en la cara, per galtada cinc sous; e per puny o per peu o per pedra o per fust, deu sous…}\footnote{Ibid. Us. 13 (14). P.~62--63.} & Статья отсутствует & «Если кто ударит другого по лицу, за пощёчину пусть уплатит пять солидов, а за удар кулаком, ногой, камнем или палкой --- десять солидов…» \\ \hline
\textbf{59(62)} & \emph{…et si quis illos requisierit, }\textbf{\emph{cederit}}\emph{, }\textbf{\emph{vulneraverit }}\emph{vel deshonoraverit…} & \emph{…e si algú aquels requerrà, }\textbf{\emph{ferrà}}\emph{ ho }\textbf{\emph{nafrarà}}\emph{ ho desonrarà…}\footnote{Ibid. Us. 59 (62). P.~94--95.} & \emph{E si alcun hom requerra ho }\textbf{\emph{fara}}\emph{ }\textbf{\emph{naifra}}\emph{, o encara aquels en alguna cosa desonra…}\footnote{Gudiol i Cunill J. Traducció dels Usatges…P.~292.} & …если кто нападёт на них, ударит, \textbf{ранит} или обесчестит их… \\ \hline
\textbf{119(77)} & \emph{…si illi tam presumptuosi extiterint ut patrem aut matrem, avum vel aviam }\emph{graviter }\textbf{\emph{percusserint}}\emph{ vel deshonestaverint…} & \emph{…si són tan presumptuosos que·l pare o la mare, l’avi o }\emph{l’àvia greument }\textbf{\emph{firan}}\emph{ o·ls desonren…}\footnote{Usatges de Barcelona. Us. 119 (77). P.~154--155.} & \emph{Los pares, los auis poden deseretar los fils e·ls nets si gran }\textbf{\emph{falia ferran}}\emph{ los }\emph{pares o·ls auis o·ls faran desonor…}\footnote{Gudiol i Cunill J. Traducció dels Usatges... Us. 79. P.~289.} & …если они окажутся настолько дерзкими, что жестоко ударят или обесчестят отца, мать, деда либо бабку… \\ \hline
\end{longtable}
\end{center}

В версии Ж.~Бастардеса-и-Пареры формы слова «ударять» (кат. --- \emph{ferrer}) передают два латинских глагола --- «бить, наносить удары» (лат. --- \emph{caedere}) и «ударять, поражать» (лат. --- \emph{percutere}). При этом латинские слова сами по себе не всегда указывают на конкретный способ нападения. В статье 5~(6) орудие и обстоятельства действия названы отдельно: палка, кулак, ладонь, нога или камень; различаются заранее подготовленная засада и удар во время ссоры. Поэтому старокаталанский глагол сохраняет общее значение физического удара, тогда как юридически значимые характеристики действия выражаются другими элементами нормы.

Та же структура прослеживается в статье 13(14). Общему сказуемому «ударит» соответствует подробная градация способов и результатов воздействия: пощёчина, удар кулаком, ногой, камнем или палкой, появление крови, падение потерпевшего. Размер возмещения зависит не от выбора более узкого глагола, а от названного орудия, части тела и последствия удара. Поэтому использование одного старокаталанского обозначения нельзя трактовать как утрату правового различия: оно сохраняется на уровне всей нормативной конструкции.

Статья 119(77) показывает употребление того же обозначения за пределами норм о денежном возмещении за телесное повреждение. Удар по отцу, матери, деду или бабке рассматривается здесь как основание для лишения наследства. Правовое значение действия определяется сочетанием указания на его тяжесть и особый статус потерпевших. Рукопись, опубликованная Ж.~Гудиолем-и-Кунилем, подтверждает устойчивость этого переводческого решения: латинскому «жестоко ударят» (лат. --- \emph{graviter percusserint}) соответствует старокаталанская конструкция с формой \emph{ferran}.

Вместе с тем статья 59(62) показывает пределы устойчивости данного соответствия. В версии Ж.~Бастардеса-и-Пареры удар и ранение выражены раздельно --- «ударит или ранит» (кат. --- \emph{ferrà ho nafrarà}), тогда как в Викской рукописи они объединены конструкцией «нанесёт рану» (кат. --- \emph{fara naifra}). Следовательно, соответствие латинских глаголов «бить» (лат. --- \emph{caedere}) и «ударять» (лат. --- \emph{percutere}) старокаталанскому «ударять» (кат. --- \emph{ferrer}) устойчиво прежде всего в версии Ж.~Бастардеса-и-Пареры подтверждается по рукописи, опубликованной Ж.~Гудиолем-и-Кунилем, лишь частично.

Рассмотрим ещё один вариант устойчивого соответствия в переводе «Барселонских обычаев»:

\emph{Таблица }\emph{X}\emph{. Употребление обозначений «рогоносец» и «кугуция» в «Барселонских обычаях».}

\begin{center}
\small
\setlength{\tabcolsep}{2pt}
\begin{longtable}{|p{\dimexpr0.94\textwidth/5\relax}|p{\dimexpr0.94\textwidth/5\relax}|p{\dimexpr0.94\textwidth/5\relax}|p{\dimexpr0.94\textwidth/5\relax}|p{\dimexpr0.94\textwidth/5\relax}|}
\hline
\textbf{Статья} & \textbf{Латинский текст} & \textbf{Старокаталанский текст по рукописи K, изданной Ж. Бастардесом-и-Парерой} & \textbf{Старокаталанский текст }\textbf{по рукописи, }\textbf{изданной Ж. Гудиолем-и-Кунилем} & \textbf{Русский перевод} \\ \hline
Us. 1 (1--2) & Homicidium uel cugucia, que non possunt neclectari, sunt secundum leges et mores iudicata et emendata siue uindicata. & Homicidi o cugúcia, que no poden ésser menyspreats, sien segons lig e costumes jutjats e esmenats o venjats\footnote{} & Статья отсутствует & Убийство или кугуция, которые не могут быть оставлены без внимания, должны быть согласно законам и обычаям рассмотрены судом, возмещены либо отомщены». \\ \hline
Us. 70 (75) & …uel appellauerit aliquem cuguz… & …o l’apelarà cuguç…\footnote{} & …o apelara altre caguz…\footnote{} & …либо назовёт другого рогоносцем… \\ \hline
Us. 87 (110) & Similiter de rebus et possessionibus cucuciorum, si earum maritis nolentibus erit facta cugucia… Si uero […] maritis uolentibus uel precipientibus uel assentientibus fuerit facta ipsa cucucia… & Semblantment de les coses e de les possessions dels cuguços, si la cugúcia és feta los marits no volents… E si […] ab volentat o ab manament o ab consentiment del marit serà feta la cugúcia…\footnote{} & Exament de les coses et de les possessions dels cuguces, si la cugucia es feta los marits no uolens… E si ab uolentat e ab manament et ab consent del marit sera feta la cugucia…\footnote{} & Подобным образом, относительно имущества и владений рогоносцев: если кугуция совершена против воли мужей… Если же кугуция совершена по воле, приказанию или с согласия мужа… \\ \hline
\end{longtable}
\end{center}

Сопоставление статей показывает последовательное юридическое развёртывание одной словообразовательной основы: обозначение «рогоносец» (лат. --- cuguz; старокат. --- cuguç, caguz) называет человека и употребляется как оскорбление, тогда как производное «кугуция» (лат. --- cucucia, cugucia; старокат. --- cugúcia, cugucia) обозначает уже значимый с правовой точки зрения случай супружеской неверности, влекущий имущественные последствия\footnote{Usatges de Barcelona. Us. 1 (1--2), 70 (75), 87 (110).}. Формы cucuciorum и cuguços / cuguces показывают, что от той же основы образовывалось наименование мужей женщин, уличённых в прелюбодеянии.

В латинском тексте «Барселонских обычаев» данный термин каталанского происхождения получил латинизированные падежные окончания и был вписан в латинский текст морфологическими средствами\footnote{\emph{Prieto Espinosa C.A}. Op.cit.Vol.~I.P.~72.}. По нашему мнению, справедливо считать использование данного термина в рамках латинского текста «Барселонских обычаев» проявлением переключения кодов, когда слово из языка, которым носитель пользуется в повседневной жизни, проникает в текст, создаваемый им же на языке официальном.

В каталонской традиции XI--XIII~вв. этот термин также вошёл в число «дурных обычаев»\footnote{\textbf{Кугусия} (лат. \emph{cugutia}) --- право сеньора на долю имущества крестьянина в случае супружеской измены его жены, трактуемой как оскорбление чести сеньора. \emph{Пискорский~В.К.} Вопрос о значении и происхождении шести «дурных обычаев» в Каталонии: С прил. неизд. документов. Киев: тип. Ун-та св. \emph{Владимира Н.Т.}~Корчак-Новицкого, 1899. С.~8--9; \emph{Варьяш~И.И., Астахов~М.А.} Арагонская Корона в XV~в. \emph{// }История Испании / под ред. \emph{В.А.~Ведюшкина}, \emph{Г.А.~Поповой}. М.: «Индрик», 2012. Т.~1. С.~346.}. При этом его история особенно показательна, поскольку он связан с вернакулярным обозначением \emph{«}рогоносец» и демонстрирует, как слово разговорного происхождения могло стать основой для полноценного юридического термина.

Примеры репрезентативны для выявленного способа перевода, поскольку охватывают разные правовые ситуации: нападение на рыцаря, возмещение за телесные повреждения, охрану передвижения по публичным дорогам и лишение наследства за насилие над родителями. Во всех случаях старокаталанское слово обозначает общий тип физического действия, а его способ, тяжесть, результат и правовые последствия раскрываются контекстом статьи. Однако сопоставление с рукописью, опубликованной Ж. Гудиолем-и-Кунилем, показывает, что этот способ перевода следует рассматривать как повторяющуюся тенденцию в рукописной традиции.

Обратимся к анализу переводческой эквивалентности в других случаях. В двух статьях латинский текст использует разные наречия для характеристики осознанного и неслучайного действия: «обдуманно» (лат. \emph{consideratamente}) и «сознательно» (лат. \emph{pensabiliter}). В старокаталанской версии, опубликованной Ж.~Бастардесом-и-Парерой, оба слова переданы одним наречием --- «преднамеренно» (кат. \emph{acordadament}). Следовательно, речь идёт не о появлении в переводе значения умысла, отсутствовавшего в латинском тексте, а о сведении двух близких латинских обозначений к одному устойчивому соответствию.

\emph{Таблица X. Передача обозначения преднамеренного действия в латинской и старокаталанской версиях «Барселонских обычаев»}

\begin{center}
\small
\setlength{\tabcolsep}{2pt}
\begin{longtable}{|p{\dimexpr0.94\textwidth/5\relax}|p{\dimexpr0.94\textwidth/5\relax}|p{\dimexpr0.94\textwidth/5\relax}|p{\dimexpr0.94\textwidth/5\relax}|p{\dimexpr0.94\textwidth/5\relax}|}
\hline
\textbf{Статья} & \textbf{Латинский текст} & \textbf{Старокаталанский текст по изданию Ж. Бастардеса-и-Пареры} & \textbf{Р}\textbf{укопись по изданию Ж. Гудиоля-и-Куниля} & \textbf{Перевод старокаталанского текста} \\ \hline
5(6) & \emph{Si quis se miserit in aguait et }\textbf{\emph{consideratamente}}\emph{ requisierit militem..}\emph{.} & \emph{Si algú se met en aguayt, e }\textbf{\emph{acordadament }}\emph{requerrà cavaler}\emph{…}\footnote{Usatges de Barcelona. Us. 5(6). P.~54--55.} & Статья отсутствует. & Если кто-либо устроит засаду и обдуманно (кат. --- \textbf{преднамеренно}\textbf{)} нападёт на рыцаря \\ \hline
36(39) & \emph{Qui seniorem suum despexerit et per superbiam eum }\textbf{\emph{pensabiliter }}\emph{diffidauerit}\emph{…} & \emph{Qui son senyor menysprearà e per ergull él }\textbf{\emph{acordadament}}\emph{ desfiarà}\emph{…}\footnote{Ibid. Us. 36(39). P.~80--81.} & \emph{Si alcun son senyor menyspreara ho per erguyl }\textbf{\emph{acordadament}}\emph{ la cundara}\emph{…}\footnote{\emph{Gudiol i Cunill J. }Traducció dels Usatges. Us. 38.\textbf{ }P.~289.} & «Кто проявит презрение к своему сеньору и, движимый гордыней, сознательно (кат. ---\textbf{преднамеренно}\textbf{)} откажет ему в верности… \\ \hline
\end{longtable}
\end{center}

В статье 5(6) наречие отличает заранее подготовленное нападение из засады от удара, нанесённого при иных обстоятельствах. В статье 36(39) оно характеризует сознательный отказ вассала от верности сеньору. В обоих случаях обдуманный характер действия имеет непосредственное значение для его правовой квалификации и последствий. Старокаталанский перевод поэтому закрепляет единое обозначение для преднамеренного действия.

Рукопись, опубликованная Ж.~Гудиолем-и-Кунилем, подтверждает устойчивое употребление наречия «преднамеренно» (кат. \emph{acordadament}) в статье 36(39) «Барселонских обычаев». Поскольку статья 5(6) в сохранившейся части этой рукописи отсутствует, распространение данного соответствия на оба контекста можно установить только по старокаталанскому переводу «Барселонских обычаев», опубликованному Ж.~Бастардесом-и-Парерой.

Вместе с тем даже в тех случаях, когда переводчик явно стремится удержать латинскую терминологическую основу, соответствие между двумя версиями текста не всегда оказывается полным. Иногда сохранение самого термина или словообразовательной модели сочетается с изменением его грамматического статуса и способа включения в норму. Показателен в этом отношении текст 117-й (138-й) статьи «Обычаев»:\emph{ «О лицах, умерших без завещания (}лат. \emph{intestatis}\emph{---}кат. \emph{D}\emph{’}\emph{aquels}\emph{ }\emph{qui}\emph{ }\emph{passen}\emph{ }\emph{sens}\emph{ }\emph{testament}\emph{) <…> То же следует применять и к женщинам, умершим без завещания (лат. }\emph{intestatis}\emph{---кат. }\emph{qui}\emph{ }\emph{moren}\emph{ }\emph{entestades}\emph{), как выше сказано то о мужчинах»}\footnote{Usatges de Barcelona. Us. 117(138): «DE INTESTATIS ab hoc seculo discessis <…> Ita sit de uxoribus intestatis quemadmodum dicitur superius de uiris$\rightarrow$D’aquels qui passen sens testament d’aquest segle <…> atresí sia dc les mulers qui moren entestades com és dels marits”». P.~150--151.}.

Нас интересует окончание статьи. В латинском варианте мы видим, что термин \emph{“}\emph{intestat}\emph{us}\emph{”} применяется, чтобы обозначить лицо, умершее без завещания, и в первом случае даже субстантивируется; во втором случае мы встречаем старокаталанское причастие “\emph{entestades}\emph{”}, образованное в рамках старокаталанской грамматики, однако по латинской модели, через отрицательную приставку лат. \emph{--in} кат. \emph{--}\emph{en}\emph{.} Важно, что это нерегулярное явление. Каталанское причастие, в отличие от латинского, не является самостоятельным, оно вводится в текст в качестве придаточного определительного, в состав которого включена личная форма глагола \emph{“}\emph{morir}\emph{”}. Это указывает на то, что понимание части смыслов причастия в правовых представлениях каталонских юристов и переводчиков текста «Барселонских обычаев» было утрачено, а сам термин требовал ряда дополнительных пояснений, которые включались в текст посредством придаточных предложений и усложнения синтаксической структуры текста статьи.

Наблюдать, как переводчик добавляет нехарактерные окончания к словам латинского происхождения и как это нарушает структуру фразы, можно и в 73-й статье «Барселонских обычаев»: «\emph{…поскольку творить правосудие над злодеями предоставлено (лат.}\emph{ }\emph{datum}\emph{ }\emph{est}\emph{---кат. }\emph{dades}\emph{ }\emph{s}\emph{ó}\emph{n}\emph{) лишь властям…»}\footnote{Usatges de Barcelona. Us. 73 (94): «…QVIA IVSTICIAM facere de malefactoribus datum est solummodo potestatibus …$\rightarrow$...Car fer justícia dels malsfeytors dades solament són a les postats...». P.~112--113.}\emph{.} Синтагма \emph{“}\emph{datum}\emph{ }\emph{est}\emph{”} переходит во множественное число \emph{“}\emph{dades}\emph{ }\emph{s}\emph{ó}\emph{n}\emph{”}, однако нетрудно заметить ошибку с точки зрения согласования подлежащего и сказуемого. Являющееся сказуемым словосочетание \emph{“}\emph{iusticiam}\emph{ }\emph{facere}\emph{” }или \emph{“}\emph{fer}\emph{ }\emph{just}\emph{ícia”} не может получить вид множественного числа ни в латинской, ни в каталанской грамматике, поскольку выражено инфинитивом глагола и существительным. Однако в данном случае переводчик согласовывает его либо со словом \emph{“}\emph{malsfeytors}\emph{”}, либо со словом \emph{“}\emph{postats}\emph{”}, которые стоят во множественном числе. Невозможно определить, является ли это ошибкой по невнимательности или свидетельством непонимания грамматических отношений между членами оригинального латинского предложения в стремлении сохранить его оригинальную синтаксическую и лексическую структуру.

С другой стороны, в ряде случаев терминология попадала практически без изменений в старокаталанский перевод «Барселонских обычаев». Это касается, например, термина \emph{“}\emph{assalt}\emph{”} в 6(7--8) статье этого сборника установлений обычного права:\emph{ «За гайту, энкаль}\emph{с}\emph{ и штурм (лат. }\emph{assalt}\emph{---кат. }\emph{assalt}\emph{) замка надлежит }\emph{возмещать [ущерб.---А.К.] посредством оммажа и алискары»}\footnote{Ibid. Us. 6 (7--8)~: «GAITA ET ENCALZ de cauallario et assalt de castello emendetur per hominiaticum et per aliscaram $\rightarrow$ [A]guayt e encalç de cavaler e assalt de casteyl sie emenat per homanatge ho per [al]is[ca]ra». P.~56--57.}. Он обозначал нападение на замок противника, данный термин, в народной латыни произошедший от второго супина глагола \emph{“}\emph{assalire}\emph{” --- “}\emph{assaltu}\emph{”}\footnote{Assaltus 1 // \emph{d}\emph{u Cange C. du Fresne et al.} Glossarium mediae et infimae latinitatis / éd. augm. par les bénédictins de Saint-Maur. Niort : L. Favre, 1883--1887. T.~1. Col. 427b. URL: \url{http://ducange.enc.sorbonne.fr/ASSALTUS1} (дата обращения: 29.03.2023).} был распространён в Западном Средиземноморье в эпоху Высокого и Позднего Средневековья , это один из примеров прямого заимствования без изменения форм термина из народной латыни старокаталанским языком.

Рассмотренные случаи не исчерпывают всех терминологических соответствий между латинским и старокаталанским текстами «Барселонских обычаев», однако представляют собой не случайную подборку разночтений, а примеры повторяющихся способов передачи правовой лексики. Каждое соответствие прослеживается в нескольких статьях и нормативных контекстах. Поэтому приведённые случаи репрезентативны не для всего словаря памятника, а для модели перевода, при которой нескольким близким латинским лексемам регулярно соответствует одно старокаталанское слово.

Чтобы установить, не являются ли обнаруженные соответствия особенностью рукописной версии, положенной в основу издания Ж. Бастардеса-и-Пареры, старокаталанский текст был сопоставлен с рукописью из Архива Викского епископства, опубликованной Ж. Гудиолем-и-Кунилем. Совпадение переводческих решений в разных списках позволяет говорить об устойчивой тенденции переводческой традиции. Расхождения, напротив, следует рассматривать как варианты отдельных рукописных редакций, а не приписывать единому переводчику «Барселонских обычаев».

Сопоставление показывает, что устойчивость соответствий не предполагала буквального воспроизведения лексической структуры латинского текста. Несколько обозначений смерти передавались словом «мёртвый» или «смерть» (кат. --- \emph{mort}), глаголы причинения телесного повреждения --- словом «ранить» (кат. --- \emph{nafrar}), а обозначения удара --- глаголом «ударить» (кат. --- \emph{ferrer}). Старокаталанский перевод, таким образом, частично нейтрализовал различия между близкими латинскими словами и формировал более регулярную систему соответствий.

Эту нейтрализацию нельзя интерпретировать как упрощение правового содержания. Значимые различия между действиями и их последствиями сохранялись на уровне всей нормы --- в указаниях на обстоятельства, орудие, тяжесть повреждения, положение потерпевшего и санкцию. Следовательно, переводческая традиция не только закрепляла повторяющиеся способы передачи правовой лексики, но и перераспределяла значения между отдельным словом и контекстом. Сопоставление изданий Ж.~Бастардеса-и-Пареры и Ж.~Гудиоля-и-Куниля позволяет рассматривать эту модель как свойство более широкой рукописной традиции, сохранявшей при этом вариативность отдельных списков.

\subsection{Синтаксическая переработка латинского оригинала}

Важной особенностью старокаталанского перевода «Барселонских обычаев» является не только переработка лексики, но и глубокое изменение синтаксической организации правового текста. При сопоставлении двух версий памятника становится очевидно, что перевод затрагивает сам способ выражения нормы: компактные и во многом абстрактные конструкции латинского оригинала последовательно разворачиваются в более аналитические и разъяснительные структуры. Это проявляется в передаче причастных и герундивных форм, в изменении способов выражения долженствования, а также в переходе от номинативного стиля к глагольному описанию действия. В результате старокаталанский текст оказывается в большей степени ориентирован на конкретную ситуацию, субъект и результат правового действия. Такой сдвиг позволяет рассматривать перевод как синтаксическую и функциональную адаптацию латинского оригинала к иной правовой и коммуникативной среде.

При этом, как отмечает немецкий лингвист, занимавшийся количественными исследованиями структуры переводов «Барселонских обычаев», Р.~Мейер-Германн, синтаксис переводов на старокаталанский язык XIII~в. очень чётко следовал латинскому оригиналу, в том числе и в области порядка глаголов при образовании сложных времён и конструкций. Так, латинский вариант предполагал следующий порядок: инфинитив + вспомогательный глагол, причастие + вспомогательный глагол; старокаталанский -- наоборот. Как показали исследования Р.~Мейера-Германна, в 31,9 \% случаев с инфинитивными конструкциями в старокаталанском тексте из рукописи K\footnote{По принятой в научной традиции исследования «Барселонских обычаев» номенклатуре рукописей, та, с которой работали Р.~Мейер-Германн и Х.~Бастардес-и-Парера имеет следующие данные: ACA. Cancillería. Papeles por incorporar. Legislació. Caixa 1. Núm.1.} используется характерный для латинского языка порядок глаголов (против 41\% в оригинальном латинском тексте), для оборотов с причастиями данные соответственно 12,8 \% для характерен латинских конструкций (в оригинале -- 40,6\%)\footnote{\emph{Meyer-Hermann R.} Auxiliar-Konstruktionen im llatí medieval katalanischer…S.~507--534.}. По мнению исследователя, это показывает, что на момент осуществления переводов «Барселонских обычаев» на старокаталанский язык, последний ещё не полностью вытеснил латынь из языка законодательных актов на территории Каталонии и этот процесс был далёк от завершения\footnote{\emph{Idem.} Zur Emergenz katalanischer Syntax in den Usatici Barchinonae / Usatges de Barcelona (12.---15. Jhdt.). // Zeitschrift für Katalanistik. 2020. Bd. 33. S. 185.}. Порядок глаголов в подобных конструкциях и его изменение в старокаталанском переводе уже подробно исследованы лингвистами. Отметим лишь, что выводы Р.~Мейера-Германна о «Барселонских обычаях» согласуются с результатами нашего анализа синтаксической структуры перевода.

Кроме того, при сравнении латинского текста с его старокаталанским переводом мы находим, что некоторые формы глаголов, выраженные в латинском тексте первым лицом множественного числа как бы от лица составителя кодекса, в переводе выражены множественным числом, но уже третьего лица: \emph{«Отсюда постановляем, что если кто сеньору своему поклянётся в чём-то, что не захочет выполнять <…>»}\footnote{Usatges de Barcelona. Us. 63(67): «ITEM STATVIMVS ut si quis seniori suo iurauerit aliquid quod tenere non curet... $\rightarrow$[D]e cap establíren que si algú a son senyor jurarà alguna cosa que tenir no vula...». P.~100--101.}.

Мы встречаемся с аналогичным явлением ещё в ряде статей, рассмотрим это явление подробнее.

\emph{Таблица }\emph{X}\emph{. Сопоставление форм лица и способов выражения предписания в латинской и старокаталанских версиях «Барселонских обычаев».}

\begin{center}
\small
\setlength{\tabcolsep}{2pt}
\begin{longtable}{|p{\dimexpr0.94\textwidth/5\relax}|p{\dimexpr0.94\textwidth/5\relax}|p{\dimexpr0.94\textwidth/5\relax}|p{\dimexpr0.94\textwidth/5\relax}|p{\dimexpr0.94\textwidth/5\relax}|}
\hline
\textbf{Статья} & \textbf{Латинский текст по изданию Ж. Бастардеса-и-Пареры} & \textbf{Старокаталанский текст по изданию Ж.}\textbf{~}\textbf{Бастардеса-и-Пареры} & \textbf{Рукопись по изданию Ж. Гудиоля-и-Куниля} & \textbf{Единый русский перевод с указанием разночтений} \\ \hline
\textbf{60(64)} & \emph{…}\textbf{\emph{decernimus atque mandamus}}\emph{ ut }\emph{omnes principes qui in hoc principatu nobis sunt successuri…} & \emph{…}\textbf{\emph{jutgam e manam}}\emph{ que tots los prínceps que en aquest }\emph{principat són ni aprés nos vendran…}\footnote{Ibid. Us. 60 (64). P.~96--97.} & \emph{…nos damont dits princeps En Ramon e Na }\emph{Alamur […] }\textbf{\emph{ordonam e manam e stablim}}\emph{ que tots los princeps qui seran en aquesta terra aquels qui vendran apres dels…}\footnote{\emph{Gudiol i Cunill J. }Traducció dels Usatges. Us. 65--66. P.~292--293.} & …мы \textbf{постановляем и приказываем }(кат1.--- \textbf{мы судим и приказываем}… [тем, кто] придёт после нас; кат2. --- мы \textbf{распоряжаемся, приказываем и устанавливаем}… [относительно] тех, кто придёт после них), чтобы все принцепсы, которые станут нашими преемниками в этом принципате… \\ \hline
\textbf{62 (66.2)} & \emph{…ita }\textbf{\emph{stabiliendo precipimus}}\emph{ ut persone eorum cum omni }\emph{honore et auere ueniant in manu principis…} & \emph{…aquestes coses }\textbf{\emph{establim, manan}}\emph{ que les persones d’aquels, ab tota lur honor e ab lur aver, }\emph{vingen en la man del príncep…}\footnote{Usatges de Barcelona. Us. 62 (66.2). P.~98--101.} & \emph{…}\textbf{\emph{stablim e manam}}\emph{ que les persones daquels ab tota la honor et ab }\emph{lauer vagen en la ma del princep…}\footnote{\emph{Gudiol i Cunill J. }Traducció dels Usatges. Us. 67. P.~293.} & …\textbf{устанавливая }это, \textbf{мы предписываем}\textbf{ }(кат1.\textbf{--- }мы устанавливаем, и они приказывают; кат2. --- мы устанавливаем и приказываем, чтобы они вместе со всеми своими владениями и имуществом перешли во власть принцепса… \\ \hline
\textbf{63 (67)} & \textbf{\emph{Item statuimus}}\emph{ ut si quis seniori suo iurauerit aliquid quod tenere non curet…} & \textbf{\emph{De cap establíren}}\emph{ que si algú a son senyor jurarà alguna cosa que tenir no vula…}\footnote{Usatges de Barcelona. Us. 63 (67). P.~100--101.} & \textbf{\emph{Establim e manam}}\emph{ que si alcun jurara a son senyor alcuna cosa que jurada li aura tenir no li vuyla…}\footnote{\emph{Gudiol i Cunill J. }Traducció dels Usatges. Us. 68. P.~293.} & Кроме того, мы постановляем (кат1. --- вновь постановили; кат2. --- мы постановляем и приказываем) если кто-либо поклянётся своему сеньору в чём-либо, чего не пожелает соблюдать… \\ \hline
\textbf{96 (119)} & \emph{Solidos de composicione arborum incisarum aliquociens }\textbf{\emph{precipimus}}\emph{ esse aureos, aliquociens denarios……et hoc}\emph{ }\textbf{\emph{concedimus}}\emph{ in arbitrium iudicis crescere uel minuere…} & \emph{Lo sou de la composició dels arbres trencats, a la vegada }\textbf{\emph{mane·s}}\emph{ éser d’or, a la vegada de diners……e açò }\textbf{\emph{atorguen}}\emph{ que sie en albir del jutge de crexer ho de minvar aquesta composició…}\footnote{Usatges de Barcelona. Us. 96 (119). P.~132--133} & \emph{La esmena de la composicio dels arbres trencats, Auegades }\textbf{\emph{mana}}\emph{ esser de ce axi con la lig ho ensenya. Auegades de diners…E }\textbf{\emph{atorga}}\emph{ que sia esgart del jutge de crexer, o de menuar aquesta composicio…}\footnote{\emph{Gudiol i Cunill J. }Traducció dels Usatges. Us. 114. P.~301.} & Мы предписываем, чтобы солиды, уплачиваемые в возмещение за срубленные деревья, иногда были золотыми, а иногда выплачивались денариями; право увеличить или уменьшить возмещение мы предоставляем судье. \\ \hline
\end{longtable}
\end{center}

Обратимся к тексту 2(3)-й статьи «Барселонских обычаев», традиционно относимой исследователями к одной из «преамбул» к основной части данного сборника правовых установлений: \emph{«Когда господин Рамон Беренгер Старый, граф и маркиз Барселоны и также покоритель Испании, обрел владения, [кат.---я.---А.К.] увидел и узнал, }\emph{что во всех судебных случаях и делах этой страны законы готов не могут соблюдаться, и еще увидел множество жалоб и судебных решений <…>»}\footnote{Usatges de Barcelona. Us. 2 (3): «CVM DOMINVS Raimundus Berengarii uetus, comes ei marchio Barchinone atque Ispanie subiugator, habuit honorem et uidit et cognouit quod in omnibus causis et negociis ipsius patrie leges gotice non possent obseruari …$\rightarrow$[Q]uan lo senyor Ramon Berenguer, veyl compte e marquès de Barçelona e subjus[...]ador d’Espanya, ac honor e viu e conec que les ligs godes no podien ésser gardades [e]n tots los pleyts e en totes les faenes d’aquesta terra…». P.~50--51.}\emph{.}

В данном случае два из трёх латинских глаголов в форме 3-го лица единственного числа переводятся на старокаталанский формами 1-го лица единственного числа (\emph{“}\emph{aver}\emph{” }и \emph{“}\emph{coinexer}\emph{”}), в то же время форму 3-го лица единственного числа сохраняет глагол \emph{“}\emph{viure}\emph{”}. В данном случае нас не интересует произошедшее в старокаталанском тексте упрощение временных форм, которое привело к смене перфектных форм формами настоящего времени. Смена лица в данном случае могла быть сделана переводчиком, чтобы подчеркнуть, что данные статьи являются частью «преамбулы» и имеют легитимирующий характер по отношению ко всему кодексу. Однако данное изменение фиксируется лишь однажды.

В дополнение к приведённым выше соображениям стоит отметить, что на протяжении всего текста «Барселонских обычаев», а не только касающихся деятельности и положения правителя, используется притяжательное местоимение \emph{“}\emph{nos}\emph{tres}\emph{”}\footnote{Ibid. Us.60 (64), 70 (75). P.~96--97, 108--109.}, вероятно, поскольку субъект, к которому оно относится, определить довольно трудно: с одной стороны, это могут быть вышеназванные принцепсы и их наследники, с другой же, всё население графства, названное \emph{«жителями»}\emph{ }(кат.\emph{ --- }\emph{(}\emph{h}\emph{)}\emph{abitadors})\footnote{Ibid. Us.60 (64). P.~96--97.}.

В данном случае такая замена показательна, ведь встречая такую форму мы понимаем, что переводчик очень хорошо осознавал вторичность создаваемого текста, поэтому как бы подразумевал: \emph{«В тексте, который я перево}\emph{ж}\emph{у, устанавливается, что <…>»}. Это обстоятельство позволяет пролить свет на его отношение к переводимому, на некие устойчивые представления о «Барселонских обычаях», существовавшие в обозначенный период, трепет и почтение, с которым к данному тексту относились легисты. Переводчик, по нашему мнению, вполне осознавал, что с «оригинальным» текстом «Барселонских обычаев», покрытым ореолом важности и даже «сакральности», ничего не происходит.

При переводе практически всех статей «Барселонских обычаев» переводчик или переводчики стремились сохранить исходные наклонения и времена глаголов (в большинстве случаев это конъюнктив со значением долженствования)\footnote{Подробнее об этом на португальском материале см. \emph{Варьяш О.И.} Язык средневекового права... С.~90.}. Для передачи значения в рамках этой тенденции в отдельных случаях используются уже новые аналитические формы, характерные для романских языков. В качестве иллюстрации приведём фрагмент текста 50(53)-й статьи «Барселонских обычаев»: \emph{«Прочим же крестьянам, называемым баккалариями}\footnote{\textbf{Баккаларии} --- особый разряд крестьян «Барселонских обычаях», противопоставленный державшим манс и обрабатывавшим его парой быков; их свидетельские показания допускались лишь по менее значительным тяжбам (подробнее об истории термина см. \emph{Филиппов И.С. }Средиземноморская Франция в раннее средневековье: Проблема становления феодализма. Москва: Скрипторий, 2000. С.~497).}\emph{, доверяют }(лат\emph{.}\emph{ }\emph{credantur}---кат. \emph{sien}\emph{ }\emph{creeguts}\emph{) на основании клятвы лишь до суммы в четыре манкуса}\footnote{Манкус // Нумизматический словарь / ред. \emph{Зварич В.В. }/ Львов, 1980. C.~104--105: «\textbf{Манкус} --- первоначально название арабского золотого динара (приблизительно 30 серебряным денариям) в Западной Европе, а затем --- денежно-счетная единица и золотая монета ряда европейских стран. <…> В Каталонии в XI~в. выпускался золотой манкус с арабской надписью и крестом. Вес монеты 1,9 - 1,95~г. В большинстве случаев она легче арабского золотого динара».}\emph{ валенсийского }\emph{золота}\emph{. В}\emph{сё, что свыше, в отношении чего они приносят клятву, должно быть доказано посредством испытания кипятком»}\textsuperscript{ }\footnote{Usatges de Barcelona. Us. 50 (53): «DE ALIIS NAMQVE rusticis qui dicuntur bachallarii, credantur per sacramentum usque ad .IIII. mancusos auri Valencie$\rightarrow$Dels altres pages[es] qui an nom bacalars, sien creeguts lur sagrament entrò ha .iiii. maçmodines d’aur de València». P.~88--89.}\emph{.}

Однако в латинском тексте «Обычаев» присутствуют и другие способы выражения значения долженствования, которые получают в переводе своё выражение:

\begin{itemize}
\item через глагол «быть должным» \emph{“}\emph{deber}\emph{”} --\emph{ “}\emph{deuer}\emph{”. }Это наиболее частотный случай, в основном сохраняется в старокаталанском переводе дословно, на нём мы останавливаться не будем;
\item посредством пассивного описательного спряжения -- специальной аналитической формы латинского глагола со значением долженствования, состоящей из формы герундива и личной формы глагола «быть»\emph{ }(лат. ---\emph{ }\emph{esse}\emph{)}. В старокаталанском переводе передаётся традиционно через глагол «быть должным» (кат. --- \emph{deuer}), т.е. аналогичное лексическое значение предложения передаётся иными грамматическими средствами.
\end{itemize}
На этом мы остановимся отдельно и приведём ряд примеров, чтобы продемонстрировать присутствие грамматической трансформации, свидетельствующей об определённой склонности в восприятии латинского текста переводчиком. Такое явление наблюдается в 49-й\footnote{Usatges de Barcelona. Us.49 (52): «SACRAMENTA RVSTICI qui teneat mansum et laboret cum pare bouum, \textbf{sunt}\textbf{ }\textbf{credenda} usque ad .VII. solidos plate$\rightarrow$Sagrament de pageses qui tinge mas e laru ab .i. parel de bous \textbf{deu}\textbf{ é}\textbf{ser}\textbf{ }\textbf{creegut}\textbf{ }entrò ha .vii. sous de plata». P.~88--89.}, 92-й\footnote{Ibid. Us. 92 (115): «...sin autem, \textbf{expectandum est }usquequo pupilli sint talis etatis... $\rightarrow$..si no, \textbf{esperar deu hom} entrò que·ls pobils sien de tal e<s>dat...». P.~128--129.} и 123-й\footnote{Ibid. Us. 123 (81): «..Et ideo \textbf{facienda} [опущен инфинитив “est”.--- А.К.] que sunt secundum usaticum... $\rightarrow$...E per ço \textbf{deu hom fer} ço que és segons l’usatge...». P.~158--159.} статьях «Барселонских обычаев». Несмотря на это, в 119(77)-й статье мы видим перевод пассивного описательного спряжения уже со значением возможности через глагол \emph{“po}\emph{der}\emph{”}, что говорит о широких возможностях для различных переводческих трактовок данной конструкции и о полисемантичности отдельных грамматических единиц в латинском тексте «Барселонских обычаев»: \emph{«…если тот же самый дед или бабка, отец или мать захотят, должны быть отстранены }(лат. \emph{sunt}\emph{ }\emph{repellendi}---кат.\emph{ }\emph{los}\emph{ }\emph{poden}\emph{ }\emph{gitar}\emph{ }\emph{de}\emph{ }\emph{si}\emph{~}\emph{)…»}\footnote{Ibid. Us. 119 (77): «...si idem auus uel auia, pater uel mater uoluerint, \textbf{sunt repellendi}$\rightarrow$...si lo pare o la mare, l’avi o l'avia se volran, \textbf{los poden gitar de si}». P.~154--155.}.\textsuperscript{ }

Если говорить о переводе неличных форм глагола в «Барселонских обычаях», то необходимо подчеркнуть, что и инфинитивы, и герундив в составе герундивной конструкции заменяются инфинитивом соответствующего глагола с предлогом \emph{a}, например форма «для осуществления» (лат. --- \emph{faciendam}) передаётся сочетанием «для совершения» (кат. --- \emph{a fer})\footnote{Ibid. Us. 37 (40), 62 (66.2), 80 (105).} и т. д. Зачастую такие конструкции переводились с помощью личных форм глаголов, как, например, в 118(76)-й статье: \emph{«…постановили, чтобы такого рода вышеуказанное держание <…> обладало бы (лат. }\emph{obtinere}\emph{---кат. }\emph{deu}\emph{ }\emph{aver}\emph{)}\emph{ во всём такой прочностью, что впредь никакой обманной хитростью не могло бы быть подорвано или каким-либо способом изменено»}\footnote{Ibid. Us. 118 (76): «…constituerunt huiusmodi pretaxatam tenedonem<…> talem in omnibus firmitatem obtinere, quod deincebs ulla fraudulenta calliditate subuerti, uel aliquo ingenio non possit mutari$\rightarrow$...establiren aquesta davant dita teneó<…> aytal fermetat en totes coses deu aver que d’aqui enant per nula moxonia no pusca eser trastomat ni per algun enguiny no pusca eser camiat». P.~152--153.}. В данном случае инфинитив был передан через конструкцию со значением долженствования, также встречается перевод и других латинских инфинитивов при помощи аналитических конструкций с иными оттенками значения, нежели оригинал.

Для того чтобы проиллюстрировать это, обратимся к 54(57) статье «Барселонских обычаев»: \emph{«Феоды, которые держали рыцари, }\emph{если их сеньоры будут отрицать, что передавали их им }(лат.\emph{ }\textbf{non}\textbf{ }\textbf{eos}\textbf{ }\textbf{illis}\textbf{ }\textbf{dedisse}---кат.\textbf{\emph{ }}\textbf{\emph{no}}\textbf{\emph{·}}\textbf{\emph{ls}}\textbf{\emph{ }}\textbf{\emph{los}}\textbf{\emph{ }}\textbf{\emph{agen}}\textbf{\emph{ }}\textbf{\emph{donats}}\emph{) <...>, пусть докажут это <...> и пусть владеют ими. Те же <...>, либо пусть докажут, что приобрели} (лат. \textbf{\emph{adquisisse}}\emph{---}кат.\emph{ }\textbf{\emph{agen}}\textbf{\emph{ }}\textbf{\emph{acaptats}}\emph{) их от своих сеньоров, либо пусть откажутся от них»}\footnote{Ibid. Us. 54 (57): «FEVOS quos tenuerint milites, si seniores eorum negauerint non eos illis dedisse, auerent illos per sacramentum et per bataiam, et habeant illos. Illos autem quos non tenuerint et clamauerint, aut probent per testes uel per scripturas eos a senioribus eorum adquisisse, aut dimittant eos $\rightarrow$ [L]os feus que tenen los cavalers, si los senyors los neguen a éls que no·ls los agen donats, averen-los per sagrament ho per batala, e agen aquels; aquels, però, que no tendran e[·] clamaran, o·ls proven per testimonis ho per escriptures qu·els agen acaptats de lur senyor, ho·ls los jaquesquen». P.~90--91.}. В данном случае мы видим, что инфинитив в рамках придаточного предложения в старокаталанском переводе приобретает личную форму. Причём важно понимать, что подобное происходит не только с глаголами, включёнными в конструкции типа accusativus cum infinitivo или nominativus cum infinitivo.

Таким образом, синтаксическая переработка латинского оригинала в старокаталанском переводе затрагивает несколько взаимосвязанных уровней: способ выражения авторской позиции, средства передачи долженствования и общую организацию юридического высказывания. Колебания в передаче лица глаголов показывают, что переводчик не всегда стремился сохранить связь нормы с голосом её предполагаемого установителя: в одних случаях эта связь удерживается, в других --- уступает место более нейтральному или безличному изложению. Это позволяет предположить, что в старокаталанской версии правовая норма всё чаще воспринималась не как непосредственно произнесённое распоряжение властителя, а как авторитетный письменный текст.

Аналогичным образом передача долженствования показывает, что нормативная направленность латинского оригинала в целом сохраняется, но выражается иными грамматическими средствами: вместо компактных латинских конструкций с герундивом и пассивным описательным спряжением чаще используются аналитические обороты с глаголами «быть должным»\emph{ }(кат. \emph{--- deure}) и «мочь» (кат. ---\emph{ poder})\emph{ }и личными формами глагола. В результате модальность нормы становится более явной и легче расчленяется на значения долга, обязанности, возможности или запрета. Наконец, замена неличных и синтаксически сжатых латинских форм более развёрнутыми конструкциями с личным глаголом свидетельствует о более общей тенденции: латинский текст в большей степени опирается на компактные, абстрактные и номинативные структуры, тогда как старокаталанская версия чаще разворачивает норму аналитически --- через действие, его носителя и его практический результат. В более широком плане это позволяет видеть в переводе не только языковую адаптацию, но и изменение самой модели правовой коммуникации: текст становится в большей степени ориентированным на применение нормы в конкретной социальной среде, где важны не только её формальная авторитетность, но и ясность, расчленённость и практическая считываемость.

\subsection{Ошибки и смещение смыслов при переводе с латинского на старокаталанский}

В старокаталанском переводе «Барселонских обычаев» одним из наиболее заметных проявлений переводческой переработки текста являются опущения и, значительно реже, добавления. Особенно часто сокращению подвергаются длинные перечисления, пояснительные формулы и отдельные части статей, которые, по-видимому, не воспринимались переводчиком как смысловое ядро нормы. Поэтому анализ таких случаев важен не только для характеристики техники перевода, но и для понимания того, каким образом латинский текст переосмыслялся в процессе его адаптации к новой языковой и правовой среде.

Одним из наиболее крупных пропусков является сокращение обширного перечисления в статье 60(64)\footnote{Usatges de Barcelona. Us. 60 (64): «... habeant omni tempore sinceram et perfectam fidem et ueram locutionem; ita ut omnes homines nobiles et ignobiles,\textbf{ reges et principes, magnates et milites, rustici et pagenses, mercerii et negociatores, peregrini et camina tenentes, amici et inimici, christiani et sarraceni, iudei et heretici, possint se in illis fidare et credere non solum illorum personas set eciam ciuitates et castella, honorem et auere, uxores et filios et cuncta que habuerint, sine timore et absque ulla mala suspicione}; et omnes homines nobiles et ignobiles,magnates, milites et pedites, marinarii et cur- sarii et monetarii, in illorum terra stantes uel aliunde aduenientes, adiuuent predictos principes eorum fidem et locutionem tenere …$\rightarrow$per tots temps ayen cincera e acabada fe e vera paraula, en axí que tots los hòmens, nobles e ignobles (...) tots los hòmens, nobles e ignobles), grans hòmens, cavalers e peons, mariners e cosaris e moneders, en aquela terra estans ho alondre vivens ajuden ais damunt dits a lur fe e aquels tenir...». P.~96--97.}. В старокаталанской версии выпадает значительный фрагмент латинского текста, построенный как ряд однородных членов. Дж.~Бастардес-и-Парера предполагал, что такой пропуск мог быть связан с гомеотелевтом\footnote{\textbf{Гомеотелевт }(греч. (греч. {\unicodefallback ὁ}μοιοτέλευτος: {\unicodefallback ὅ}μοιος --- «подобный» --- и τελευτή --- «окончание») --- вид морфемного повтора, при котором на относительно небольшом отрезке текста встречается значительное количество слов с одинаковой финальной частью. См. \emph{Аверинцев С.С.}~Поэтика ранневизантийской литературы. Спб.: Азбука-Классика, 2004. С.~231. Об этом конкретном случае: Usatges de Barcelona. P.~97.}.

Иной характер имеет пропуск в 11(13) статье «Барселонских обычаев», где в старокаталанской версии отсутствует не часть перечисления, а целое предложение, касающееся менее тяжких правонарушений\emph{: }\emph{«}\emph{Увечье и избиение должны возмещаться по закону соответствующим количеством солидов денариев}\emph{»}\footnote{Ibid. Us. 11 (13): «RVSTICVS INTERFECTVS seu alius homo qui nullam habet dignitatem, preter quod christianus est, emendetur per .vi. uncias; plaga quoque uncias duas. \textbf{Debilitacio}\textbf{ }\textbf{et}\textbf{ }\textbf{cedis}\textbf{ }\textbf{emendetur}\textbf{ }\textbf{per}\textbf{ }\textbf{legem}\textbf{ }\textbf{secundum}\textbf{ }\textbf{solidos}\textbf{ }\textbf{denariorum}$\rightarrow$[P]ages mort ho altre hom que nula dignitat no ha, esters que és crestis, sie emenat per .vi. onges; plaga .ii. onges ()». P.~60--61.}. В отличие от предыдущего случая, здесь труднее объяснить сокращение исключительно ошибкой при копировании рукописи. Не исключено, что данный фрагмент был воспринят как второстепенный по отношению к основной норме об убийстве крестьянина, однако надёжно установить мотив такого опущения по одному примеру невозможно.\emph{ }

Сходный случай наблюдается в статье 9(11), где в старокаталанском тексте отсутствует упоминание об убийстве иудея\footnote{Ibid. Us. 9 (11): «IVDEI CESI uel uulnerati, capti aut debilitati siue eciam interfecti, ad uoluntatem potestatis sint emendati$\rightarrow$[J]ueus t(r)enquats ho nafrats, preses e debilitats ( ), a la volentat de la postat sien emenats». P.~58--59.}. Такое опущение может быть значимым, однако делать на его основании вывод о прямом изменении нормы было бы преждевременно. Скорее можно говорить о том, что отдельные элементы санкционной формулы в переводе сокращались, и вопрос о том, было ли это связано с правовой актуализацией текста, требует дополнительной проверки по рукописному материалу и более широкому контексту каталонского законодательства XIII~в.

Наряду с крупными пропусками перевод демонстрирует и более мелкие сокращения, не меняющие основного содержания статьи, но устраняющие то, что переводчик, по-видимому, считал избыточным. Так, в 71(91)статье  развёрнутая латинская формула \emph{«трудящихся на их службе»} (лат. --- \emph{in eorum servicio laborantibus}) передана как \emph{«находящихся у них на службе»} (кат. --- \emph{qui estien en lur servii}). Глагольно-именная конструкция здесь заменена более компактным обозначением службы, сохраняющим основное значение исходной формулы.\footnote{Ibid. Us. 71(91): «… cum omnibus hominibus eorum honorem tenentibus uel in illorum honores  permanentibus siue in eorum seruicio laborantibus…$\rightarrow$…ab tots lurs hómens qui lur honor tinguen ho qui estien en lur honor ho qui estien en lur servii...». P.~108--109.}.

В статье 70 (75) по рукописи K переводчик сократил пространственную формулу \emph{«в пределах наших стен или в бургах»} (лат. --- \emph{infra menia nostra uel burgos}) до слов \emph{«в пределах наших стен» }(старокат. --- \emph{dins los murs nostres})\textsuperscript{ }\footnote{Ibid. Us. 70 (75): «…uel si quis\textbf{ }\textbf{infra}\textbf{ }\textbf{menia}\textbf{ }\textbf{nostra}\textbf{ }\textbf{uel}\textbf{ }\textbf{burgo}s primus eduxerit gladium contra alium uel appellauerit aliquem cuguz …$\rightarrow$… ho si negú \textbf{dins}\textbf{ }\textbf{los}\textbf{ }\textbf{murs}\textbf{ }\textbf{nostres} primerament traurà coltel contra altre o l’apelarà cuguç …». P.~108--109.}. При этом переводчик текста, представленного в рукописи K, хорошо знал обозначение замка: в других статьях он последовательно употреблял слово «замок» (старокат. --- \emph{castell}, в рукописной графике --- \emph{casteyl}), в частности в сочетании «штурм замка» (старокат. --- \emph{assalt de casteyl}).

Следовательно, опущение слова «бурги» (лат. --- \emph{burgos}) нельзя объяснить отсутствием в его словаре наименований укреплённых поселений или смешением бурга с замком. Оно относится к конкретной передаче пространственной формулы в рукописи K. В то же время перевод, опубликованный Ж.~Гудиолем-и-Кунилем, сохранял оба элемента латинского оригинала:\emph{ «в пределах наших стен или в бургах»} (кат. --- \emph{enfre els murs nos}\emph{)}\footnote{Gudiol i Cunill J. Traducció dels Usatges...: «…o si algun enfre els murs nostres o els burcs primerament traura coltel a altre o apelara altre caguz…». P.~295.}. Поскольку здесь речь идёт в целом об укреплённых поселениях, к числу которых относятся и замки, и города, термин \emph{“burgos”}\emph{ }в данном случае кажется ему избыточным исходя из контекста. с При этом, важно понимать, когда речь идёт о социальном статусе того или иного человека, дериваты слова \emph{“burgos”}---например, \emph{“burgeses” }переводчик сохраняет. В таких случаях опущение выступает не как утрата нормы, а как её стилистическое и смысловое уплотнение. Таким образом, локальные сокращения выполняют разные функции: замена сочетания «трудящиеся на службе» более компактной формулой не затрагивает содержания нормы, тогда как опущение слова «бурги» (лат. --- \emph{burgos}) сокращает содержащийся в ней перечень пространств.

Сходную ситуацию можно наблюдать и в 116(137) статье «Барселонских обычаев»: \emph{«…да принесут старшие [по положению.---А.К.] клятву младшим лично, если младшие }\textbf{\emph{смогут }}\textbf{\emph{располагать/дадут}}\emph{ }(лат. \emph{habere potuerint}---кат. \emph{donen) лицами, равными им по статусу, которые принесут клятву за них…»}\footnote{Usatges de Barcelona / ed. J. Bastardas i Parera. Barcelona: Fundació Noguera, 1991. Us. 116 (137): «…iurent maiores minoribus per semedipsos, si minores habere potuerint illorum coequales qui pro eis iurent… ---…jurats los mayors als menors per si ilex e els mayors donen a és [co]nsembles que juren a és…». P.~148--149.}. Сопоставление латинской конструкции «смогут иметь» (лат. --- \emph{habere potuerint}) и старокаталанского глагола «дают» (старокат. --- \emph{donen}) показывает, что в каталанском переводе неточно передан сам смысл нормы --- речь идёт о поручительстве в рамках конкретного процесса. Помимо искажения смысла всей фразы, и, более широко, всей статьи «Барселонских обычаев» переводчик нарушил и грамматическое строение фразы, сделав её содержание неясным и, более того, весьма странным. Это отмечал и издатель источника Ж.~Бастардес-и-Парера\footnote{См. комментарии к тексту «Барселонских обычаев»: Usatges de Barcelona / ed. por Bastardas i Parera J. Barcelona, 1991. P.~149. N. 99.}.

Однако перевод как интерпретация проявляется не только в опущениях. В ряде случаев смысл нормы меняется и без прямого сокращения текста --- за счёт выбора эквивалента или синтаксической перестройки, смещающей юридический акцент. Именно такие случаи особенно важны, поскольку они показывают, что старокаталанская версия не только редуцирует латинский оригинал, но местами и перераспределяет его смысл.

На этом фоне особенно показательно, что лексические добавления в старокаталанском переводе сравнительно редки. Это позволяет предположить, что переводчик, как правило, не стремился расширять текст «Барселонских обычаев» за счёт собственных вставок, а предпочитал работать через сокращение, перестройку и объяснительную передачу уже имеющегося материала. Тем самым пределы его интервенции оказывались асимметричными: опущение использовалось значительно чаще, чем дополнение.

Таким образом, опущения и добавления в старокаталанском переводе «Барселонских обычаев» нельзя рассматривать только как технические особенности текста. В одних случаях они, по-видимому, связаны с механической передачей рукописного материала, в других --- с устранением синтаксической и лексической избыточности, а иногда и вовсе --- с более глубокой переработкой нормы в процессе её адаптации. Поэтому перевод выступает здесь, скорее, как форма сокращения и частичного переосмысления латинского оригинала. В этом смысле старокаталанская версия фиксирует не только движение от одного языка к другому, но и изменение самой практики чтения и употребления правового текста.

\textbf{Выводы к §2.4}\textbf{.}

Часто за изменением грамматических категорий в переводе «Барселонских обычаев» стояли глубокие изменения в правовой культуре носителей данной традиции. Грамматические изменения являются хорошей «оптикой» для анализа переосмысления правового наследия более ранних эпох в XIII~в.

Больше всего изменений мы видим в глагольных формах и выражениях действия. Эти изменения показывают конкретизацию мышления, постепенный отход от латинских юридических абстрактных понятий и возвращение к самым базовым глаголам, их более простым формам и т.~д. Мы вновь возвращаемся к тому обстоятельству, что для переводчика «Барселонских обычаев» результат того или иного действия был важнее его характеристик, типологии, места среди других поступков и правонарушений. Эти нюансы он уже практически не отличает.

Именно этот процесс привёл к отдельным проявлениям усиления значения долженствования, которое стало выраженно эксплицитно. В отдельных статьях долженствование, выраженное через конъюнктив в латинской версии «Барселонских обычаев», передаётся глаголом «быть должным» (кат. --- \emph{deure}).

Однако вместе с упрощением формы приходит и необходимость объяснять те или иные понятия и термины. Этим можно объяснить активное использование объяснительного перевода, который помогал дать аудитории основную информацию, а дополнительную опустить или вынести в придаточное предложение. Экспликация применяется в основном к сложным латинским (в~т.ч. и взятым из народной латыни) терминам, которые, по мнению переводчика, могли оказаться непонятными для предполагаемых читателей «Барселонских обычаев».

Таким образом, изменения в представлениях о праве и о значении в нём обычаев, в том, как их понимали и осмысляли составители и переводчики, ярко проявляются в грамматическом строе этого законодательного кодекса и могут быть выявлены только путём сравнения двух версий «Барселонских обычаев» на разных языках.

\section*{Выводы к Главе 2}
\addcontentsline{toc}{section}{Выводы к Главе 2}

Проведённый анализ показывает, какую роль языковые формы играли в правовой культуре Каталонии XII--XIII~вв. Латинская письменная традиция, старокаталанское словоупотребление и устные речевые практики образовывали общую коммуникативную среду, внутри которой происходили запись обычая, составление документов, перевод правовых текстов и регулирование судебной процедуры. Выбор языковой формы определял способы выражения нормы, роли участников и порядок совершения юридически значимых действий.

Сравнение каталонских сборников обычного права показывает разнообразие форм письменной фиксации устного слова. В «Барселонских обычаях» словесное действие обычно передавалось кратким указанием на обвинение, клятву, заявление или распоряжение. «Обычаи Льейды» и «Обычаи Орты» содержали лаконичные формулы возбуждения дела, «Обычаи Миравета» определяли сведения, которые следовало назвать при предъявлении требования, а «Кутюмы Перпиньяна» закрепляли значение публичного документа для словесного оформления иска. В королевской привилегии «Знатные мужи Барселоны признали…» (Recognoverunt proceres, 1284) устное волеизъявление включалось в порядок свидетельского и нотариального удостоверения.

Наиболее развёрнутая система речевых формул представлена в «Обычаях Тортосы». Вопросы, ответы, заявления и распоряжения распределялись между судьями, сторонами и должностными лицами, организуя ход судебного разбирательства. Прямая речь определяла содержание вызова в суд, форму обвинения, порядок распоряжения имуществом, назначения опекуна и наследника. Такие реплики представляли собой нормативно обработанные модели коммуникации, закреплявшие юридически достаточный состав высказывания и его правовые последствия.

Актовый материал XII~в. показывает, как словесное действие оформлялось документально. Формулы окончательного отказа от притязаний, признания права другой стороны и уступки имущества стали устойчивыми элементами формуляра. В отдельных грамотах фиксировались момент произнесения слов, способность лица говорить, а также присутствие тех, кто видел и слышал совершённое действие. Указания на сохранную память, здравый рассудок и речь обозначали способность человека осознанно выразить свою волю. Документ закреплял содержание волеизъявления и обеспечивал возможность его последующего предъявления в суде.

Глаголы речи выполняли в правовых текстах структурную функцию. Формы со значениями «говорить», «отвечать», «спрашивать», «жаловаться», «называть» и «обвинять» обозначали тип совершаемого действия и распределяли процессуальные роли. Их употребление обеспечивало определённость обвинения, требования, ответа или распоряжения и включало высказывание в установленный порядок правового действия.

Анализ иноязычных элементов выявил внутреннюю многослойность латинского языка каталонского права. В «Барселонских обычаях» и актовом материале местные слова обозначали военные действия, технические приспособления, формы землепользования, хозяйственные объекты, имущественные институты и словесные оскорбления. Одни термины сопровождались указаниями на народное или крестьянское употребление, другие непосредственно включались в латинский текст как устойчивые элементы правового словаря. Такие способы фиксации отражали представления составителей о социальных и функциональных регистрах речи.

Особое значение имели слова, посредством которых совершалось правонарушение. Обозначения вероотступника, ренегата и рогоносца сохранялись в формах, близких к старокаталанскому употреблению, поскольку правовая квалификация зависела от произнесения узнаваемого оскорбительного наименования. Сопоставление латинского текста 70 статьи «Барселонских обычаев» с переводом рукописи K, опубликованным Ж.~Бастардесом-и-Парерой, и текстом рукописи, изданным Ж.~Гудиолем-и-Кунилем, показывает устойчивость передачи самого оскорбительного обозначения при вариативности других элементов нормы.

Локальная терминология приданого, выморочного имущества, ирригационных сооружений, угодий и военных действий закрепляла понятия, сформировавшиеся в каталонской правовой и хозяйственной практике. Арабские, окситанские и романские по происхождению слова обеспечивали точное обозначение объектов и отношений, входивших в сферу правового регулирования. Их включение в латинский текст придавало местным категориям письменную и терминологическую устойчивость.

В каталаноязычных памятниках продолжалось освоение латинского понятийного аппарата. В старокаталанском переводе «Барселонских обычаев» латинизмы употреблялись наряду с местными эквивалентами, а выбор между ними зависел от контекста и характера действия. В «Обычаях Тортосы» латинские элементы использовались в названиях рубрик, формулах подписей, дефинициях, нормативных положениях и цитатах из авторитетных текстов. Старокаталанские пояснения раскрывали их содержание применительно к местной практике. Редакционная переработка памятника сопровождалась отбором латинских фрагментов в соответствии с их нормативной и практической функцией.

Сопоставление латинского текста «Барселонских обычаев» с двумя старокаталанскими версиями выявило устойчивость основных переводческих решений и вариативность рукописной передачи. Тексты, опубликованные Ж.~Бастардесом-и-Парерой и Ж.~Гудиолем-и-Кунилем, содержат общие терминологические и синтаксические соответствия, а также разночтения в выборе слов, составе формул и объёме передаваемого материала. Перевод затрагивал лексику, грамматику и синтаксис нормы: изменялись способы выражения долженствования и возможности, грамматическое лицо, структура перечислений и смысловые акценты. Старокаталанские версии сохраняли правовой авторитет памятника и формировали собственные способы выражения нормы.

В целом языковая многослойность составляла устойчивую характеристику правовой культуры Каталонии XII--XIII~вв. Латынь обеспечивала связь с книжной, документальной и учёной традицией, старокаталанский язык служил средством перевода, толкования и формулирования местных норм, а речевые формулы задавали словесную форму процессуальных и имущественных действий.
